\section{Phương trình, bất phương trình mũ và lôgarit}

\begin{dang}{Phương trình mũ và phương trình logarit cơ bản}
	\begin{enumerate}
		\item Phương trình mũ cơ bản có dạng $a^x=b\ \left(a>0,a\neq 1\right)$.
		\begin{itemize}
			\item Nếu $b\le0$ thì phương trình vô nghiệm.
			\item Nếu $b>0$ thì phương trình có nghiệm duy nhất $x=\log_a b$. 
		\end{itemize}
		\item Phương trình logarit cơ bản có dạng $\log_a x=b\ \left(a>0,a\neq 1\right)$.\\ 
		Phương trình đó có một nghiệm là $x=a^b$
	\end{enumerate}
	
	
\end{dang}
\subsubsection{Các ví dụ}
\Opensolutionfile{ans}[ans/ans-1K6-21-Dang3]
\begin{vd}%[1K6KK-2] 
	Giải các phương trình sau:
	\begin{enumEX}{2}
		\item $3^{x-3}=5$;
		\item $10^{2x+1}=2023$;
		\item $5^{2x}=\dfrac{1}{5^{4-x}}$;
		\item $10^{x+1}-4\cdot10^x=12$.
	\end{enumEX}
\loigiai{
	\begin{enumerate}
		\item $3^{x-3}=5 \Leftrightarrow x-3=\log_3 5 \Leftrightarrow x=\log_3 5+3$.\\
		Vậy phương trình có nghiệm là $x=\log_3 5+3$. 
		\item $10^{2x+1}=2023 \Leftrightarrow 2x+1=\log 2023 \Leftrightarrow x=\dfrac{1}{2}\left(\log 2023-1\right)$.\\
		Vậy phương trình có nghiệm là $x=\dfrac{1}{2}\left(\log 2023-1\right)$. 
		\item $5^{2x}=\dfrac{1}{5^{4-x}} \Leftrightarrow 5^{2x}=5^{x-4} \Leftrightarrow 2x=x-4 \Leftrightarrow x=-4$.\\
		Vậy phương trình có nghiệm là $x=-4$.
		\item {\allowdisplaybreaks
			\begin{eqnarray*}
			10^{x+1}-4\cdot10^x=12 &\Leftrightarrow& 10\cdot 10^x-4\cdot 10^x=12\\ &\Leftrightarrow& 6\cdot 10^x=12\\ &\Leftrightarrow& 10^x=2\\ &\Leftrightarrow& x=\log 2.	
		\end{eqnarray*}}
	Vậy phương trình có nghiệm là $x=\log 2$.
	\end{enumerate}
}
\end{vd}
\begin{vd}%[1K6KK-2] 
	Giải các phương trình sau:
	\begin{enumEX}{2}
	\item $\log_3 x=5$;
	\item $3\log_2 {(2x-4)}+1=10$;
	\item $\log_8(3x-6)=\log_8(2x-2)$;
	\item $\log_4 (x+2)+\log_4 (x-1)=1$.
	\end{enumEX}
\loigiai{
	\begin{enumerate}
	\item Điều kiện: $x>0$.\\
	Khi đó $\log_3 x=5 \Leftrightarrow x=3^5$ (thỏa).\\
	Vậy phương trình có nghiệm là $x=3^5$.
	\item Điều kiện: $2x-4>0 \Leftrightarrow x>2$.
	Khi đó $$3\log_2 {(2x-4)}+1=10 \Leftrightarrow \log_2 (2x-4)=3 \Leftrightarrow 2x-4=2^3 \Leftrightarrow x=6 \text{ (thỏa).}$$
	Vậy phương trình có nghiệm là $x=6$.
	\item Điều kiện: $\heva{&3x-6>0\\&2x-2>0} \Leftrightarrow x>2$.\\
	Khi đó $$\log_8(3x-6)=\log_8(2x-2) \Leftrightarrow 3x-6=2x-2 \Leftrightarrow x=4 \text{ (thỏa)}.$$
	Vậy phương trình có nghiệm là $x=4$.
	\item Điều kiện: $\heva{&x+2>0\\&x-1>0} \Leftrightarrow x>1$.\\
	Khi đó {\allowdisplaybreaks
		\begin{eqnarray*}
		&&\log_4 (x+2)+\log_4 (x-1)=1\\ &\Leftrightarrow& \log_4 \left[(x+2)(x-1)\right]=1 \\
		&\Leftrightarrow& (x+2)(x-1)=4\\
		&\Leftrightarrow& x^2+x-6=0 \ \Leftrightarrow  \hoac{&x=2 \text{ (thỏa)}\\&x=-3 \text{ (loại).}}	
	\end{eqnarray*}}
	Vậy phương trình có nghiệm là $x=2$.
	\end{enumerate}
}
\end{vd}
\subsubsection{Bài tập vận dụng}
\begin{bt}%[1K6KK-2] 
	Giải các phương trình sau:
	\begin{enumEX}{2}
		\item $3^{5x-2}=81$;
		\item $\mathrm{e}^{2x-1}=5$;
		\item $5^{2x+3}=125^{x-2}$;
		\item $100^{2x^2-3}=0{,}1^{2x^2-18}$;
		\item $3^{3x}=\dfrac{1}{3^{6-2x}}$;
		\item $5^x=3^{2x-1}$.
	\end{enumEX}
\loigiai{
	\begin{enumerate}
		\item $3^{5x-2}=81 \Leftrightarrow 3^{5x-2}=3^4 \Leftrightarrow 5x-2=4 \Leftrightarrow x=\dfrac{6}{5}$.\\
		Vậy phương trình có nghiệm là $x=\dfrac{6}{5}$.
		\item $\mathrm{e}^{2x-1}=5 \Leftrightarrow 2x-1=\ln 5 \Leftrightarrow x=\dfrac{\ln 5+1}{2}$.\\
		Vậy phương trình có nghiệm là $x=\dfrac{\ln 5+1}{2}$.
		\item {\allowdisplaybreaks
			\begin{eqnarray*}
			5^{2x+3}=125^{x-2} &\Leftrightarrow& 5^{2x+3}=5^{3(x-2)}\\
			&\Leftrightarrow& 2x+3=3x-6\\
			&\Leftrightarrow&  x=9.
		\end{eqnarray*}}
		Vậy phương trình có nghiệm là $x=9$.
		\item {\allowdisplaybreaks
			\begin{eqnarray*}
			100^{2x^2-3}=0{,}1^{2x^2-18} &\Leftrightarrow& 10^{2\left(2x^2-3\right)}=10^{-\left(2x^2-18\right)}\\ 
			&\Leftrightarrow & 4x^2-6=-2x^2+18 \\
			&\Leftrightarrow& x=\pm 2.	
		\end{eqnarray*}}
		Vậy phương trình có nghiệm là $x=\pm 2$.
		\item $3^{3x}=\dfrac{1}{3^{6-2x}} \Leftrightarrow 3^{3x}=3^{2x-6} \Leftrightarrow 3x=2x-6 \Leftrightarrow x=-6$.\\
		Vậy phương trình có nghiệm là $x=-6$.
		\item Lấy logarit cơ số 3 hai vế ta được: {\allowdisplaybreaks
			\begin{eqnarray*}
			x\cdot \log_3 5=2x-1 &\Leftrightarrow& \left(2-\log_3 5\right)x=1\\ &\Leftrightarrow& x=\dfrac{1}{2-\log_3 5}.	
		\end{eqnarray*}}
		Vậy phương trình có nghiệm là $x=x=\dfrac{1}{2-\log_3 5}$.
	\end{enumerate}
}
\end{bt}
\begin{bt}%[1K6KK-2] 
	Giải các phương trình sau: 
	\begin{enumEX}{2}
		\item $\log(3x+10)=2$;
		\item $\log_{\tfrac{1}{3}}\left(2x-3\right)=-3$;
		\item $\log_5{\left(3x-5\right)}=\log_5{\left(2x+1\right)}$;
		\item $2\log_4 x+\log_2\left(x-3\right)=2$;
		\item $\ln x+\ln \left(x-1\right)=\ln 4x$;
		\item $\log_5 \left(x^2-3x+2\right)=-2\log_{\tfrac{1}{5}}\left(x-1\right)$.	
	\end{enumEX}
\loigiai{
	\begin{enumerate}
		\item Điều kiện: $x>-\dfrac{10}{3}$.\\
		Khi đó $$\log(3x+10)=2 \Leftrightarrow 3x+10=10^2 \Leftrightarrow x=30 \text{ (thỏa).}$$
		Vậy phương trình có nghiệm là $x=30$.
		\item Điều kiện: $x>\dfrac{3}{2}$.\\
		Khi đó {\allowdisplaybreaks
			\begin{eqnarray*}
			\log_{\tfrac{1}{3}}\left(2x-3\right)=-3 &\Leftrightarrow& -\log_3 \left(2x-3\right)=-3\\
			&\Leftrightarrow& 2x-3=3^3 \Leftrightarrow x=15 \text{ (thỏa)}.	
		\end{eqnarray*}}
		Vậy phương trình có nghiệm là $x=15$.
		\item Điều kiện: $\heva{&3x-5>0\\&2x+1>0} \Leftrightarrow x>\dfrac{5}{3}$.\\
		Khi đó $$\log_5{\left(3x-5\right)}=\log_5{\left(2x+1\right)} \Leftrightarrow 3x-5=2x+1 \Leftrightarrow x=6 \text{ (thỏa).}$$
		Vậy phương trình có nghiệm là $x=6$.
		\item Điều kiện: $\heva{&x>0\\&x-3>0} \Leftrightarrow x>3$.\\
		Khi đó{\allowdisplaybreaks
			\begin{eqnarray*}
			2\log_4 x+\log_2\left(x-3\right)=2 &\Leftrightarrow& \log_2x+\log_2\left(x-3\right)=2\\
			&\Leftrightarrow& \log_2 \left[x\left(x-3\right)\right]=2\\
			&\Leftrightarrow& x\left(x-3\right)=2^2\\
			&\Leftrightarrow& x^2-3x-4=0 \Leftrightarrow \hoac{&x=-1 \text{ (loại)}\\& x=4 \text{ (thỏa).}}	
		\end{eqnarray*}} 
		Vậy phương trình có nghiệm là $x=4$.
		\item Điều kiện: $\heva{&x>0\\& x-1>0} \Leftrightarrow x>1$.\\
		Khi đó: {\allowdisplaybreaks
			\begin{eqnarray*}
			\ln x+\ln \left(x-1\right)=\ln 4x &\Leftrightarrow& \ln \left[x\left(x-1\right)\right]=\ln 4x\\
			&\Leftrightarrow& x\left(x-1\right)=4x\\ &\Leftrightarrow& x^2-5x=0 \Leftrightarrow \hoac{&x=0 \text{ (loại)}\\& x=5 \text{ (thỏa).}}
		\end{eqnarray*}}
	Vậy phương trình có nghiệm là $x=5$.
		\item Điều kiện: $\heva{&x^2-3x+2>0\\& x-1>0} \Leftrightarrow x>2$.\\
		Khi đó: {\allowdisplaybreaks
			\begin{eqnarray*}
			\log_5 \left(x^2-3x+2\right)=-2\log_{\tfrac{1}{5}}\left(x-1\right) &\Leftrightarrow& \log_5 \left(x^2-3x+2\right)=\log_{5}\left(x-1\right)^2 \\
			&\Leftrightarrow& x^2-3x+2=\left(x-1\right)^2\\ &\Leftrightarrow& x=1 \text{ (loại).}	
		\end{eqnarray*}} 
		Vậy phương trình vô nghiệm.
	\end{enumerate}
}
\end{bt}
\subsubsection{Bài tập trắc nghiệm}
\begin{ex}%[1K6BK-2] 
	Nghiệm của phương trình $3^{2x-1}=27$ là
	\choice
	{$1$}
	{\True$2$}
	{$3$}
	{$4$}
	\loigiai{
	$3^{2x-1}=27 \Leftrightarrow 3^{2x-1}=3^3 \Leftrightarrow 2x-1=3 \Leftrightarrow x=2$.\\
	Vậy phương trình có nghiệm $x=2$.	
	}
\end{ex}
\begin{ex}%[1K6BK-2]
	Nghiệm của phương trình $5^{x-3}=7$ là
	\choice
	{$3$}
	{$5^7-3$}
	{\True $\log_5 7 + 3$}
	{$\log_7 5 - 3$}
	\loigiai{
	$5^{x-3}=7 \Leftrightarrow x-3=\log_5 7 \Leftrightarrow x=\log_5 7+3$.\\
	Vậy phương trình có nghiệm $x=\log_5 7+3$.	
	}
\end{ex}
\begin{ex}%[1K6KK-2]
	Tập nghiệm của phương trình $\mathrm{e}^{x^2-5x+4}-1=0$ là
	\choice
	{$S=\{1\}$}
	{$S=\{1;-4\}$}
	{\True $S=\{1;4\}$}
	{$S=\{-1;5\}$}
	\loigiai{
	$\mathrm{e}^{x^2-5x+4}-1=0 \Leftrightarrow \mathrm{e}^{x^2-5x+4}={e}^0 \Leftrightarrow x^2-5x+4=0 \Leftrightarrow \hoac{&x=1\\& x=4.}$	\\
	Vậy tập nghiệm của phương trình là $S=\{1;4\}$.
	}
\end{ex}
\begin{ex}%[1K6KK-2]
	Phương trình $2^{3x+1}=\dfrac{1}{2^{3-x^2}}$ có hai nghiệm $x_1,x_2$ khi đó giá trị của $P=x_1+x_2$ bằng
	\choice
	{$-3$}
	{$-1$}
	{\True $3$}
	{$4$}
	\loigiai{
	$2^{3x+1}=\dfrac{1}{2^{3-x^2}} \Leftrightarrow 2^{3x+1}=2^{x^2-3} \Leftrightarrow 3x+1=x^2-3 \Leftrightarrow \hoac{&x=-1\\& x=4.}$\\
	Vậy $P=x_1+x_2=-1+4=3$.	
	}
\end{ex}
\begin{ex}%[1K6KK-2]
	Giá trị tổng các nghiệm của phương trình $10000^x=0{,}1^{5-x^2}$ là 
	\choice
	{$0$}
	{$2$}
	{\True $4$}
	{$5$}
	\loigiai{
	$10000^x=0{,}1^{5-x^2} \Leftrightarrow 10^{4x}=10^{-\left(5-x^2\right)} \Leftrightarrow 4x=x^2-5 \Leftrightarrow \hoac{&x=-1\\& x=5.}$\\
	Vậy tổng nghiệm của phương trình bằng $-1+5=4$.	
	}
\end{ex}
\begin{ex}%[1K6BK-2]
	Nghiệm của phương trình $\log_3 (4x-7)=2$ là
	\choice
	{$1$}
	{$2$}
	{$3$}
	{\True $4$}
	\loigiai{
		Điều kiện: $x>\dfrac{7}{4}$.\\
		Khi đó $$\log_3 (4x-7)=2 \Leftrightarrow 4x-7=3^2 \Leftrightarrow x=4 \text{ (thỏa).}$$
		Vậy phương trình có nghiệm $x=4$.	
	}
\end{ex}
\begin{ex}%[1K6BK-2]
	Nghiệm của phương trình $\log_{\tfrac{1}{5}}(4x+5)=-2$ là
	\choice
	{$0$}
	{\True $5$}
	{$25$}
	{$2$}
	\loigiai{
		Điều kiện: $x>-\dfrac{5}{4}$.\\
		Khi đó: $$\log_{\tfrac{1}{5}}(4x+5)=-2 \Leftrightarrow -\log_5 (4x+5)=-2 \Leftrightarrow 4x+5=5^2 \Leftrightarrow x=5 \text{ (thỏa).}$$
		Vậy phương trình có nghiệm $x=5$.	
	}
\end{ex}
\begin{ex}%[1K6KK-2]
	Tập nghiệm của phương trình $\log (12-x)=2\log x$ là
	\choice
	{$S=\{-4;3\}$}
	{\True $S=\{3\}$}
	{$S=\{4\}$}
	{$S=\{3;12\}$}
	\loigiai{
		Điều kiện: $0<x<12$.\\
		Khi đó{\allowdisplaybreaks
			\begin{eqnarray*}
				\log (12-x)=2\log x &\Leftrightarrow& \log (12-x)=\log x^2 \\
				&\Leftrightarrow& x^2+x-12=0 \Leftrightarrow \hoac{&x=-4 \text{ (loại)}\\& x=3 \text{ (thỏa)}.}	
		\end{eqnarray*}} 
		Vậy phương trình có nghiệm $x=3$.	
	}
\end{ex}
\begin{ex}%[1K6KK-2]
	Nghiệm của phương trình $\ln (2x-1)+ \ln(x+3)=\ln 4x$ là 
	\choice
	{$0$}
	{$1$}
	{$2$}
	{\True $3$}
	\loigiai{
		Điều kiện: $x>\dfrac{1}{2}$.\\
		Khi đó {\allowdisplaybreaks
			\begin{eqnarray*}
				\ln (2x-1)+ \ln(x+3)=\ln 4x &\Leftrightarrow& \ln \left[(2x-1)(x+3)\right]=\ln 4x \\
				&\Leftrightarrow& 2x^2+x-3=0 \Leftrightarrow \hoac{&x=1 \text{ (thỏa)}\\& x=-\dfrac{3}{2} \text{ (loại).}}
		\end{eqnarray*}}
		Vậy phương trình có nghiệm $x=3$.		
	}
\end{ex}
\begin{ex}%[1K6KK-2]
	Nghiệm của phương trình $\log_{\tfrac{1}{7}}(x+9)=\log_{\tfrac{1}{7}}(2x-1)$ là
	\choice
	{$1$}
	{$5$}
	{\True $10$}
	{$12$}
	\loigiai{
		Điều kiện: $x>\dfrac{1}{2}$.\\
		Khi đó {\allowdisplaybreaks
			\begin{eqnarray*}
				\log_{\tfrac{1}{7}}(x+9)=\log_{\tfrac{1}{7}}(2x-1)
				&\Leftrightarrow & -\log_7 (x+9)=-\log_7 (2x-1) \\
				& \Leftrightarrow & \log_7 (x+9)=\log_7 (2x-1) \\
				& \Leftrightarrow & 	x+9=2x-1 \Leftrightarrow x=10 \text{ (thỏa).}
		\end{eqnarray*}}
		Vậy phương trình có nghiệm $x=10$.
	}
\end{ex}
\Closesolutionfile{ans}
\begin{indapan}{10}
	{ans/ans-1K6-21-Dang3}
\end{indapan}


%\chapter{Hàm số  lượng giác và phương trình lượng giác}%Hoài Nam


\begin{dang}{Bất phương trình mũ, lôgarit cơ bản}
	\begin{itemize}
		\item Bất phương trình mũ cơ bản có dạng $a^x>b$ (hoặc $a^x\geq b$, $a^x<b$, $a^x\leq b$) với $(a>0,\,a\neq 1)$.
		\item Bất phương trình lôgarit cơ bản có dạng $\log_a x>b$ (hoặc $\log_a x\geq b$, $\log_a x<b$, $\log_a x\leq b$) với $(a>0,\,a\neq 1)$.
	\end{itemize}
\end{dang}
\subsubsection{Ví dụ mẫu}

\begin{vd}%[DCHT Toán 11 - KNTT -Nguyễn Hoài Nam]%[1K6BK-3]
	Giải các bất phương trình sau
	\begin{listEX}[3]
		\item $\left(\dfrac{1}{5}\right)^{x^2-5x+3}>125$.
		\item $3^{x+2}+3^{x-1}\le 28$.
		\item $\left(\dfrac{4}{7}\right)^{2x^2-3x}<\dfrac{7}{4}$.
	\end{listEX}
	\loigiai{
		\begin{enumerate}[a)]
			\item $\left(\dfrac{1}{5}\right)^{x^2-5x+3}>125  \Leftrightarrow x^2-5x+3<-3 \Leftrightarrow x^2-5x+6 <0 \Leftrightarrow 2<x<3$.\\
			Vậy tập nghiệm của bất phương trình là $S=(2;3)$.
			\item $3^{x+2}+3^{x-1}\le 28 \Leftrightarrow 3^x\left(3^2+3^{-1}\right)\le 28 \Leftrightarrow 3^x \le 3 \Leftrightarrow x\le 1$.\\
			Vậy tập nghiệm của bất phương trình là $S=(-\infty;1]$.
			\item $\left(\dfrac{4}{7}\right)^{2x^2-3x}<\dfrac{7}{4} \Leftrightarrow 2x^2-3x > \log_{\frac{4}{7}} \frac{7}{4} \Leftrightarrow 2x^2-3x>-1 \Leftrightarrow 2x^2-3x+1>0 \Leftrightarrow \hoac{&x<\frac{1}{2}\\&x>1.}$\\
			Vậy tập nghiệm của bất phương trình là $S=\left(-\infty;\dfrac{1}{2}\right)\cup (1;\infty)$.
		\end{enumerate}
	}
\end{vd}

\begin{vd}%[DCHT Toán 11 - KNTT -Nguyễn Hoài Nam]%[1K6BK-3]
	Giải bất phương trình $\log_{\sqrt{3}} (x^2-3x+11)\le4$.
	\loigiai{
		Ta có 
		\begin{eqnarray*}
			\log_{\sqrt{3}} (x^2-3x+11)\le4 \Leftrightarrow 0<x^2-3x+11\le 9 &\Leftrightarrow& \heva{&x^2-3x+11>0 \text{ (đúng với mọi $ x \in \mathbb{R}$)}\\&x^2-3x+11\le9}\\
			&\Leftrightarrow& x^2-3x+2\le0 \\
			&\Leftrightarrow& 1 \le x \le 2.
		\end{eqnarray*}
		Vậy tập nghiệm của bất phương trình là $S=\left[1;2\right]$.
	}
\end{vd}

\begin{vd}%[DCHT Toán 11 - KNTT -Nguyễn Hoài Nam]%[1K6BK-3]
	Giải bất phương trình $ \log_{\frac{1}{15}}( x-2)+\log_{\frac{1}{15}}( 10-x) \ge -1 $.
	\loigiai{
		Điều kiện $ 2<x<10 $.\\
		Ta có
		\begin{eqnarray*}
			\log_{\frac{1}{15}}( x-2)+\log_{\frac{1}{15}}( 10-x) \ge -1& \Leftrightarrow& \log_{\frac{1}{15}}\left[( x-2)( 10-x)\right]\ge \log_{\frac{1}{15}}15\\
			&\Leftrightarrow &( x-2)( 10-x)\le 15\\ 
			&\Leftrightarrow & x^2-12x+35\le 0 \\
			&\Leftrightarrow &x\le 5 \text{ hoặc }x\ge 7.
		\end{eqnarray*}
		Kết hợp với điều kiện ta được tập nghiệm của bất phương trình là $ S=\left(2;5\right]\cup \left[7;10\right)$.
	}	
\end{vd}
\begin{vd}%[DCHT Toán 11 - KNTT -Nguyễn Hoài Nam]%[1K6BK-3]
	Giải bất phương trình
	$\left(\sqrt{10}+3\right)^{\frac{x-3}{x-1}} < \left(\sqrt{10}-3\right)^{\frac{x+1}{x+3}}$.
	\loigiai{
		Điều kiện $\heva{&x\neq -3\\&x\neq 1.}$ \\
		Khi đó
		\begin{eqnarray*}
			\left(\sqrt{10}+3\right)^{\frac{x-3}{x-1}} < \left(\sqrt{10}-3\right)^{\frac{x+1}{x+3}} &\Leftrightarrow&
			\left(\sqrt{10}+3\right)^{\frac{x-3}{x-1}} < \left(\sqrt{10}+3\right)^{-\left(\frac{x+1}{x+3}\right)}\\
			&		\Leftrightarrow & \dfrac{x-3}{x-1} <-\dfrac{x+1}{x+3}\\
			&\Leftrightarrow &\dfrac{x^2-5}{(x-1)(x+3)}<0\\
			&\Leftrightarrow &-3 < x < -\sqrt{5} \text{ hoặc } 1 <x< \sqrt{5}.
		\end{eqnarray*}
		Kết hợp với điều kiện ta được tập nghiệm của bất phương trình là
		$S= (-3; -\sqrt{5}) \cup (1; \sqrt{5})$.
	}
\end{vd}
\subsubsection{Bài tập rèn luyện}
\centerline{\fcolorbox{red}{yellow!50}{\bf {BÀI TẬP TỰ LUẬN }}}
\begin{bt}%[1K6BK-3]
	Giải bất phương trình $\log_{\tfrac{1}{2}}(2x)>-2$.
	\loigiai{
		Vì $0<\dfrac{1}{2}<1$ nên $\log_{\tfrac{1}{2}}(2x)>-2\Leftrightarrow 0<2x<4\Leftrightarrow 0<x<2$.
	}	
\end{bt}
\begin{bt}%[DCHT Toán 11 - KNTT -Nguyễn Hoài Nam]%[1K6BK-3]
	Giải bất phương trình $\left(\dfrac{\mathrm{e}}{\pi}\right)^x>1$.
	\loigiai{
		Vì $0<\dfrac{\mathrm{e}}{\pi}<1$ nên $\left(\dfrac{\mathrm{e}}{\pi}\right)^x>1\Leftrightarrow x<0$.\\
		Vậy tập nghiệm của bất phương trình là $S=\left(-\infty;0\right)$.
	}
\end{bt}

\begin{bt}%[DCHT Toán 11 - KNTT -Nguyễn Hoài Nam]%[1K6BK-3]
	Giải bất phương trình $\log_{\frac{1}{2}}(x+1)<\log_{\frac{1}{2}}(2x-1)$.
	\loigiai{
		Điều kiện xác định $\heva{& x+1>0 \\ & 2x-1>0}\Leftrightarrow x>\frac{1}{2}$.	Ta có $$\log_{\frac{1}{2}}(x+1)<\log_{\frac{1}{2}}(2x-1) \Leftrightarrow x+1>2x-1 \Leftrightarrow x<2.$$
		Kết hợp với điều kiện xác định, ta có tập nghiệm của bất phương trình là $S=\left(\dfrac{1}{2};2\right)$.
	}
\end{bt}
\begin{bt}%[DCHT Toán 11 - KNTT -Nguyễn Hoài Nam]%[1K6BK-5]
	Giải bất phương trình $\left(\dfrac{2}{3}\right)^{4x}\le \left(\dfrac{2}{3}\right)^{x-2}$.
	\loigiai{
		Ta có
		$$\left(\dfrac{3}{2}\right)^{-4x}\le \left(\dfrac{3}{2}\right)^{2-x}
		\Leftrightarrow -4x\le 2-x
		\Leftrightarrow x\ge -\dfrac{2}{3}.$$	
	}
\end{bt}

\centerline{\fcolorbox{red}{yellow!50}{\bf {CÂU HỎI TRẮC NGHIỆM (Tầm 10 - 20 câu theo theo tỉ lệ 4:3:2:1)}}}
\Opensolutionfile{ans}[ans/ans-1K6-21-Dang4]
\begin{ex}%[DCHT Toán 11 - KNTT -Nguyễn Hoài Nam]%[1K6BK-3] 
	Tập nghiệm của bất phương trình $3^{x-1}>5$ là
	\choice
	{\True $\left( \log_{3}15; + \infty\right) $}
	{$\left( - \infty; \log_{3}15 \right) $}
	{$\left( \log_3{5}; + \infty\right) $}
	{$\left( \log_5{3}; + \infty\right) $}
	\loigiai
	{
		Ta có
		\[3^{x-1}>5 \Leftrightarrow x - 1 > \log_{3}5  \Leftrightarrow x >  1 +\log_{3}5 \Leftrightarrow x > \log_{3} 15.\]
		Vậy tập nghiệm của bất phương trình là $S =\left( \log_{3}15; + \infty\right)$.
	}
\end{ex}
\begin{ex}%[DCHT Toán 11 - KNTT -Nguyễn Hoài Nam]%[1K6BK-3] 
	Số nghiệm nguyên của bất phương trình $\log_{6}(2x+4)<2$ là
	\choice
	{$19$}
	{\True $17$}
	{$18$}
	{$20$}
	\loigiai{
		Điều kiện: $2x+4>0\Leftrightarrow x>-2$.\\
		Với điều kiện trên bất phương trình đã cho tương đương
		$$2x+4<6^2\Leftrightarrow x<16.$$
		So sánh điều kiện, tập nghiệm của phương trình là $S=(-2;16)$.\\
		Số nghiệm nguyên của phương trình là $17$.
	}
\end{ex}
\begin{ex}%[DCHT Toán 11 - KNTT -Nguyễn Hoài Nam]%[1K6BK-3] 
	Tập nghiệm $S$ của bất phương trình $\log_2(x+2)\le 0$ là
	\choice
	{$S=[-1;+\infty)$}
	{$S=(-2;+\infty)$}
	{\True $S=(-2;-1]$}
	{$S=(0;-1]$}
	\loigiai
	{Điều kiện: $x>-2$.
		{\allowdisplaybreaks
			\begin{eqnarray*}
				\log_2(x+2)\le 0\Leftrightarrow x+2\le 1
				\Leftrightarrow
				x\le -1.
			\end{eqnarray*}
		}	
		Kết hợp với điều kiện, tập nghiệm của bất phương trình là $S=(-2;-1]$.
	}
\end{ex}


\begin{ex}%[DCHT Toán 11 - KNTT -Nguyễn Hoài Nam]%[1K6BK-3] 
	Tập nghiệm của bất phương trình là $5^x>\dfrac{1}{25}$ là
	\choice
	{$(-1;+\infty)$}
	{\True $(-2;+\infty)$}
	{$(5;+\infty)$}
	{$(2;+\infty)$}
	\loigiai{
		Ta có 
		$$5^x>\dfrac{1}{25}\Leftrightarrow 5^x>5^{-2}\Leftrightarrow x>-2.$$
		Vậy tập nghiệm của bất phương trình là $(-2;+\infty)$.
	}
\end{ex}
\begin{ex}%[DCHT Toán 11 - KNTT -Nguyễn Hoài Nam]%[1K6BK-3] 
	Tập nghiệm của bất phương trình $3^{x^2-23}<9$ là
	\choice
	{\True $(-5;5)$}
	{$(-\infty;5)$}
	{$(5;+\infty)$}
	{$(0;5)$}
	\loigiai{
		$3^{x^2-23}<9 \Leftrightarrow x^2-23 < 2 \Leftrightarrow x^2-25<0 \Leftrightarrow -5<x<5$.
	}
\end{ex}


\begin{ex}%[DCHT Toán 11 - KNTT -Nguyễn Hoài Nam]%[1K6BK-3] 
	Tập nghiệm của bất phương trình $\left(\dfrac{1}{2}\right)^{x-1}\ge 8$ là
	\choice
	{$[-3;+\infty)$}
	{$(-\infty;-3]$}
	{\True $(-\infty;-2]$}
	{$[-2;+\infty)$}
	\loigiai{
		$\left(\dfrac{1}{2}\right)^{x-1}\ge 8 \Leftrightarrow x-1 \le -3 \Leftrightarrow x\le -2$.
	}
\end{ex}


\begin{ex}%[DCHT Toán 11 - KNTT -Nguyễn Hoài Nam]%[1K6BK-3] 
	Tập nghiệm của bất phương trình $\log_3 (36-x^2)\ge 3$ là
	\choice
	{$(-\infty;-3]\cup[3;+\infty)$}
	{$(-\infty;3]$}
	{\True $[-3;3]$}
	{$(0;3]$}
	\loigiai{
		Điều kiện: $36-x^2 > 0 \Leftrightarrow -6<x<6$.\\
		$$\log_3 (36-x^2)\ge 3 \Leftrightarrow 36-x^2 \ge 27 \Leftrightarrow 9-x^2 \ge 0 \Leftrightarrow -3\le x\le 3.$$\\
		Kết hợp với điều kiện, ta có tập nghiệm $S=[-3;3]$.
	}
\end{ex}


\begin{ex}%[DCHT Toán 11 - KNTT -Nguyễn Hoài Nam]%[1K6BK-3] 
	Tập nghiệm của bất phương trình $\log_{\frac{1}{4}}(4x-2) + 1\ge 0$ là
	\choice
	{$\left[\dfrac{3}{2};+\infty\right)$}
	{$\left(\dfrac{1}{2};\dfrac{3}{2}\right)$}
	{$\left[\dfrac{1}{2};\dfrac{3}{2}\right]$}
	{\True $\left(\dfrac{1}{2};\dfrac{3}{2}\right]$}
	\loigiai{
		Điều kiện: $4x-2>0 \Leftrightarrow x>\dfrac{1}{2}$.
		$$ \log_{\frac{1}{4}}(4x-2) + 1\ge 0 \Leftrightarrow \log_{\frac{1}{4}}(4x-2) \ge -1 \Leftrightarrow 4x-2 \le 4 \Leftrightarrow x\le \dfrac{3}{2}.$$
		Kết hợp với điều kiện, ta có tập nghiệm $S=\left( \dfrac{1}{2}; \dfrac{3}{2} \right]$.
	}
\end{ex}
\begin{ex}%[DCHT Toán 11 - KNTT -Nguyễn Hoài Nam]%[1C6Y4-3]
	Giải bất phương trình $\log_2(3x-1)>3$. 
	\choice
	{\True $x>3$}
	{$\dfrac{1}{3}<x<3$}
	{$x<3$}
	{$x>\dfrac{10}{3}$}
	\loigiai{
		Ta có $\log_2(3x-1)>3\Leftrightarrow\log_2(3x-1)>\log_28\Leftrightarrow 3x-1>8\Leftrightarrow x>3$.}
\end{ex}
\begin{ex}%[DCHT Toán 11 - KNTT -Nguyễn Hoài Nam]%[1C6B4-3]
	Tìm tập nghiệm $S$ của bất phương trình $\log_2(3x-2)>\log_2(6-5x)$. 
	\choice
	{\True $S=\left(1;\dfrac{6}{5}\right)$}
	{$S=\left(\dfrac{2}{3};1\right)$}
	{$S=(1;+\infty)$}
	{$S=\left(\dfrac{2}{3};\dfrac{6}{5}\right)$}
	\loigiai{
		Điều kiện: $\heva{&3x-2>0\\&6-5x>0}\Leftrightarrow\dfrac{2}{3}<x<\dfrac{6}{5}$. \\
		Ta có $\log_2(3x-2)>\log_2(6-5x)\Leftrightarrow 3x-2>6-5x\Leftrightarrow x>1$.\\
		Kết hợp điều kiện suy ra $ 1<x<\dfrac{6}{5}\Rightarrow S=\left(1;\dfrac{6}{5}\right)$.}
\end{ex}
\begin{ex}%[DCHT Toán 11 - KNTT -Nguyễn Hoài Nam]%[1C6B4-3]
	Tìm tập nghiệm T của bất phương trình $\log_{\tfrac{1}{4}}(4x-2)\geq-1$. 
	\choice
	{$\left[\dfrac{3}{2};+\infty\right)$}
	{$\left(\dfrac{1}{2};\dfrac{3}{2}\right)$}
	{$\left[\dfrac{1}{2};\dfrac{3}{2}\right]$}
	{\True $\left(\dfrac{1}{2};\dfrac{3}{2}\right]$}
	\loigiai{
		Bất phương trình tương đương $0<4x-2\leq\left(\dfrac{1}{4}\right)^{-1}\Leftrightarrow\dfrac{1}{2}<x\leq\dfrac{3}{2}$.\\
		Vậy tập nghiệm là $T=\left(\dfrac{1}{2};\dfrac{3}{2}\right]$.}
\end{ex}
\begin{ex}%[DCHT Toán 11 - KNTT -Nguyễn Hoài Nam]%[1C6B4-3]
	Có bao nhiêu số nguyên trên $[0;10]$ nghiệm đúng bất phương trình ${\log_e(3x-4)>\log_e(x-1)}$?
	\choice
	{10}
	{11}
	{\True 9}
	{8}
	\loigiai{
		Ta có $\log_e(3x-4)>\log_e(x-1)\Leftrightarrow\heva{&x-1>0\\&3x-4>x-1}\Leftrightarrow\heva{&x>1\\&x>\dfrac{3}{2}}\Leftrightarrow x>\dfrac{3}{2}$. \\
		Mà $\heva{&x\in\mathbb{Z}\\&x\in [0;10]}$ nên $x\in\{2,3,4,5,6,7,8,9,10\}$.\\
		Vậy có 9 số.}
\end{ex}
\begin{ex}%[DCHT Toán 11 - KNTT -Nguyễn Hoài Nam]%[1K6BK-5]
	Tập nghiệm của bất phương trình $3^{2x-1}>27$ là
	\choice
	{\True $(2;+\infty)$}
	{$(3;+\infty)$}
	{$\left( \dfrac{1}{3};+\infty \right)$}
	{$\left( \dfrac{1}{2};+\infty \right)$}
	\loigiai{
		Ta có $3^{2x-1}>27 \Leftrightarrow 2x-1>3 \Leftrightarrow x>2$.
	}
\end{ex}

\begin{ex}%[DCHT Toán 11 - KNTT -Nguyễn Hoài Nam]%[1K6BK-4]
	Tìm tập nghiệm của bất phương trình $4^{2x+5}<2^{2-x}$.
	\choice
	{\True $\left(-\infty;-\dfrac{8}{5}\right)$}
	{$\left(-\infty;-\dfrac{8}{5}\right]$}
	{$\left(-\infty;\dfrac{8}{5}\right)$}
	{$\dfrac{8}{5}$}
	\loigiai{
		Ta có
		$$4^{2x+5}<2^{2-x}\Leftrightarrow 2^{2(2x+5)}<2^{2-x}\Leftrightarrow 2(2x+5)<2-x\Leftrightarrow x<-\dfrac{8}{5}.$$
		Vậy tập nghiệm của phương trình là $S=\left(-\infty;-\dfrac{8}{5}\right)$.		
	}
\end{ex}

\begin{ex}%[DCHT Toán 11 - KNTT -Nguyễn Hoài Nam]%[1K6BK-4]
	Tập nghiệm của bất phương trình $ \log_2 x>\log_2 {(x^2-x)} $ là
	\choice
	{\True $ (1;2) $}
	{$ (0;2) $}
	{$ \{0;2\} $}
	{$ \{2\} $}
	\loigiai{Điều kiện xác định của phương trình là $ \heva{&
			x>0 \\ &x^2-x >0 }
		\Leftrightarrow x>1$.\\
		Bất phương trình đã cho tương đương  $x>x^2-x \Leftrightarrow x^2 -2x< 0 \Leftrightarrow 0<x<2$\\
		Kết hợp với điều kiện, tập nghiệm của bất phương trình là $ S=(1;2) $.
	}
\end{ex}
\begin{ex}%[DCHT Toán 11 - KNTT -Nguyễn Hoài Nam]%[1K6BK-5]
	Giải  bất  phương trình $\log_2(3x-2)>\log_2(6-5x)$ được tập nghiệm là $(a;b)$. Hãy tính   tổng $S=a+b$.
	\choice
	{$S=\dfrac{26}{5}$}
	{$S=\dfrac{8}{3}$}
	{$S=\dfrac{28}{15}$}
	{\True $S=\dfrac{11}{5}$}
	\loigiai{
		$\log_2(3x-2)>\log_2(6-5x)\Leftrightarrow \heva{&3x-2>0\\&6-5x>0\\&3x-2>6-5x}\Leftrightarrow \heva{&x>\dfrac{2}{3}\\&x<\dfrac{6}{5}\\&x>1}\Leftrightarrow 1<x<\dfrac{6}{5}\Rightarrow S=1+\dfrac{6}{5}=\dfrac{11}{5}.$	
	}
\end{ex}
\begin{ex}%[DCHT Toán 11 - KNTT -Nguyễn Hoài Nam]%[1K6BK-3]
	Tập nghiệm của bất phương trình $3^{x^2-23}<9$ là
	\choice
	{\True $(-5;5)$}
	{$(-\infty;5)$}
	{$(5;+\infty)$}
	{$(0;5)$}
	\loigiai{
		Ta có $3^{x^2-23}<9\Leftrightarrow 3^{x^2-23}<3^2\Leftrightarrow x^2-23<2\Leftrightarrow x^2<25\Leftrightarrow -5<x<5$.\\
		Vậy tập nghiệm của bất phương trình là $S=(-5;5)$. }
\end{ex}
\begin{ex}%[DCHT Toán 11 - KNTT -Nguyễn Hoài Nam]%[1K6BK-3]
	Bất phương trình $\left(\dfrac{1}{3}\right)^{x^2-4x-12}>1$ có tất cả bao nhiêu nghiệm nguyên?
	\choice
	{$3$}
	{$5$}
	{\True $7$}
	{Vô số}
	\loigiai{
		BPT $\Leftrightarrow x^2-4x-12<0\Leftrightarrow -2<x<6$.\\
		Vì $x\in\mathbb{Z}$ nên $x\in\lbrace -1; 0; 1; 2; 3; 4; 5\rbrace$.}
\end{ex}
\begin{ex}%[DCHT Toán 11 - KNTT -Nguyễn Hoài Nam]%[1K6BK-3]
	Tập nghiệm của bất phương trình $\log_2{(2x+3)}\ge 0$ là
	\choice
	{$S=(-\infty;-1]$}
	{\True $S=[-1;+\infty)$}
	{$S=(-\infty;-1)$}
	{$S=(-\infty;0]$}
	\loigiai{
		Điều kiện: $2x+3>0\Leftrightarrow x>-\dfrac{3}{2}$.\\
		Ta có $\log_2{(2x+3)}\ge 0\Leftrightarrow 2x+3\ge 1\Leftrightarrow x\ge -1$.\\
		Vậy tập nghiệm của bất phương trình $S=[-1;+\infty)$. }
\end{ex}

\begin{ex}%[DCHT Toán 11 - KNTT -Nguyễn Hoài Nam]%[1K6BK-3]
	Số nghiệm nguyên thuộc đoạn $[0;10]$ của bất phương trình $7^{\sqrt{x+6}}\ge 7^x$ là
	\choice
	{$3$}
	{\True $4$}
	{$11$}
	{$10$}
	\loigiai{
		Điều kiện: $x\ge -6$.\\
		Ta có $7^{\sqrt{x+6}}\ge 7^x\Leftrightarrow \sqrt{x+6}\ge x\Leftrightarrow\hoac{&x\le 0\\&\heva{&x>0\\&x-6\ge x^2}}\Leftrightarrow\hoac{&x\le 0\\&\heva{&x>0\\&-2\le x\le 3}}\Leftrightarrow\hoac{&x\le 0\\&0<x\le 3}$.\\
		Khi đó các nghiệm nguyên thuộc đoạn $[0;10]$ là $x\in\lbrace 0; 1;2;3\rbrace$.}
\end{ex}

\Closesolutionfile{ans}
\begin{indapan}{10}
{ans/ans-1K6-21-Dang4}
\end{indapan}

\begin{dang}{Phương trình mũ, phương trình logarit đưa về cùng cơ số}
	\begin{itemize}
		\item Phương pháp giải phương trình mũ đưa về cùng cơ số $a^{u(x)} = a^{v(x)}$ ($0<a \ne 1$):
		$$a^{u(x)} = a^{v(x)} \Leftrightarrow u(x) = v(x) \Leftrightarrow \ldots$$
		\item Phương pháp giải phương trình logarit đưa về cùng cơ số: $\log_a u(x) = \log_a v(x)$ ($0<a \ne 1$):
		\begin{enumEX}[\faCheckCircleO]{1}
			\item Tìm điều kiện xác định: $\begin{cases}
				u(x) > 0 \\ v(x) > 0
			\end{cases}$.
			\item Biến đổi tương đương phương trình
			$$\log_a u(x) = \log_a v(x) \Leftrightarrow u(x) = v(x) \Leftrightarrow \ldots$$
			\item Kiểm tra nghiệm theo điều kiện xác định và kết luận.
		\end{enumEX}
	\end{itemize}
\end{dang}
\subsubsection{Ví dụ mẫu}
\begin{vd}[NB]%[DCHT Toán 11 - KNTT -Phong Lê] %[1K6YK-4]
	Giải phương trình $3^x = 3^{2-x}$.
	\dapso{$S = \{1\}$}
	\loigiai{$$3^x = 3^{2-x}
		\Leftrightarrow x = 2-x
		\Leftrightarrow 2x=2
		\Leftrightarrow x=1.$$
		Vậy tập nghiệm của phương trình là $S = \{1\}$.}
\end{vd}
\begin{vd}[TH]%[DCHT Toán 11 - KNTT -Phong Lê] %[1K6BK-4]
	Giải phương trình $4^x = 2^{x^2-2x}$.
	\dapso{$S = \{0;4\}$}
	\loigiai{$$4^x = 2^{x^2-2x} \Leftrightarrow \left(2^2\right)^x = 2^{x^2-2x} \Leftrightarrow2^{2x} = 2^{x^2-2x} \Leftrightarrow 2x = x^2-2x \Leftrightarrow x^2-4x = 0 \Leftrightarrow \left[\begin{matrix}
			x=0\\x=4
		\end{matrix}\right.. $$
		Vậy tập nghiệm của phương trình là $S = \{0;4\}$.}
\end{vd}
\begin{vd}[TH]%[DCHT Toán 11 - KNTT -Phong Lê] %[1K6BK-4]
	Giải phương trình $\dfrac{1}{5^{1+x}} = 25^{-x^2}$.
	\dapso{$S = \left\{1;\dfrac{1}{2}\right\}$}
	\loigiai{$$ \dfrac{1}{5^{1+x}} = 25^{-x^2} \Leftrightarrow \left(5^{-1}\right)^{1+x} = \left(5^2\right)^{-x^2}\Leftrightarrow 5^{-x-1} = 5^{-2x^2}\Leftrightarrow -x-1=-2x^2 \Leftrightarrow 2x^2-x-1=0 \Leftrightarrow \left[\begin{matrix}
			x=1\\ x= -\dfrac{1}{2}
		\end{matrix}\right.. $$
		Vậy tập nghiệm của phương trình là $S = \left\{1;\dfrac{1}{2}\right\}$.}
\end{vd}
\begin{vd}[NB]%[DCHT Toán 11 - KNTT -Phong Lê] %[1K6YK-4]
	Giải phương trình $\log_3\left(x^2\right)= \log_3(-x+2)$.
	\dapso{$S = \{1;-2\}$}
	\loigiai{
		Điều kiện xác định: $\begin{cases}
			x^2>0\\ -x+2>0
		\end{cases} \Leftrightarrow \begin{cases}
			x \ne 0\\ x<2
		\end{cases}$.\\
		$\log_3\left(x^2\right) = \log_3(-x+2) \Leftrightarrow x^2 = -x+2 \Leftrightarrow x^2+x-2 = 0 \Leftrightarrow \left[\begin{matrix}
			x=1\\ x= -2
		\end{matrix}\right..$\\
		Cả hai nghiệm trên đều thỏa mãn điều kiện. Vậy tập nghiệm của phương trình là $S = \{1;-2\}$.}
\end{vd}
\begin{vd}[TH]%[DCHT Toán 11 - KNTT -Phong Lê] %[1K6BK-4]
	Giải phương trình $\log_2x + \log_2(x+2)= \log_2(3x+2)$.
	\dapso{$S = \{2\}$}
	\loigiai{
		Điều kiện xác định: $\begin{cases}
			x>0\\ x+2>0\\3x+2>0
		\end{cases} \Leftrightarrow \begin{cases}
			x > 0\\ x>-2\\x>-\dfrac{2}{3}
		\end{cases} \Leftrightarrow x>0$.\\
		$ \log_2x + \log_2(x+2)= \log_2(3x+2) \Leftrightarrow \log_2\left[x(x+2)\right] = \log_2(3x+2)\\
		\Leftrightarrow x(x+2) = 3x+2 \Leftrightarrow x^2-x-2 = 0 \Leftrightarrow \left[\begin{matrix}
			x=-1\\ x= 2
		\end{matrix}\right..$\\
		Chỉ có nghiệm $x=2$ thỏa mãn điều kiện. Vậy tập nghiệm của phương trình là $S = \{2\}$.}
\end{vd}
\begin{vd}[TH]%[DCHT Toán 11 - KNTT -Phong Lê] %[1K6BK-4]
	Giải phương trình $\log_2\left(x-1\right)= \log_4(3x-3)$.
	\dapso{$S = \{5\}$}
	\loigiai{
		Điều kiện xác định: $\begin{cases}
			x-1>0\\ 3x-3>0
		\end{cases} \Leftrightarrow \begin{cases}
			x > 1\\ x>1
		\end{cases}\Leftrightarrow x>1$.\\
		$\log_2(x-1) = \log_4(3x-3) \Leftrightarrow \log_4\left[(x-1)^2\right] = \log_4(3x-3) \\ \Leftrightarrow (x-1)^2 = 3x-3 \Leftrightarrow x^2-5x+4=0 \Leftrightarrow \left[\begin{matrix}
			x=1\\ x= 5
		\end{matrix}\right..$\\
		Chỉ có nghiệm $x=5$ thỏa mãn điều kiện. Vậy tập nghiệm của phương trình là $S = \{5\}$.}
\end{vd}
\subsubsection{Bài tập rèn luyện}
\centerline{\fcolorbox{red}{yellow!50}{\bf {BÀI TẬP TỰ LUẬN}}}
\begin{bt}[NB]%[DCHT Toán 11 - KNTT -Phong Lê] %[1K6YK-4]
	Giải phương trình $5^{2x+1} = 5^{x^2+2}$.
	\dapso{$S = \{1\}$}
	\loigiai{$$5^{2x+1} = 5^{x^2+2}
		\Leftrightarrow 2x+1=x^2+2
		\Leftrightarrow x^2-2x+1=0
		\Leftrightarrow x=1.$$
		Vậy tập nghiệm của phương trình là $S = \{1\}$.}
\end{bt}
\begin{bt}[NB]%[DCHT Toán 11 - KNTT -Phong Lê] %[1K6YK-4]
	Giải phương trình $\log_2\left(x^2+1\right)= \log_2(x+7)$.
	\dapso{$S = \{3;-2\}$}
	\loigiai{
		Điều kiện xác định: $\begin{cases}
			x^2+1>0\\ x+7>0
		\end{cases} \Leftrightarrow x>-7$.\\
		$\log_2\left(x^2+1\right) = \log_2(x+7) \Leftrightarrow x^2+1 = x+7 \Leftrightarrow x^2-x-6 = 0 \Leftrightarrow \left[\begin{matrix}
			x=3\\ x= -2
		\end{matrix}\right..$\\
		Cả hai nghiệm trên đều thỏa mãn điều kiện. Vậy tập nghiệm của phương trình là $S = \{3;-2\}$.}
\end{bt}
\begin{bt}[TH]%[DCHT Toán 11 - KNTT -Phong Lê] %[1K6BK-4]
	Giải phương trình $2^{2x^3+x} = 4^{x^3+2}$.
	\dapso{$S = \{4\}$}
	\loigiai{$$2^{2x^3+x} = 4^{x^3+2}
		\Leftrightarrow 2^{2x^3+x} = \left(2^2\right)^{x^3+2}
		\Leftrightarrow 2^{2x^3+x} = 2^{2x^3+4}
		\Leftrightarrow 2x^3+x=2x^3+4
		\Leftrightarrow x=4.$$
		Vậy tập nghiệm của phương trình là $S = \{4\}$.}
\end{bt}
\begin{bt}[TH]%[DCHT Toán 11 - KNTT -Phong Lê] %[1K6BK-4]
	Giải phương trình $\left(\dfrac{8}{27}\right)^{x^2-1} = \left(\dfrac{9}{4}\right)^{x+1}$.
	\dapso{$S = \left\{-1;\dfrac{1}{3}\right\}$.}
	\loigiai{$\left(\dfrac{8}{27}\right)^{x^2-1} = \left(\dfrac{9}{4}\right)^{x+1}
		\Leftrightarrow \left[\left(\dfrac{2}{3}\right)^3\right]^{x^2-1} = \left[\left(\dfrac{2}{3}\right)^{-2}\right]^{x+1}
		\Leftrightarrow \left(\dfrac{2}{3}\right)^{3x^2-3} = \left(\dfrac{2}{3}\right)^{-2x-2}$\\
		$\Leftrightarrow 3x^2-3=-2x-2
		\Leftrightarrow 3x^2+2x-1=0
		\Leftrightarrow \left[\begin{matrix}
			x=-1\\ x= \dfrac{1}{3}
		\end{matrix}\right..$\\
		Vậy tập nghiệm của phương trình là $S = \left\{-1;\dfrac{1}{3}\right\}$.}
\end{bt}
\begin{bt}[TH]%[DCHT Toán 11 - KNTT -Phong Lê] %[1K6BK-4]
	Giải phương trình $\log_3\left(x^3+x\right)= \log_3\left(x^3+x^2-1\right)$.
	\dapso{$S = \left\{\dfrac{1+\sqrt{5}}{2}\right\}$}
	\loigiai{
		Điều kiện xác định: $\begin{cases}
			x^3+x>0\\ x^3+x^2-1>0
		\end{cases}$.\\
		$\log_3\left(x^3+x\right)= \log_3\left(x^3+x^2-1\right) \Leftrightarrow x^3+x = x^3+x^2-1 \Leftrightarrow x^2-x-1 = 0 \Leftrightarrow \left[\begin{matrix}
			x=\dfrac{1+\sqrt{5}}{2}\\ x= \dfrac{1-\sqrt{5}}{2}
		\end{matrix}\right..$\\
		Thử lại, chỉ có nghiệm $x = \dfrac{1+\sqrt{5}}{2}$ thỏa mãn điều kiện.\\
		Vậy tập nghiệm của phương trình là $S = \left\{\dfrac{1+\sqrt{5}}{2}\right\}$.}
\end{bt}
\begin{bt}[TH]%[DCHT Toán 11 - KNTT -Phong Lê] %[1K6BK-4]
	Giải phương trình $\log (3x+1) + \log (2x+4)= \log \left(7x^2+37\right)$.
	\dapso{$S = \{3;11\}$}
	\loigiai{
		Điều kiện xác định: $\begin{cases}
			3x+1>0\\ 2x+4>0\\7x^2+37>0
		\end{cases} \Leftrightarrow \begin{cases}
			x > -\dfrac{1}{3}\\ x>-2
		\end{cases} \Leftrightarrow x>-\dfrac{1}{3}$.\\
		$\log (3x+1) + \log (2x+4)= \log \left(7x^2+37\right)  \Leftrightarrow \log\left[(3x+1)(2x+4)\right] =  \log \left(7x^2+37\right)$ \\ $\Leftrightarrow (3x+1)(2x+4) = 7x^2+37 \Leftrightarrow x^2-14x+33 = 0 \Leftrightarrow \left[\begin{matrix}
			x=3\\ x= 11
		\end{matrix}\right..$\\
		Cả hai nghiệm thỏa mãn điều kiện. Vậy tập nghiệm của phương trình là $S = \{3;11\}$.}
\end{bt}
\begin{bt}[TH]%[DCHT Toán 11 - KNTT -Phong Lê] %[1K6BK-4]
	Giải phương trình $\log_2\left(x+2\right)= \log_4\left(4+6x\right)$.
	\dapso{$S = \{0;2\}$}
	\loigiai{
		Điều kiện xác định: $\begin{cases}
			x+2>0\\ 4+6x>0
		\end{cases}\Leftrightarrow \begin{cases}
			x>-2\\ x>-\dfrac{2}{3}
		\end{cases} \Leftrightarrow x>-\dfrac{2}{3}$. \\
		$\log_2\left(x+2\right)= \log_4\left(4+6x\right) \Leftrightarrow \log_4\left[(x+2)^2\right]= \log_4\left(4+6x\right)$\\ $\Leftrightarrow (x+2)^2 = 4+6x \Leftrightarrow x^2-2x = 0 \Leftrightarrow \left[\begin{matrix}
			x=0\\ x= 2
		\end{matrix}\right..$\\
		Cả hai nghiệm đều thỏa mãn điều kiện. Vậy tập nghiệm của phương trình là $S = \{0;2\}$.}
\end{bt}
\centerline{\fcolorbox{red}{yellow!50}{\bf {CÂU HỎI TRẮC NGHIỆM}}}
\Opensolutionfile{ans}[ans/ans-1K6-21-Dang5]
\begin{ex}[NB]%[DCHT Toán 11 - KNTT -Phong Lê] %[1K6YK-4]
	Nghiệm của phương trình $3^{x-2} = 3^{4-x}$ là
	\choice{\True $x=3$}{$x=4$}{$x=2$}{$x=6$}
	\loigiai{$$3^{x-2} = 3^{4-x}
		\Leftrightarrow x-2=4-x
		\Leftrightarrow 2x=6
		\Leftrightarrow x=3.$$}
\end{ex}
\begin{ex}[NB]%[DCHT Toán 11 - KNTT -Phong Lê] %[1K6YK-4]
	Tập nghiệm của phương trình $7^{2x^2-2} = 7^{x^2+2}$ là
	\choice{\True $\{2;-2\}$}{$\{-4\}$}{$\{4\}$}{$\{4;-4\}$}
	\loigiai{$$7^{2x^2-2} = 7^{x^2+2}
		\Leftrightarrow 2x^2-2=x^2+2
		\Leftrightarrow x^2=4
		\Leftrightarrow \left[\begin{matrix}
			x=2\\x=-2
		\end{matrix}\right..$$
		Vậy tập nghiệm của phương trình là $S=\{2;-2\}$.}
\end{ex}
\begin{ex}[NB]%[DCHT Toán 11 - KNTT -Phong Lê] %[1K6YK-4]
	Tập nghiệm của phương trình $\mathrm{e}^{x^2+4x} = \mathrm{e}^{x^2+5x+5}$ là
	\choice{\True $\{-5\}$}{$\{-4\}$}{$\{-5;-4\}$}{$\{0;-4\}$}
	\loigiai{$$\mathrm{e}^{x^2+4x} = \mathrm{e}^{x^2+5x+5} \Leftrightarrow x^2+4x = x^2+5x+5 \Leftrightarrow x= -5.$$
		Vậy tập nghiệm của phương trình là $S=\{-5\}$.}
\end{ex}
\begin{ex}[NB]%[DCHT Toán 11 - KNTT -Phong Lê] %[1K6YK-4]
	Số nghiệm của phương trình $\left(\dfrac{2}{5}\right)^{x^2-2} = \left(\dfrac{2}{5}\right)^{4x-6}$ là
	\choice{0}{\True 1}{2}{3}
	\loigiai{$$\left(\dfrac{2}{5}\right)^{x^2-2} = \left(\dfrac{2}{5}\right)^{4x-6}
		\Leftrightarrow x^2-2=4x-6
		\Leftrightarrow x^2-4x+4=0
		\Leftrightarrow x=2.$$
		Vậy phương trình có nghiệm duy nhất $x=2$.}
\end{ex}
\begin{ex}[NB]%[DCHT Toán 11 - KNTT -Phong Lê] %[1K6YK-4]
	Tập nghiệm của phương trình $\log_2\left(x^2+3x-1\right)= \log_2(4x-1)$ là
	\choice{\True $\{1\}$}{$\{0\}$}{$\{0;1\}$}{$\varnothing$}
	\loigiai{
		Điều kiện xác định: $\begin{cases}
			x^2+3x-1>0\\ 4x-1>0
		\end{cases}$.\\
		$\log_2\left(x^2+3x-1\right) = \log_2(4x-1) \Leftrightarrow x^2+3x-1 = 4x-1 \Leftrightarrow x^2-x = 0 \Leftrightarrow \left[\begin{matrix}
			x=0\\ x= 1
		\end{matrix}\right..$\\
		Chỉ có nghiệm $x=1$ thỏa mãn điều kiện. Vậy tập nghiệm của phương trình là $S = \{1\}$.}
\end{ex}
\begin{ex}[NB]%[DCHT Toán 11 - KNTT -Phong Lê] %[1K6YK-4]
	Tập nghiệm của phương trình $\log_2\left(x^2-2x-1\right)= \log_2(x^2-x-3)$ là
	\choice{\True $\varnothing$}{$\{2\}$}{$\{0;2\}$}{$\{0\}$}
	\loigiai{
		Điều kiện xác định: $\begin{cases}
			x^2-2x-1>0 \\ x^2-x-3>0
		\end{cases}$.\\
		$\log_2\left(x^2-2x-1\right)= \log_2(x^2-x-3) \Leftrightarrow x^2-2x-1 = x^2-x-3 \Leftrightarrow x = 2.$\\
		Đối chiếu điều kiện, $x=2$ không thỏa mãn điều kiện. Vậy tập nghiệm của phương trình là $S = \varnothing$.}
\end{ex}
\begin{ex}[NB]%[DCHT Toán 11 - KNTT -Phong Lê] %[1K6YK-4]
	Tập nghiệm của phương trình $\ln (-x) = \ln \left(x^2\right)$ là
	\choice{$\varnothing$}{$\{0\}$}{\True $\{-1\}$}{$\{0;-1\}$}
	\loigiai{
		Điều kiện xác định: $\begin{cases}
			-x>0 \\ x^2>0
		\end{cases}\Leftrightarrow x<0$.\\
		$\ln (-x) = \ln \left(x^2\right) \Leftrightarrow -x = x^2 \Leftrightarrow x^2+x=0 \Leftrightarrow \left[\begin{matrix}
			x = 0 \\ x= -1
		\end{matrix}\right..$\\
		Chỉ có $x=-1$ thỏa mãn điều kiện. Vậy tập nghiệm của phương trình là $S =\{-1\}$.}
\end{ex}
\begin{ex}[NB]%[DCHT Toán 11 - KNTT -Phong Lê] %[1K6YK-4]
	Phương trình $\log_2\left(x-4\right)= \log_2\left(x^2+2x-6\right)$ có bao nhiêu nghiệm?
	\choice{3}{2}{1}{\True 0}
	\loigiai{
		Điều kiện xác định: $\begin{cases}
			x-4>0 \\ x^2+2x-6>0
		\end{cases}\Leftrightarrow \begin{cases}
			x>4 \\ x^2+2x-6>0
		\end{cases} $.\\
		$\log_2\left(x-4\right)= \log_2\left(x^2+2x-6\right) \Leftrightarrow x-4 = x^2+2x-6 \Leftrightarrow x^2+x-2 = 0 \Leftrightarrow \left[\begin{matrix}
			x=1\\ x= -2
		\end{matrix}\right..$\\
		Có $x=1$ và $x=-2$ không mãn điều kiện. Vậy phương trình vô nghiệm.}
\end{ex}
\begin{ex}[TH]%[DCHT Toán 11 - KNTT -Phong Lê] %[1K6BK-4]
	Tập nghiệm của  phương trình $6^{x+1} = 36^{x-1}$ là
	\choice{$\{0\}$}{$\varnothing$}{\True $\{3\}$}{$\{0;3\}$}
	\loigiai{$$6^{x+1} = 36^{x-1}
		\Leftrightarrow 6^{x+1} = \left(6^2\right)^{x-1}
		\Leftrightarrow 6^{x+1} = 6^{2x-2}
		\Leftrightarrow x+1=2x-2
		\Leftrightarrow x=3.$$
		Vậy tập nghiệm của phương trình là $S = \{3\}$.}
\end{ex}
\begin{ex}[TH]%[DCHT Toán 11 - KNTT -Phong Lê] %[1K6BK-4]
	Tập nghiệm của phương trình $\left(\dfrac{3}{4}\right)^{x+5} = \left(\dfrac{4}{3}\right)^{-x^2+x+10}$ là
	\choice{\True $\{5;-3\}$}{$\{5;3\}$}{$\{5\}$}{$\{3\}$}\
	\loigiai{$\left(\dfrac{3}{4}\right)^{x+5} = \left(\dfrac{4}{3}\right)^{-x^2+x+10}
		\Leftrightarrow \left(\dfrac{3}{4}\right)^{x+5} = \left[\left(\dfrac{3}{4}\right)^{-1}\right]^{-x^2+x+10}
		\Leftrightarrow \left(\dfrac{3}{4}\right)^{x+5} = \left(\dfrac{3}{4}\right)^{x^2-x-10}$\\
		$\Leftrightarrow x+5=x^2-x-10
		\Leftrightarrow x^2-2x-15=0
		\Leftrightarrow \left[\begin{matrix}
			x=5\\ x=-3
		\end{matrix}\right..$\\
		Vậy tập nghiệm của phương trình là $S = \{5;-3\}$.}
\end{ex}
\begin{ex}[TH]%[DCHT Toán 11 - KNTT -Phong Lê] %[1K6BK-4]
	Tập nghiệm của phương trình $\dfrac{3^x}{3^{-x}} =3^{x^2}$ là
	\choice{\True $\{2;0\}$}{$\{1;0\}$}{$\{0\}$}{$\{1\}$}\
	\loigiai{$\dfrac{3^x}{3^{-x}} =3^{x^2} \Leftrightarrow 3^{2x} = 3^{x^2} \Leftrightarrow 2x = x^2 \Leftrightarrow x^2-2x = 0 \Leftrightarrow \left[x = 0 \\x = 2\right..$\\
		Vậy tập nghiệm của phương trình là $S = \{0;2\}$.}
\end{ex}
\begin{ex}[TH]%[DCHT Toán 11 - KNTT -Phong Lê] %[1K6BK-4]
	Tập nghiệm của phương trình $\log_8\left(x^4+2x\right)= \log_8\left(x^4+x^2-3\right)$.
	\choice{$\{3;-1\}$}{$\{-1\}$}{\True $\{3\}$}{$\varnothing$}
	\loigiai{
		Điều kiện xác định: $\begin{cases}
			x^4+2x>0 \\ x^4+x^2-3>0
		\end{cases}$\\
		$\log_8\left(x^4+2x\right)= \log_8\left(x^4+x^2-3\right) \Leftrightarrow x^4+2x = x^4+x^2-3 \Leftrightarrow x^2-2x-3=0 \Leftrightarrow \left[\begin{matrix}
			x=3\\ x= -1
		\end{matrix}\right..$\\
		Thử lại, chỉ có nghiệm $x = 3$ thỏa mãn điều kiện.\\
		Vậy tập nghiệm của phương trình là $S = \{3\}$.}
\end{ex}
\begin{ex}[TH]%[DCHT Toán 11 - KNTT -Phong Lê] %[1K6BK-4]
	Số nghiệm của phương trình $\ln x + \ln (x+2)= 2\ln (x+1)$ là
	\choice{3}{2}{1}{\True 0}
	\loigiai{
		Điều kiện xác định: $\begin{cases}
			x>0\\ x+2>0\\x+1>0
		\end{cases} \Leftrightarrow \begin{cases}
			x > 0\\ x>-2\\x>-1
		\end{cases} \Leftrightarrow x>0$.\\
		$\ln x + \ln (x+2)= 2\ln (x+1) \Leftrightarrow \ln \left[x(x+2)\right]  = \ln \left[(x+1)^2\right] \Leftrightarrow x(x+2) = (x+1)^2 \Leftrightarrow 0 = 1$ (vô lý).\\
		Vậy phương trình vô nghiệm.}
\end{ex}
\begin{ex}[TH]%[DCHT Toán 11 - KNTT -Phong Lê] %[1K6BK-4]
	Tập nghiệm của phương trình $\log \left(x^2+2\right) + \log (x+8)= 3\log (x+2)$ là
	\choice{\True $\{1;4\}$}{$\{1\}$}{$\{4\}$}{$\varnothing$}
	\loigiai{
		Điều kiện xác định: $\begin{cases}
			x^2+2>0\\ x+8>0\\x+2>0
		\end{cases} \Leftrightarrow \begin{cases}
			x>-8\\x>-2
		\end{cases} \Leftrightarrow x>-2$.\\
		$\log \left(x^2+2\right) + \log (x+8)= 3\log (x+2) \Leftrightarrow \log \left[\left(x^2+2\right)(x+8)\right] = \log \left[(x+2)^3\right] \Leftrightarrow \left(x^2+2\right)(x+8) = (x+2)^3 \Leftrightarrow 2x^2-10x+8 = 0 \Leftrightarrow \left[\begin{matrix}
			x = 4 \\ x= 1
		\end{matrix}\right..$\\
		Cả hai nghiệm đều thỏa mãn điều kiện. Vậy tập nghiệm của phương trình là $S = \{1;4\}.$}
\end{ex}
\begin{ex}[TH]%[DCHT Toán 11 - KNTT -Phong Lê] %[1K6BK-4]
	Giải phương trình $\log_2\left(x+2\right)= \log_8\left(x^3+6x+8\right)$.
	\choice{\True $\{0;-1\}$}{$\varnothing$}{$\{0\}$}{$\{-1\}$}
	\loigiai{
		Điều kiện xác định: $\begin{cases}
			x+2>0\\ x^3+6x+8>0
		\end{cases}\Leftrightarrow \begin{cases}
			x>-2\\ x^3+6x+8>0
		\end{cases}$. \\
		$\log_2\left(x+2\right)= \log_8\left(x^3+6x+8\right) \Leftrightarrow \log_8\left[(x+2)^3\right]= \log_8\left(x^3+6x+8\right)$\\ $\Leftrightarrow (x+2)^3 = x^3+6x+8 \Leftrightarrow 6x^2+6x = 0 \Leftrightarrow \left[\begin{matrix}
			x=0\\ x= -1
		\end{matrix}\right..$\\
		Cả hai nghiệm đều thỏa mãn điều kiện.\\
		Vậy tập nghiệm của phương trình là $S = \{0;-1\}$.}
\end{ex}
\Closesolutionfile{ans}
\begin{indapan}{10}
{ans/ans-1K6-21-Dang5}
\end{indapan}

\begin{dang}{Bất phương trình mũ, lôgarit đưa về cùng cơ số}
	
\end{dang}

\subsubsection{Ví dụ mẫu}
\begin{vd}%[Nguyễn Thành Nhân]%[Dự án 11 chương trình mới]%%[1K6BK-5]
Tìm nghiệm của bất phương trình $3^{x^2-23}<9$.
\loigiai{
Ta có $3^{x^2-23}<9\Leftrightarrow 3^{x^2-23}<3^2\Leftrightarrow x^2-23<2\Leftrightarrow x^2<25\Leftrightarrow -5<x<5$.\\
Suy ra, tập ngiệm của bất phương trình là $(-5;5)$.
}
\end{vd}

\begin{vd}%[Nguyễn Thành Nhân]%[Dự án 11 chương trình mới]%[1K6BK-5]
Giải bất phương trình $\log_2(3x-2)+\log_{\frac{1}{2}}(6-5x)>0$.
\loigiai{
Điều kiện $\dfrac{2}{3}<x<\dfrac{6}{5}$.\\
Bất phương trình đã cho tương đương với $\log_2(3x-2)>\log_2(6-5x)\Leftrightarrow 3-2x>6-5x\Leftrightarrow x>1$.\\
Kết hợp với điều kiện, tập nghiệm của bất phương trình là $S=\left(1;\dfrac{6}{5}\right)$.
}
\end{vd}

\begin{vd}%[Nguyễn Thành Nhân]%[Dự án 11 chương trình mới]%[1K6BK-5]
Tìm tập nghiệm của bất phương trình $(0{,}6)^{\tfrac{1}{x}}\leq (0{,}6)^{\tfrac{1}{6}}$.
\loigiai{
Ta có $(0{,}6)^{\tfrac{1}{x}}\leq (0{,}6)^{\tfrac{1}{6}}\Leftrightarrow\dfrac{1}{x}\geq\dfrac{1}{6}\Leftrightarrow 0<x\leq 6$.	\\
Vậy tập nghiệm của bất phương trình đã cho là $S=\left(0;6\right]$.
}
\end{vd}
\begin{vd}%[Nguyễn Thành Nhân]%[Dự án 11 chương trình mới]%[1K6BK-5]
Tìm tập nghiệm $S$ của bất phương trình $\ln x^2 > \ln (4x-4)$.
\loigiai
{Ta có
\[\ln x^2 > \ln (4x-4) \Leftrightarrow \left\{\begin{aligned}&4x-4 > 0 \\&x^2>4x-4\end{aligned}\right. \Leftrightarrow \left\{\begin{aligned}&x>1 \\&(x-2)^2>0\end{aligned}\right. \Leftrightarrow \left\{\begin{aligned}&x>1 \\&x\neq 2.\end{aligned}\right.\]
Vậy bất phương trình đã cho có tập nghiệm là $S=(1;+\infty)\setminus\{2\}$.
}
\end{vd}
\begin{vd}%[Nguyễn Thành Nhân]%[Dự án 11 chương trình mới]%[1K6BK-5]
Giải bất phương trình $\log_{\frac{1}{2}}\left(x^2-3x+2\right)\geq-1$.
\loigiai{ Bất phương trình đã cho tương đương $\heva{&x^2-3x+2 \le 2\\&x^2-3x+2>0 \\} \Leftrightarrow \heva{& 0 \le x \le 3\\&\hoac{&x>2\\&x<1} \\} \Leftrightarrow  \hoac{& 0 \le x <1 \\&2<x\le 3.}$}
\end{vd}


\begin{vd}%[Nguyễn Thành Nhân]%[Dự án 11 chương trình mới]%[1K6BK-5]
Tìm tập nghiệm của bất phương trình $\left(2-\sqrt{3}\right)^{x^2+4x-14} \geqslant 7+4\sqrt{3}$.
\loigiai{
Ta có $\left(2-\sqrt{3}\right)^{x^2+4x-14} \geqslant 7+4\sqrt{3}$
$ \Leftrightarrow x^2+4x-14 \leqslant \log_{\left(2-\sqrt{3}\right)}\left(7+4\sqrt{3}\right)$
$ \Leftrightarrow x^2+4x-14 \leqslant -2$ \\
$ \Leftrightarrow x^2+4x-12 \leqslant 0$
$ \Leftrightarrow x \in [-6;2]$.}
\end{vd}
\begin{vd}%%[Nguyễn Thành Nhân]%[Dự án 11 chương trình mới]%[1K6BK-5]
Bất phương trình $3\log_8(x+1)-\log_2(3-x)\le 1$ có bao nhiêu nghiệm nguyên?
\loigiai{
Điều kiện: $\heva{&x+1>0\\&3-x>0} \Leftrightarrow -1<x<3$.\\
Với điều kiện trên, ta có
\allowdisplaybreaks{
\begin{eqnarray*}
3\log_8(x+1)-\log_2(3-x)\le 1 &\Leftrightarrow& \log_2(x+1)-\log_2(3-x)\le 1\\
&\Leftrightarrow&\log_2\dfrac{x+1}{3-x}\le 1 \\
&\Leftrightarrow& \dfrac{x+1}{3-x}\le 2\\
&\Leftrightarrow& x+1\le 2(3-x)\\
&\Leftrightarrow& x\le\dfrac{5}{3}.
\end{eqnarray*}
}
Đối chiếu với điều kiện ta được $-1<x\le \dfrac{5}{3}$. Vậy phương trình có $2$ nghiệm nguyên.
}
\end{vd}

\begin{vd}%%[Nguyễn Thành Nhân]%[Dự án 11 chương trình mới]%[1K6KK-5]
Bất phương trình $\log_{\tfrac{1}{5}}(4x-3)>\dfrac{1}{2}\log_{\tfrac{1}{5}}(37-6x)^2$ có bao nhiêu nghiệm nguyên trong đoạn $[ -2020;2020 ]$?
\loigiai{
Điều kiện $ \heva{& 4x-3>0\\ &  37-6x\ne 0 }\Leftrightarrow\heva{&  x>\dfrac 34\\ & x\ne\dfrac{37}{6}.}$\\
Ta có
\begin{eqnarray*}
& &\log_{\tfrac{1}{5}}(4x-3)>\dfrac{1}{2}\log_{\tfrac {1}{5}}(37-6x)^2\\
&\Leftrightarrow &\log_{\tfrac{1}{5}}(4x-3)>\dfrac{1}{2}\cdot 2 \cdot\log_{\tfrac{1}{5}}| 37-6x|\\
&\Leftrightarrow&\log_{\tfrac{1}{5}}(4x-3)>\log_{\tfrac{1}{5}}| 37-6x|\\
&\Leftrightarrow& 4x-3<|37-6x|\\
&\Rightarrow&(4x-3)^2<(37-6x)^2\\
&\Leftrightarrow& 16x^2-24x+9<1369-444x+36x^2\\
&\Leftrightarrow& 20x^2-420x+1360>0\Leftrightarrow\hoac{ &x>17\\&x<4.}
\end{eqnarray*}
Mà $x>\dfrac{3}{4}$ và $x\in [ -2020;2020 ]$ nên $\hoac{&\dfrac{3}{4}<x<4 \\ & 17<x\le 2020 .} $\\
Vì $x\in \mathbb{Z}$ nên $\hoac{& x\in \{ 1;2;3 \} \\ & x\in \{ 18;19;20;\ldots;2020 \} }$, khi đó có tất cả $2006$ giá trị.}
\end{vd}
\begin{vd}%%[Nguyễn Thành Nhân]%[Dự án 11 chương trình mới]%[1K6GK-5]
Trong tất cả các cặp số thực $x$, $y$ thỏa mãn $\log_{\frac{1}{2}}(x+2y)-1\leq \log_{\frac{1}{2}}\left(x^2+y^2+1\right)$ chỉ có duy nhất một cặp số $(x;y)$ sao cho $x-2y-m=0$, $(m\in \mathbb{R})$. Tính tổng tất cả các giá trị của $ m$ thỏa mãn.
\loigiai{
Ta có
\begin{align*}
\log_{\frac{1}{2}}(x+2y)-1\leq \log_{\frac{1}{2}}\left(x^2+y^2+1\right)&\Leftrightarrow \log_2(2x+4y)\geq \log_2\left(x^2+y^2+1\right)\\
&\Leftrightarrow x^2+y^2+1\leq 2x+4y\\
&\Leftrightarrow (x-1)^2+(y-2)^2\leq 4.
\end{align*}
Chỉ có duy nhất một cặp số $(x;y)$ thỏa mãn yêu cầu bài toán khi và chỉ khi đường thẳng $\Delta\colon x-2y-m=0$ tiếp xúc với đường tròn tâm $I(1;2)$, bán kính $R=2 \Leftrightarrow d(I,\Delta)=R\Leftrightarrow \dfrac{|m+3|}{\sqrt{5}}=2\Leftrightarrow m=-3\pm 2\sqrt{5}$.\\
Vậy tổng tất cả các giá trị của $m$ thỏa mãn yêu cầu bài toán là $-6$.
}
\end{vd}
\begin{vd}%[Nguyễn Thành Nhân]%[Dự án 11 chương trình mới]%[1K6GK-5]
Cho dãy số $(u_n)$ thỏa mãn $\log u_2+\sqrt{2+\log u_2-2\log u_{10}}=2\log u_{10}$ và $u_{n+1}=3u_n$ với mọi $n\ge 1$. Giá trị nhỏ nhất của $n$ để $u_n>3^{2019}$ bằng bao nhiêu?
\loigiai{
\begin{itemize}
\item Ta có $u_{n+1}=3u_n,\,\forall n\ge 1$ nên $(u_n)$ là cấp số nhân có số hạng đầu là $u_1$ và công bội là $3$.\hfill{(1)}
\item Đặt $t=2\log u_{10}-\log u_2$, ta được
\[\sqrt{2-t}=t\Leftrightarrow\heva{&t\ge 0\\&2-t=t^2}\Leftrightarrow\heva{&t\ge 0\\&\hoac{&t=1\\&t=-2}}\Leftrightarrow t=1.\tag{2}\]
\item Từ $(1)$ và $(2)$, ta có
\begin{align*}2\log u_{10}-\log u_2=1&\Leftrightarrow 2\log u_1\cdot 3^9-\log u_1\cdot 3=1\\
&\Leftrightarrow 2\log u_1+2\log 3^9-\log u_1-\log 3=1\\
&\Leftrightarrow \log u_1=1-\log 3^{17}.\tag{3}
\end{align*}
\item Ta có
\begin{eqnarray*}u_n>3^{2019}&\Leftrightarrow& u_1\cdot 3^{n-1}>3^{2019}\\
&\Leftrightarrow& \log u_1+\log 3^{n-1}>\log 3^{2019}\\
&\Leftrightarrow&(n-1)\log 3>\log 3^{2019}-1+\log 3^{17}\,\,\text{do\, (3)}\\
&\Leftrightarrow& n-1>2019-\log_3 {10}+17\Leftrightarrow n>2037-\log_3{10}\approx 2034{,}9.
\end{eqnarray*}
\item Vậy giá trị nhỏ nhất của $n$ thỏa mãn đề bài là $2035$.
\end{itemize}
}
\end{vd}
\begin{vd}%[Nguyễn Thành Nhân]%[Dự án 11 chương trình mới]%[1K6GK-5]
Có bao nhiêu giá trị nguyên của tham số $m$ để tập nghiệm của bất phương trình $\ln(x^2+2x+m)-2\ln(2x-1)>0$ chứa đúng $2$ số nguyên?
\loigiai{
\begin{itemize}
\item Ta có
\begin{eqnarray*}
\ln(x^2+2x+m)-2\ln(2x-1)>0&\Leftrightarrow&\ln(x^2+2x+m)>2\ln(2x-1)\\
&\Leftrightarrow&\heva{&x>\dfrac{1}{2}\\&\ln(x^2+2x+m)>\ln(2x-1)^2}\\
&\Leftrightarrow&\heva{&x>\dfrac{1}{2}\\&x^2+2x+m>(2x-1)^2}\\
&\Leftrightarrow &\heva{&x>\dfrac{1}{2}\\&3x^2-6x-m+1<0.\quad\quad (1)}
\end{eqnarray*}
\item Bất phương trình $(1)$ có nghiệm khi tam thức $3x^2-6x-m+1$ có hai nghiệm phân biệt $\Leftrightarrow \Delta'>0\Leftrightarrow m>-2$.
\item Với $m>-2$, tam thức $3x^2-6x-m+1$ có hai nghiệm là $x_1=\dfrac{3+\sqrt{3m+6}}{3}>1$; $x_2=\dfrac{3-\sqrt{3m+6}}{3}$.
\item Kết hợp điều kiện $x>\dfrac{1}{2}$ để xác định tập nghiệm thỏa mãn đề bài, ta xét $2$ trường hợp:\\
TH1: $x_2<\dfrac{1}{2}\Leftrightarrow\sqrt{3m+6}>\dfrac{3}{2}\Leftrightarrow m>-\dfrac{5}{4}$, khi đó tập nghiệm của $(1)$ là
$\left(\dfrac{1}{2}; x_1\right)$.\\
Khi đó yêu cầu bài toán $\Leftrightarrow 2<x_1\leq 3\Leftrightarrow 3<\sqrt{3m+6}\leq 6\Leftrightarrow 1<m\leq 10$.\\
Kết hợp các điều kiện, ta có $9$ giá trị nguyên của $m$ thỏa mãn.\\
TH2: $x_2\geq\dfrac{1}{2}\Leftrightarrow\sqrt{3m+6}\leq\dfrac{3}{2}$. Khi đó tập nghiệm của $(1)$ là $\left(x_2; x_1\right)$.\\
Mà $x_1-x_2=\dfrac{2\sqrt{3m+6}}{3}\leq 1$ nên tập nghiệm không thể chứa hai giá trị nguyên.
\item Vậy có $9$ giá trị nguyên của tham số $m$ thỏa mãn yêu cầu bài toán.
\end{itemize}
}
\end{vd}
\subsubsection{Bài tập rèn luyện}
\centerline{\fcolorbox{red}{yellow!50}{\bf {BÀI TẬP TỰ LUẬN }}}
\begin{bt}%[Nguyễn Thành Nhân]%[Dự án 11 chương trình mới]%[1K6BK-5]
Số nghiệm nguyên của bất phương trình $7^{-x^2-5x+7}>\left(\dfrac{1}{7}\right)^{2x+3}$ là
\choice
{$3$}
{$2$}
{\True $6$}
{$8$}
\loigiai{
Ta thấy
\begin{eqnarray*}
7^{-x^2-5x+7}>\left(\dfrac{1}{7}\right)^{2x+3}
&\Leftrightarrow& 7^{-x^2-5x+7}>7^{-2x-3}
\Leftrightarrow-x^2-5x+7>-2x-3\\
&\Leftrightarrow&-x^2-3x+10>0
\Leftrightarrow -5<x<2.
\end{eqnarray*}
Vậy số nghiệm nguyên của bất phương trình là $6$.
}
\end{bt}
\begin{bt}%[Nguyễn Thành Nhân]%[Dự án 11 chương trình mới]%[1K6BK-5]
Tìm tập nghiệm của bất phương trình $\left(\sqrt{5}+2\right)^{x-1}\leq \left(\sqrt{5}-2\right)^{x-1}$.
\loigiai{
$(\sqrt{5}+2)(\sqrt{5}-2)=1 \Rightarrow \sqrt{5}-2=\dfrac{1}{\sqrt{5}+2}$. Ta có:
\begin{eqnarray*}
&&	\left(\sqrt{5}+2\right)^{x-1}\leq \left(\sqrt{5}-2\right)^{x-1}\\ &\Leftrightarrow&\left(\sqrt{5}+2\right)^{x-1}\leq
\left(\sqrt{5}+2\right)^{-x+1}\\
&\Leftrightarrow& x-1\leq -x+1\\
&\Leftrightarrow& x\leq 1.
\end{eqnarray*}
Vậy tập nghiệm của bất phương trình $\left(\sqrt{5}+2\right)^{x-1}\leq \left(\sqrt{5}-2\right)^{x-1}$ là $S=(-\infty; 1]$.
}
\end{bt}



\begin{bt}%[Nguyễn Thành Nhân]%[Dự án 11 chương trình mới]%[1K6BK-5]
Giải bất phương trình\, $2\log_3(4x-3)+\log_{\tfrac{1}{9}}(2x+3)^2\leq 2$.
\loigiai{ ĐKXĐ: $\heva{&4x-3>0\\ &2x+3\ne 0}$ $\Leftrightarrow x>\dfrac{3}{4}$.\\
Với $x>\dfrac{3}{4}$ thì $2x+3>0$ nên bất phương trình tương đương với
\begin{eqnarray*}
& & \log_3(4x-3)^2-\log_3(2x+3)\leq 2\\
&\Leftrightarrow & \log_3(4x-3)^2\leq \log_3(2x+3)+\log_39\\
&\Leftrightarrow & \log_3(4x-3)^2\leq \log_3[9(2x+3)]\\
&\Leftrightarrow & (4x-3)^2\leq 9(2x+3)\\
&\Leftrightarrow & 16x^2-42x-18\leq 0\\
&\Leftrightarrow & -\dfrac{3}{8}\leq x\leq 3.
\end{eqnarray*}
Kết hợp ĐKXĐ ta được\, $\dfrac{3}{4}<x\leq 3$.
}
\end{bt}

\begin{bt}%[Nguyễn Thành Nhân]%[Dự án 11 chương trình mới]%[1K6BK-5]
Tìm tập nghiệm $S$ của bất phương trình\, $\log_4(x-1)^2-\log_2(x+2)\leq 1$.
\loigiai{ ĐKXĐ: $\heva{& x>-2\\ & x\ne 1.}$ \\
Bất phương trình tương đương với
\begin{eqnarray*}
& & \log_4(x-1)^2-2\log_4(x+2)\leq \log_44\\
&\Leftrightarrow &\log_4(x-1)^2\leq \log_4\left[4(x+2)^2\right]\\
&\Leftrightarrow & (x-1)^2\leq 4(x+2)^2\\
&\Leftrightarrow &3x^2+18x+15\geq 0\\
&\Leftrightarrow &\hoac{& x\leq -5\\ & x\geq -1.}
\end{eqnarray*}
Kết hợp ĐKXĐ ta được tập nghiệm của bất phương trình là $S=[-1;1)\cup (1;+\infty)$.
}
\end{bt}
\begin{bt}%[Nguyễn Thành Nhân]%[Dự án 11 chương trình mới]%[1K6BK-5]
Tìm tập nghiệm của bất phương trình $4^x-6\cdot 2^x+8<0$.
\loigiai{
Ta có $4^x-6\cdot 2^x+8<0 \Leftrightarrow 2<2^x<4 \Leftrightarrow 1<x<2$.\\
Vậy tập nghiệm của bất phương trình là $S=(1;2)$.
}
\end{bt}

\begin{bt}%[Nguyễn Thành Nhân]%[Dự án 11 chương trình mới]%[1K6BK-5]
Có bao nhiêu số nguyên thỏa mãn bất phương trình $\log(x-40)+\log(60-x)<2$?
\loigiai{
Điều kiện $\heva{&x-40>0\\&60-x>0}\Leftrightarrow 40<x<60$. Ta có
\allowdisplaybreaks
{
\begin{eqnarray*}
&&\log(x-40)+\log(60-x)<2\\
&\Leftrightarrow& \log(x-40)(60-x)<2\\
&\Leftrightarrow& (x-40)(60-x)<10^2	\\
&\Leftrightarrow& x^2-100x+2500>0\Leftrightarrow x\neq 50.
\end{eqnarray*}
}
Vì $x\in (40;60)$ và $x\in \mathbb{Z}$ nên có $18$ số nguyên thỏa mãn.
}
\end{bt}
\begin{bt}%[Nguyễn Thành Nhân]%[Dự án 11 chương trình mới]%[1K6BK-5]
Cho các số thực dương $x, y$ thỏa mãn $\left(\dfrac{10}{9}\right)^{2x^2-5xy}\leq \left(\dfrac{3}{\sqrt{10}}\right)^{xy+5y^2}$. Tính hiệu giữa giá trị lớn nhất và giá trị nhỏ nhất của biểu thức $\dfrac{x}{y}$.
\loigiai{
Ta có \begin{eqnarray*}
\left(\dfrac{10}{9}\right)^{2x^2-5xy}\leq \left(\dfrac{3}{\sqrt{10}}\right)^{xy+5y^2}&\Leftrightarrow &\left(\dfrac{\sqrt{10}}{3}\right)^{2(2x^2-5xy)}\leq \left(\dfrac{\sqrt{10}}{3}\right)^{-xy-5y^2}\\
&\Leftrightarrow & 4x^2-10xy\leq -xy -5y^2\\
&\Leftrightarrow & 4x^2 -9xy +5y^2\leq 0\\
&\Leftrightarrow &4\left(\dfrac{x}{y}\right)^2 -9\left(\dfrac{x}{y}\right) +5\leq 0\\
&\Leftrightarrow & 1\leq \dfrac{x}{y}\leq \dfrac{5}{4}.
\end{eqnarray*}
Vậy $\min\dfrac{x}{y}=1$, $\max\dfrac{x}{y}=\dfrac{5}{4}$. Nên hiệu giữa giá trị lớn nhất và giá trị nhỏ nhất của biểu thức $\dfrac{x}{y}$ bằng $\dfrac{1}{4}$.
}
\end{bt}

\begin{bt}%[Nguyễn Thành Nhân]%[Dự án 11 chương trình mới]%[1K6KK-5]
Có tất cả bao nhiêu giá trị nguyên của tham số $ m $ thuộc khoảng $ (1;20) $ để $ \forall x\in\left( \dfrac{1}{3};1 \right) $ đều là nghiệm của bất phương trình $ \log_mx>\log_xm $?
\loigiai
{
Điều kiện xác định là $\heva{&x>0\\&x\ne 1.}$\\
Bất phương trình đã cho trở thành
\[\log_mx>\dfrac{1}{\log_mx}\Leftrightarrow \dfrac{(\log_mx)^2-1}{\log_mx}>0. \tag{$*$}\]
Do $ x\in\left( \dfrac{1}{3};1 \right) $ nên $ \log_mx<0 $. Suy ra
$$ (*)\Leftrightarrow-1<\log_mx<1\Leftrightarrow \dfrac{1}{m}<x<m.$$
Mọi $ x\in\left( \dfrac{1}{3};1 \right) $ đều là nghiệm của bất phương trình khi chỉ khi
$$\dfrac{1}{m}\le\dfrac{1}{3}<1\le m\Leftrightarrow m\ge 3.$$
Vì $m$ là số nguyên thuộc $(1;20)$ và $m\geq 3$ nên $m\in \{3;4;\ldots;19\}$.\\
Vậy có $17$ giá trị nguyên $m$ thuộc khoảng $(1;20)$ thỏa mãn yêu cầu bài toán.
}
\end{bt}

\begin{bt}%[Nguyễn Thành Nhân]%[Dự án 11 chương trình mới]%[1K6KK-5]
Có bao nhiêu giá trị nguyên của tham số $m$ để bất phương trình \\ $\log5+\log(x^2+1)\ge\log(mx^2+4x+m)$  đúng với mọi $x\in\mathbb{R}$?
\loigiai
{Vì bất phương trình đã cho đúng với mọi $x\in\mathbb{R}$ nên điều kiện của bất phương trình là
\[mx^2+4x+m>0,~\forall x\in\mathbb{R}\Leftrightarrow \heva{&m>0\\&4-m^2<0}\Leftrightarrow\heva{&m>0\\&m<-2\text{ hoặc }m>2}\Leftrightarrow m\in(2;+\infty).\]
Khi đó,
\allowdisplaybreaks
\begin{eqnarray*}
& &\log5+\log(x^2+1)\ge\log(mx^2+4x+m),~\forall x\in\mathbb{R}\\
&\Leftrightarrow & \log5(x^2+1)\ge\log(mx^2+4x+m),~\forall x\in\mathbb{R}\\
& \Leftrightarrow & 5(x^2+1)\ge mx^2+4x+m,~\forall x\in\mathbb{R}\\
& \Leftrightarrow & (5-m)x^2-4x+5-m\ge0,~\forall x\in\mathbb{R}\\
& \Leftrightarrow & \heva{&5-m>0\\&4-(5-m)^2\le0}\Leftrightarrow\heva{&m<5\\&m\le3\text{ hoặc }m\ge7}\Leftrightarrow m\le3.
\end{eqnarray*}
Kết hợp với điều kiện trên ta được $2<m\le3$. Suy ra $m=3$ thỏa mãn đề bài.
}
\end{bt}

\begin{bt}%[Nguyễn Thành Nhân]%[Dự án 11 chương trình mới]%[1K6BK-5]%
Cho $x,\,y$ là các số thực dương thỏa mãn $\log x+\log \left(20y\right)\ge 1+\log \left(x+16y^3\right)$. Tìm giá trị nhỏ nhất của biểu thức $P={\log_2}x-{\log_2}\left( 2y\right)$.
\loigiai{
Điều kiện $x>0,\,y>0$.
\allowdisplaybreaks
\begin{eqnarray*}
&&\log x+\log \left(20y\right)\ge 1+\log \left(x+16y^3\right)\\
&\Leftrightarrow& \log \left(20xy\right)\ge \log \left( 10\left(x+16y^3\right) \right)\\
&\Leftrightarrow& 20xy\ge 10\left(x+16y^3\right)\\
&\Leftrightarrow& 2xy\ge x+16y^3\\
&\Leftrightarrow& x\left(2y-1\right)\ge 16y^3. \quad (1)
\end{eqnarray*}
Vì $x,\,y>0$ $\Rightarrow 2y-1>0$, do đó (1) $\Leftrightarrow x\ge \dfrac{16y^3}{2y-1}$.\\
Khi đó $P={\log_2}x-{\log_2}\left(2y\right)={\log_2}\dfrac{x}{2y}\ge {\log_2}\dfrac{8y^2}{2y-2}$.\\
Xét hàm số $f\left(y\right)=\dfrac{8y^2}{2y-1}$ trên $\left(\dfrac{1}{2};+\infty\right)$
$\Rightarrow {f}'\left(y\right)=\dfrac{16y^2-16y}{\left(2y-1\right)^2} \Rightarrow f'\left(y\right)=0\Leftrightarrow y=1$.\\
Bảng biến thiên
\begin{center}
\begin{tikzpicture}[scale=1, font=\footnotesize, line join=round, line cap=round, >=stealth]
\tkzTabInit[lgt=1.5,espcl=3,deltacl=0.5]{$y$/1 ,$f'(y)$/.7,$f(y)$/2}
{$\dfrac{1}{2}$,$1$,$+\infty$}
\tkzTabLine{ ,-,0, +,}
\tkzTabVar{+/$+\infty$,-/$8$,+/$+\infty$}
\end{tikzpicture}
\end{center}
$\Rightarrow f\left(y\right)\ge 8\Leftrightarrow \dfrac{8y^2}{2y-1}\ge 8\Leftrightarrow {\log_2}\dfrac{8y^2}{2y-1}\ge 3$.\\
Dấu \lq\lq=\rq\rq\ xảy ra khi $y=1$.\\
Vậy $P_{\min}=3$.}
\end{bt}
\begin{bt}%[Nguyễn Thành Nhân]%[Dự án 11 chương trình mới]%[1K6GK-5]
Cho dãy số $(u_n)$ thỏa mãn $\sqrt{3}\log_{15}u_6=\sqrt{4-\log_{15}u_6}$ và $u_{n+1}=u_n+2$ với mọi $n\ge 1$. Giá trị nhỏ nhất của $n$ để $\left(\dfrac{6}{5}\right)^{u_n}>2019$ bằng bao nhiêu?
\loigiai{
\begin{itemize}
\item Ta có $u_{n+1}=u_n+2,\,\forall n\ge 1$ nên $(u_n)$ là cấp số cộng có số hạng đầu là $u_1$ và công sai là $2$.
\item Ta có
\[\sqrt{3}\log_{15}u_6=\sqrt{4-\log_{15}u_6}\Leftrightarrow\heva{&\log_{15}u_6\geq 0\\&3\log^2_{15}u_6+\log_{15}u_6-4=0}
\Leftrightarrow\log_{15}u_6=1.\]
\item Suy ra
\[u_6=15\Leftrightarrow u_1+5\cdot 2=15\Leftrightarrow u_1=5.\]
\item Ta có
\begin{eqnarray*}\left(\dfrac{6}{5}\right)^{u_n}>2019&\Leftrightarrow& u_n>\log_{\tfrac{6}{5}}2019\\
&\Leftrightarrow& 5+(n-1)\cdot 2>\log_{\tfrac{6}{5}}2019\\
&\Leftrightarrow& n>\dfrac{1}{2}(\log_{\tfrac{6}{5}}2019-3)\approx 19{,}37.
\end{eqnarray*}
\item Vậy giá trị nhỏ nhất của $n$ thỏa mãn đề bài là $20$.
\end{itemize}
}
\end{bt}

\centerline{\fcolorbox{red}{yellow!50}{\bf {BÀI TẬP TRẮC NGHIỆM }}}

\Opensolutionfile{ans}[ans/ans-1K6-21-Dang6]
\begin{ex}%[Nguyễn Thành Nhân]%[Dự án 11 chương trình mới]%[1K6YK-5]
Tìm tập nghiệm $S$ của bất phương trình $16-2^{2x+1}\ge0$.
\choice
{$S=\left[\dfrac{3}{2};+\infty\right)$}
{$S=\left(-\infty;\dfrac{3}{2}\right)$}
{\True $S=\left(-\infty;\dfrac{3}{2}\right]$}
{$S=\left(0;\dfrac{3}{2}\right]$}
\loigiai
{
Ta có
\allowdisplaybreaks
\begin{eqnarray*}
& &16-2^{2x+1}\ge0\Leftrightarrow2^{2x+1}\le16 \\
& \Leftrightarrow & 2x+1\le\log_2 16\Leftrightarrow2x+1\le4\\
&\Leftrightarrow&x\le\dfrac{3}{2}.
\end{eqnarray*}
Vậy $S=\left(-\infty;\dfrac{3}{2}\right]$.
}
\end{ex}

\begin{ex}%[Nguyễn Thành Nhân]%[Dự án 11 chương trình mới]%[1K6BK-5]
Tập nghiệm của bất phương trình $2^{x^2+5x+5}>2$ là
\choice
{\True $(-\infty;-4)\cup (-1;+\infty)$}
{$(1;+\infty)$}
{$(-4;-1)$}
{$(-\infty;1)\cup(4;+\infty)$}
\loigiai{
Ta thấy $$2^{x^2+5x+5}>2\Leftrightarrow x^2+5x+5>1\Leftrightarrow x^2+5x+4>0\Leftrightarrow \hoac{&x<-4\\&x>-1.}$$
}
\end{ex}
\begin{ex}%[Nguyễn Thành Nhân]%[Dự án 11 chương trình mới]%[1K6BK-5]
Cho bất phương trình $\left(\dfrac{1}{2}\right)^{4x^2-15x+13}<\left(\dfrac{1}{2}\right)^{4-3x}$. Tập nghiệm của bất phương trình là
\choice
{$\left(\dfrac{3}{2};+\infty\right)$}
{$\mathbb{R}$}
{\True $\mathbb{R}\setminus \left\{\dfrac{3}{2}\right\}$}
{$\varnothing$}
\loigiai{
Ta có\\
$\left(\dfrac{1}{2}\right)^{4x^2-15x+13}<\left(\dfrac{1}{2}\right)^{4-3x}\Leftrightarrow 4x^2-15x+13>4-3x \Leftrightarrow (2x-3)^2>0\Leftrightarrow x\neq \dfrac{3}{2}$.\\
Vậy tập nghiệm của bất phương trình là $S=\mathbb{R}\setminus\left\{\dfrac{3}{2}\right\}$.
}
\end{ex}

\begin{ex}%[Nguyễn Thành Nhân]%[Dự án 11 chương trình mới]%[1K6BK-5]
Tập nghiệm của bất phương trình $5^{x^2-2x}>\dfrac{1}{5}$ là
\choice
{\True $\mathbb{R}\setminus \{ 1\}$}
{$\mathbb{R}$}
{$(1;+\infty)$}
{$\{ 1\}$}
\loigiai{
Ta có $5^{x^2-2x}>\dfrac{1}{5}\Leftrightarrow x^2-2x>-1\Leftrightarrow x^2-2x+1>0\Leftrightarrow x\ne 1$.\\
Vậy tập nghiệm của bất phương trình là $S=\mathbb{R}\setminus \{ 1\}$.
}
\end{ex}



\begin{ex}%[Nguyễn Thành Nhân]%[Dự án 11 chương trình mới]%[1K6BK-5]
Tập nghiệm của bất phương trình $\log_{5} \left( 2x-1 \right) < \log_{5} \left( x+2 \right)$ là
\choice
{\True $S = \left( \dfrac{1}{2};3 \right)$}
{$S = \left( -2;3 \right)$}
{$S = \left( 3;+\infty \right)$}
{$S = \left( -\infty;3 \right)$}
\loigiai{
Điêu kiện xác định $x > \dfrac{1}{2}$. Ta có
$$\log_{5} \left( 2x-1 \right) < \log_{5} \left( x+2 \right) \Leftrightarrow 2x-1<x+2 \Leftrightarrow x<3.$$
Suy ra tập nghiệm $S = \left( \dfrac{1}{2};3 \right)$.
}
\end{ex}

\begin{ex}%[Nguyễn Thành Nhân]%[Dự án 11 chương trình mới]%[1K6BK-5]
Tập nghiệm của bất phương trình $3^{3x+2}>\left( \dfrac{1}{3}\right)^{x-10}$ là
\choice
{\True $ (2;+\infty) $}
{$(-3;+\infty) $}
{$(-6;+\infty) $}
{$ (-\infty;2) $}
\loigiai{Ta có $3^{3x+2}>\left( \dfrac{1}{3}\right)^{x-10}
\Leftrightarrow 3^{3x+2}>3^{10-x}
\Leftrightarrow  3x+2>10-x
\Leftrightarrow  x>2.$
}
\end{ex}

\begin{ex}%[Nguyễn Thành Nhân]%[Dự án 11 chương trình mới]%[1K6BK-5]
Tập nghiệm của bất phương trình $(0{,}5)^{x}\ge 1$ là
\choice
{$(-\infty;2]$}
{$[0;+\infty)$}
{\True $(-\infty;0]$}
{$[2;+\infty)$}
\loigiai{
Ta có $\left(0{,}5\right)^{x}\ge 1\Leftrightarrow(0{,}5)^{x}\ge (0{,}5)^{0}\Leftrightarrow x\le 0$.
}
\end{ex}

\begin{ex}%[Nguyễn Thành Nhân]%[Dự án 11 chương trình mới]%[1K6BK-5]
Tập nghiệm của bất phương trình $\log_{\frac{1}{2}}(2x+3)+\log_{2}(3x+1)>0$ là
\choice
{$-\dfrac{1}{3}<x<2$}
{$-\dfrac{2}{3}<x<2$}
{$x<2$}
{\True $x>2$}
\loigiai{
Điều kiện: $\heva{&2x+3>0\\&3x+1>0} \Leftrightarrow \heva{&x>\dfrac{-3}{2}\\&x>\dfrac{-1}{3}} \Leftrightarrow x>\dfrac{-1}{3}$.\\
Bất phương trình tương đương:
\begin{eqnarray*}
& &-\log_{2}{\left(2x+3\right)}+\log_{2}{\left(3x+1\right)}>0\\
& \Leftrightarrow & \log_{2}{\left(3x+1\right)}>\log_{2}{\left(2x+3\right)}\\
& \Leftrightarrow & 3x+1>2x+3\\
& \Leftrightarrow & x>2.
\end{eqnarray*}
}
\end{ex}

\begin{ex}%[Nguyễn Thành Nhân]%[Dự án 11 chương trình mới]%[1K6BK-5]
Số nghiệm nguyên của bất phương trình $ \log_{\tfrac{1}{2}}(x-3)\ge\log_{\tfrac{1}{2}}4 $ là
\choice
{$ 2 $}
{$ 3 $}
{$ 1 $}
{\True $ 4 $}
\loigiai{
Ta có $ \log_{\tfrac{1}{2}}(x-3)\ge\log_{\tfrac{1}{2}}4\Leftrightarrow\heva{&x-3\le 4\\&x-3>0}\Leftrightarrow\heva{&x\le7\\&x>3}\Leftrightarrow 3<x\le7 $.\\
Suy ra các nghiệm nguyên của bất phương trình đã cho là $ 4; 5; 6; 7 $.\\
Vậy bất phương trình đã cho có $ 4 $ nghiệm nguyên.
}
\end{ex}

\begin{ex}%[[Nguyễn Thành Nhân]%[Dự án 11 chương trình mới]%[1K6BK-5]
Bất phương trình $\log_9 x^2+\log_3 (2-x)<1$ có bao nhiêu nghiệm nguyên?
\choice
{$2$}
{$4$}
{\True $1$}
{$3$}
\loigiai{
Điều kiện $\heva{& x^2>0\\ & 2-x>0}\Leftrightarrow 0\neq x<2$. Bất phương trình đã cho tương đương với
\begin{align*}
\log_3|x|+\log_3(2-x)<\log_3 3\Leftrightarrow |x|(2-x)<3. \tag{$*$}
\end{align*}
\begin{itemize}
\item Nếu $x<0$ thì $(*)\Leftrightarrow -x(2-x)<3\Leftrightarrow x^2-2x-3<0 \Leftrightarrow -1<x<3$. Do đó $-1<x<0$.
\item Nếu $0<x<2$ thì $(*)\Leftrightarrow x(2-x)<3\Leftrightarrow x^2-2x+3>0$ (luôn đúng với mọi $0<x<2$).
\end{itemize}
Vậy bất phương trình có tập nghiệm là $(-1;0)\cup(0;2)$. Do đó chỉ có một nghiệm nguyên là $x=1$.
}
\end{ex}


\begin{ex}%[Nguyễn Thành Nhân]%[Dự án 11 chương trình mới]%[1K6BK-5]
Tìm tập nghiệm $S$ của bất phương trình $\log_3(4x-3)\le\log_9(18x+27)$.
\choice
{$S=\left(\dfrac{3}{4};+\infty\right)$}
{\True $S=\left(\dfrac{3}{4};3\right]$}
{$S=\left[-\dfrac{3}{8};3\right]$}
{$S=[3;+\infty)$}
\loigiai{Điều kiện: $\heva{&4x-3>0\\&18x+27>0}\Leftrightarrow x>\dfrac{3}{4}$.\\
Bất phương trình được viết lại là
\begin{eqnarray*}
&&\log_3(4x-3)\le\dfrac{1}{2}\log_3(18x+27)\Leftrightarrow\log_3(4x-3)^2\le\log_3(18x+27)\\
&\Leftrightarrow&(4x-3)^2\le18x+27\Leftrightarrow16x^2-42x-18\le0\Leftrightarrow-\dfrac{3}{8}\le x\le3.
\end{eqnarray*}
So điều kiện, tập nghiệm bất phương trình là $S=\left(\dfrac{3}{4};3\right]$.}
\end{ex}

\begin{ex}%[Nguyễn Thành Nhân]%[Dự án 11 chương trình mới]%[1K6BK-5]
Tập nghiệm $S$ của bất phương trình $\left(\dfrac{1}{2}\right)^{x^2-4x}<8$ là
\choice
{$S=(-\infty ; 3)$}
{$S=(1 ;+\infty)$}
{\True $S=(-\infty ; 1) \cup (3 ;+\infty)$}
{$S=(1 ; 3)$}
\loigiai{
Bất phương trình $\Leftrightarrow 2^{-x^2+4x}<2^3\Leftrightarrow -x^2+4x<3\Leftrightarrow x^2-4x+3>0\Leftrightarrow \hoac{& x>3\\ &x<1.}$
}
\end{ex}



\begin{ex}%[Nguyễn Thành Nhân]%[Dự án 11 chương trình mới]%[1K6BK-5]
Biết rằng bất phương trình $\left(\dfrac{2}{3}\right)^{x^2-x}\geq \left(\dfrac{9}{4}\right)^{x-1}$ có tập nghiệm là đoạn $[a;b]$. Tính $b-a$.
\choice
{$b-a=2\sqrt{5}$}
{\True $b-a=3$}
{$b-a=\sqrt{5}$}
{$b-a=2$}
\loigiai{
Ta có
\[\left(\dfrac{2}{3}\right)^{x^2-x}\geq \left(\dfrac{9}{4}\right)^{x-1}
\Leftrightarrow\left(\dfrac{2}{3}\right)^{x^2-x}\geq \left(\dfrac{2}{3}\right)^{2-2x}
\Leftrightarrow x^2-x\leq 2-2x\Leftrightarrow x^2+x-2\leq 0\Leftrightarrow -2\leq x\leq 1.\]
Vậy $b-a=3$.
}
\end{ex}




\begin{ex}%[Nguyễn Thành Nhân]%[Dự án 11 chương trình mới]%[1K6BK-5]
Tìm tập nghiệm $S$ của bất phương trình $\ln x^2>\ln (4x-4)$.
\choice
{$S=(2;+\infty)$}
{$S=(1;+\infty)$}
{$S=\mathbb{R}\setminus\left\{ 2\right\}$}
{\True $S=(1;+\infty)\setminus\left\{ 2\right\}$}
\loigiai{
Điều kiện $x>1$.\\
Khi đó $\ln x^2>\ln (4x-4)\Leftrightarrow\heva{&x^2>4x-4\\&x>1}
\Leftrightarrow\heva{&x^2-4x+4>0\\&x>1}\Leftrightarrow\heva{&x>1\\&x\neq 2.}$\\
Vậy tập nghiệm của bất phương trình là $\ln x^2>\ln(4x-4)$ là $S=(1;+\infty)\setminus\left\{ 2\right\}$.}
\end{ex}

\begin{ex}%[Nguyễn Thành Nhân]%[Dự án 11 chương trình mới]%[1K6BK-5]
Tập nghiệm $S$ của bất phương trình ${\left(\dfrac{2}{5}\right)}^{1-3x}\ge \dfrac{25}{4}$ là
\choice
{$S=\left[\dfrac{1}{3};+\infty \right)$}
{$S=\left(-\infty;\dfrac{1}{3}\right)$}
{$S=(-\infty;1]$}
{\True $S=[1;+\infty)$}
\loigiai{
Ta có
\begin{eqnarray*}
&&\left(\dfrac{2}{5}\right)^{1-3x}\ge \dfrac{25}{4}\\
&\Leftrightarrow& \left(\dfrac{2}{5}\right)^{1-3x}\ge \left(\dfrac{2}{5}\right)^{-2}\\
&\Leftrightarrow& 1-3x\le-2\Leftrightarrow 3x\ge 3\Leftrightarrow x\ge 1.
\end{eqnarray*}
Vậy tập nghiệm của bất phương trình là $S=[1;+\infty)$.}
\end{ex}

\begin{ex}%[Nguyễn Thành Nhân]%[Dự án 11 chương trình mới]%[1K6BK-5]%Câu 22
Tập nghiệm của bất phương trình $2\log_3(4x-3)+\log_{\tfrac{1}{3}}(2x+3) \leqslant 2$ là
\choice
{$\left[\dfrac{3}{4};3\right]$}
{$\left(\dfrac{3}{4};+\infty \right)$}
{\True $\left(\dfrac{3}{4};3\right]$}
{$\left[\dfrac{-8}{3};3\right]$}
\loigiai{
Điều kiện: $\heva{&4x-3>0\\ &2x+3>0} \Leftrightarrow \heva{
&x>\dfrac{3}{4}\\&x>-\dfrac{3}{2}} \Leftrightarrow x>\dfrac{3}{4}$.\\
Ta có $\begin{aligned}[t]
2\log_3(4x-3)+\log_{\tfrac{1}{3}}(2x+3) \leqslant 2 & \Leftrightarrow \log_3(4x-3)^2 \leqslant \log_3(2x+3)+\log_39 \\
&\Leftrightarrow \log_3(4x-3)^2 \leqslant \log_3\left[9 \cdot (2x+3)\right] \\
&\Leftrightarrow (4x-3)^2 \leqslant 9(2x+3)\\
&\Leftrightarrow 16x^2-42x-18 \leqslant 0 \\
&\Leftrightarrow -\dfrac{3}{8} \leqslant x \leqslant 3.
\end{aligned}$\\
Kết hợp với điều kiện ta được nghiệm là $\dfrac{3}{4}<x \leqslant 3$.}
\end{ex}

\begin{ex}%[Nguyễn Thành Nhân]%[Dự án 11 chương trình mới]%[1K6BK-5]
Số nghiệm nguyên dương của bất phương trình $0{,}5^{2x-4}>0{,}5^{x+1}$ là
\choice
{$6$}
{$5$}
{Vô số}
{\True $4$}
\loigiai{
Ta có
$$0{,}5^{2x-4}>0{,}5^{x+1} \Leftrightarrow 2x-4<x+1 \Leftrightarrow x<5.$$
Suy ra các nghiệm nguyên dương của bất phương trình là $x \in \{1,2,3,4\}$.\\
Vậy bất phương trình đã cho có $4$ nghiệm nguyên dương.
}
\end{ex}

\begin{ex}%[Nguyễn Thành Nhân]%[Dự án 11 chương trình mới]%[1K6BK-5]
Tập nghiệm của bất phương trình $2^{2x} < 2^{x^2-3}$ là
\choice
{$(-1;3)$}
{\True $(-\infty;-1) \cup  (3;+\infty)$}
{$(1;3)$}
{$(-\infty;1) \cup  (3;+\infty)$}
\loigiai{
Ta có $2^{2x} < 2^{x^2-3} \Leftrightarrow 2x < x^2 -3 \Leftrightarrow x^2-2x-3>0  \Leftrightarrow \hoac{&x>3\\&x<-1.}$
}
\end{ex}

\begin{ex}%[Nguyễn Thành Nhân]%[Dự án 11 chương trình mới]%[2D2K6-2]
\allowdisplaybreaks
Tìm số nghiệm nguyên của bất phương trình $6^x+4\leq 2^{x+1}+2\cdot 3^x$.
\choice
{$2$}
{$3$}
{\True $1$}
{$0$}
\loigiai{
Bất phương trình tương đương:
\begin{eqnarray*}
6^x+4-2\cdot 2^x-2\cdot 3^x\leq 0
&\Leftrightarrow &2^x\left( 3^x -2 \right)-2\left(3^x-2\right)\leq 0\\
&\Leftrightarrow & (3^x -2)(2^x-2)\le 0\\
&\Leftrightarrow &\hoac{&\heva{&3^x-2\geq 0\\&2^x-2\leq 0}\\&\heva{&3^x-2\leq 0\\&2^x-2\geq 0}}\\
&\Leftrightarrow & \hoac{&\heva{&x\geq\log_32\\&x\leq 1}\\&\heva{&x\leq \log_32\\&x\geq 1}\text{(vô lý)}}\\
&\Leftrightarrow & \log_32\leq x\leq 1.
\end{eqnarray*}
Vậy nghiệm nguyên của bất phương trình là $x=1$.
}
\end{ex}


\begin{ex}%[Nguyễn Thành Nhân]%[Dự án 11 chương trình mới]%[1K6BK-5]
Tập nghiệm của bất phương trình $4^{x^2}\le 2^{4x}$ là
\choice
{$(-\infty;0]\cup [2;+\infty)$}
{$(0;2]$}
{\True $[0;2]$}
{$(0;2)$}
\loigiai{
Ta có $4^{x^2}\leq 2^{4x}\Leftrightarrow 2^{2x^2}\le 2^{4x}\Leftrightarrow 2x^2\le 4x\Leftrightarrow 2x^2-4x\le 0\Leftrightarrow x(x-2)\le 0\Leftrightarrow 0\le x\le 2$.\\
Vậy tập nghiệm của bất phương trình đã cho là $[0;2]$.}
\end{ex}

\begin{ex}%[Nguyễn Thành Nhân]%[Dự án 11 chương trình mới]%[1K6BK-5]
Tập nghiệm của bất phương trình $\left(\dfrac{1}{2}\right)^{x^2 + 2x}<\dfrac{1}{8}$ là
\choice
{$(-\infty; -3)$}
{$(1;+\infty)$}
{$(-3;-1)$}
{\True $(-\infty;-3) \cup (1;+\infty)$ }
\loigiai{
$$\left(\dfrac{1}{2}\right)^{x^2 + 2x}<\dfrac{1}{8} \Leftrightarrow x^2+2x >3 \Leftrightarrow \hoac{&x<-3\\&x>1.}$$
Vậy tập nghiệm của bất phương trình là $S=(-\infty;-3) \cup (1;+\infty)$.
}
\end{ex}

\begin{ex}%[Nguyễn Thành Nhân]%[Dự án 11 chương trình mới]%[1K6BK-5]
Tập nghiệm của bất phương trình $\log_2^2(2x)+\log_2\dfrac{x}{4}<9$ chứa tập hợp nào sau đây?
\choice
{$\left(\dfrac{3}{2};6\right)$}
{$(0;3)$}
{$(1;5)$}
{\True $\left(\dfrac{1}{2};2\right)$}
\loigiai{
Điều kiện xác định $x>0$. \\
Phương trình tương đương\\ $\left(1+\log_2x\right)^2+\log_2x-2<9 \Leftrightarrow \left(\log_2x\right)^2+3\log_2x-10<0$
$ \Leftrightarrow -5<\log_2x<2 \Leftrightarrow \dfrac{1}{32}<x<4$. \\
Vậy tập nghiệm của BPT chứa tập $\left(\dfrac{1}{2};2\right)$.}
\end{ex}

\begin{ex} %[Nguyễn Thành Nhân]%[Dự án 11 chương trình mới]%[1K6BK-5]
Tập nghiệm $S$ của bất phương trình $3^{x^2-2x}<27$ là
\choice
{$S=(-\infty;-1)$}
{$S=(3;+\infty)$}
{\True $S=(-1;3)$}
{$S=(-\infty;-1)\cup(3;+\infty)$}
\loigiai{
Ta có $3^{x^2-2x}<27\Leftrightarrow 3^{x^2-2x}<3^3\Leftrightarrow x^2-2x-3<0\Leftrightarrow -1<x<3$.
}
\end{ex}

\begin{ex}%[Nguyễn Thành Nhân]%[Dự án 11 chương trình mới]%[1K6BK-5]
Tập nghiệm của bất phương trình $11^{\sqrt{x+6}}\ge 11^x$ là $S=[a;b]$. Tính $a+b$.
\choice
{$-2$}
{$3$}
{$2$}
{\True $-3$}
\loigiai{
Điều kiện $x\ge -6$.\\
Ta có $11^{\sqrt{x+6}}\ge 11^x\Leftrightarrow \sqrt{x+6}\ge x\quad (1)$.
\begin{itemize}
\item[TH1.] $x\in [-6;0)\Rightarrow (1)$ luôn đúng.
\item[TH2.] $x\ge 0$. Khi đó $(1)\Leftrightarrow x+6\ge x^2\Leftrightarrow x^2-x-6\le 0\Leftrightarrow x\in [-2;3]$.\\
Vì $x\ge 0$ nên $x\in [0;3]$.
\end{itemize}
Tập nghiệm của bất phương trình là $S= [-6;0)\cup [0;3]=[-6;3]$.\\
Suy ra $a=-6$, $b=3$. Vậy $a+b=-3$.
}
\end{ex}


\begin{ex}%[Nguyễn Thành Nhân]%[Dự án 11 chương trình mới]%[1K6BK-5]
Tìm tập nghiệm của bất phương trình $ \left( \dfrac{1}{2}\right)^x \geq 2$.
\choice
{\True $(-\infty; -1]$}
{$[-1;+\infty)$}
{$(-\infty; -1)$}
{$(-1;+\infty)$}
\loigiai{
Ta có $ \left( \dfrac{1}{2}\right)^x \geq 2 \Leftrightarrow 2^{-x} \geq 2 \Leftrightarrow -x \geq 1\Leftrightarrow x\leq -1$.\\
Vậy tập nghiệm của bất phương trình đã cho là $( -\infty; -1]$.
}
\end{ex}



\begin{ex}%[Nguyễn Thành Nhân]%[Dự án 11 chương trình mới]%[1K6BK-5]
Tìm tập nghiệm $S$ của bất phương trình $\left(\dfrac{2}{5}\right)^{1-3x}\geq \dfrac{25}{4}$.
\choice
{$S=\left(-\infty;\dfrac{1}{3}\right)$}
{$S=(-\infty;1]$}
{$S=\left[\dfrac{1}{3};+\infty\right)$}
{\True $S=[1;+\infty)$}
\loigiai{
Ta có $\left(\dfrac{2}{5}\right)^{1-3x}\geq \dfrac{25}{4}\Leftrightarrow \left(\dfrac{2}{5}\right)^{1-3x}\geq \left(\dfrac{5}{2}\right)^{2}\Leftrightarrow 1-3x\leq -2\Leftrightarrow x\geq 1$.\\
Vậy tập nghiệm của bất phương trình là $S=[1;+\infty)$.
}
\end{ex}
\begin{ex}%[Nguyễn Thành Nhân]%[Dự án 11 chương trình mới]%[1K6KK-5]
Cho bất phương trình $2 + \log_3 (x^2 + 1) \geq \log_3 (mx^2 - 2x + m)$. Có bao nhiêu giá trị nguyên của tham số $m$ để bất phương trình trên nghiệm đúng với mọi $x$ thuộc $\mathbb{R}$?
\choice
{$  8 $}
{\True $  7 $}
{$  6 $}
{$ 5  $}
\loigiai{
Điều kiện $mx^2 - 2x + m > 0$.\\
Ta có \begin{eqnarray*}
& & 2 + \log_3 (x^2 + 1) \geq \log_3 (mx^2 - 2x + m)\\
&\Leftrightarrow & 9x^2 + 9 \geq mx^2 - 2x + m\\
&\Leftrightarrow & (9 - m)x^2 + 2x + 9 - m \geq 0.
\end{eqnarray*}
Để bất phương trình có nghiệm đúng với mọi $x \in \mathbb{R}$ thì
$$\heva{&mx^2 - 2x + m > 0, \, \forall x \in \mathbb{R}\\&(9 - m)x^2 + 2x + 9 - m \geq 0, \, \forall x \in \mathbb{R}} \Leftrightarrow \heva{&m > 0\\&1 - m^2 < 0\\& 9 - m > 0\\&1 - (9 - m)^2 \leq 0} \Leftrightarrow \heva{&m > 1\\&m \leq 8} \Leftrightarrow 1 < m \leq 8.$$
Do $m$ nguyên nên $m \in \{2; 3; 4; 5; 6; 7; 8\}$. Vậy có $7$ giá trị $m$ thỏa yêu cầu bài toán.
}
\end{ex}
\begin{ex}%[Nguyễn Thành Nhân]%[Dự án 11 chương trình mới]%[1K6KK-5]
Có bao nhiêu giá trị nguyên của $m$ để bất phương trình $\log_2\left(7x^2+7\right)\geq \log_2\left(mx^2+4x+m\right)$ nghiệm đúng với mọi $x$.
\choice
{$5$}
{\True $4$}
{$0$}
{$3$}
\loigiai{
Ta có
\begin{eqnarray*}
&& \log_2\left(7x^2+7\right)\geq \log_2\left(mx^2+4x+m\right)
\\&\Leftrightarrow & \heva{& mx^2+4x+m>0\\& 7x^2+7\geq mx^2+4x+m}
\\&\Leftrightarrow & \heva{& mx^2+4x+m>0\quad &(1)\\& (m-7)x^2+4x+m-7\leq 0\quad &(2)}
\end{eqnarray*}
Bất phương trình đã cho có nghiệm đúng với mọi $x\in\mathbb{R}$ khi và chỉ khi bất phương trình $(1)$ và $(2)$ có nghiệm đúng với mọi $x\in \mathbb{R}$.
\begin{itemize}
\item Với $\hoac{& m=0\\& m=7}\Rightarrow$ bất phương trình không có nghiệm với mọi $x\in\mathbb{R}$.
\item Với $\heva{& m\neq 0\\& m\neq 7}$ bất phương trình đã cho có nghiệm với mọi $x\in\mathbb{R}$ khi và chỉ khi
\begin{eqnarray*}
&& \heva{& m>0\\& m-7<0\\& \Delta_{(1)}'=4-m^2<0\\&\Delta_{(2)}'=4-(m-7)^2<0}
\Leftrightarrow \heva{& m>2\\& m<7.}
\end{eqnarray*}
\end{itemize}
Suy ra có $4$ giá trị nguyên của $m\in\{3;4;5;6\}$ thỏa mãn yêu cầu.
}
\end{ex}
\begin{ex}%[Nguyễn Thành Nhân]%[Dự án 11 chương trình mới]%[1K6KK-5]
Cho dãy số $(u_n)$ thỏa mãn $\log u_2+\sqrt{2+\log u_2-2\log u_{10}}=2\log u_{10}$ và $u_{n+1}=3u_n$ với mọi $n\ge 1$. Giá trị nhỏ nhất của $n$ để $u_n>3^{2019}$ bằng
\choice
{$2019$}
{$2036$}
{\True $2035$}
{$2018$}
\loigiai{
\begin{itemize}
\item Ta có $u_{n+1}=3u_n,\,\forall n\ge 1$ nên $(u_n)$ là cấp số nhân có số hạng đầu là $u_1$ và công bội là $3$.\hfill{(1)}
\item Đặt $t=2\log u_{10}-\log u_2$, ta được
\[\sqrt{2-t}=t\Leftrightarrow\heva{&t\ge 0\\&2-t=t^2}\Leftrightarrow\heva{&t\ge 0\\&\hoac{&t=1\\&t=-2}}\Leftrightarrow t=1.\tag{2}\]
\item Từ $(1)$ và $(2)$, ta có
\begin{align*}2\log u_{10}-\log u_2=1&\Leftrightarrow 2\log u_1\cdot 3^9-\log u_1\cdot 3=1\\
&\Leftrightarrow 2\log u_1+2\log 3^9-\log u_1-\log 3=1\\
&\Leftrightarrow \log u_1=1-\log 3^{17}.\tag{3}
\end{align*}
\item Ta có
\begin{eqnarray*}u_n>3^{2019}&\Leftrightarrow& u_1\cdot 3^{n-1}>3^{2019}\\
&\Leftrightarrow& \log u_1+\log 3^{n-1}>\log 3^{2019}\\
&\Leftrightarrow&(n-1)\log 3>\log 3^{2019}-1+\log 3^{17}\,\,\text{do\, (3)}\\
&\Leftrightarrow& n-1>2019-\log_3 {10}+17\Leftrightarrow n>2037-\log_3{10}\approx 2034{,}9.
\end{eqnarray*}
\item Vậy giá trị nhỏ nhất của $n$ thỏa mãn đề bài là $2035$.
\end{itemize}
}
\end{ex}
\begin{ex}%[Nguyễn Thành Nhân]%[Dự án 11 chương trình mới]%[1K6KK-5]
Bất phương trình $2(2-x)+\sqrt{4 x^2+3x^3-x^4} \log _4 x > \left(x^2-2x\right) \log _4 x -2 \sqrt{4+3x-x^2}$ có tập nghiệm là $S=(a;b)$. Khi đó $a + b$ bằng
\choice
{$\dfrac{1}{2}$}
{$\dfrac{3}{2}$}
{$\dfrac{5}{2}$}
{\True $\dfrac{7}{2}$}
\loigiai{
Điều kiện: $\heva{& x>0\\& 4+3x-x^2 \ge 0} \Leftrightarrow \heva{& x>0\\& -1 \le x \le 4} \Leftrightarrow 0<x \le 4$.\\
Bất phương trình đã cho tương đương
\[\left( \sqrt{4+3x-x^2}-x+2 \right) \left( x\log _4 x+2 \right) >0 \Leftrightarrow \hoac{& \heva{& \sqrt{4+3x-x^2}-x+2>0 \\& x\log _4 x >-2} & (I)\\& \heva{& \sqrt{4+3x-x^2}-x+2<0 \\& x\log _4 x <-2.} &(II)}\]
Ta có $\sqrt{4+3x-x^2}-x+2=0 \Leftrightarrow \heva{& x-2 \ge 0\\ & 4+3x-x^2=(x-2)^2} \Leftrightarrow \heva{& x\ge 2 \\ & x=0~\text{hoặc}~x=\dfrac{7}{2}} \Leftrightarrow x=\dfrac{7}{2}$.\\
Bảng xét dấu của $\sqrt{4+3x-x^2}-x+2$
\begin{center}
\begin{tikzpicture}[scale=1, font=\footnotesize,line join=round, line cap=round,>=stealth]
\tkzTabInit[nocadre=false,lgt=3.5,espcl=2.5,deltacl=0.8]
{$x$/1,$\sqrt{4+3x-x^2}-x+2$/1}
{$0$,$\dfrac{7}{2}$,$4$}
\tkzTabLine{ ,+,0,-, }
\end{tikzpicture}
\end{center}
Xét hàm số $y=x\log _4 x$ có $y'=\log _4 x + \dfrac{1}{\ln 2}=\log _4 \mathrm{e}x$, $y'=0 \Leftrightarrow x = \dfrac{1}{\mathrm{e}} \in (0;4]$.\\
Bảng biến thiên của $y$
\begin{center}
\begin{tikzpicture}[scale=1, font=\footnotesize,line join=round, line cap=round,>=stealth]
\tkzTabInit[nocadre=false,lgt=1.2,espcl=2.5,deltacl=0.6]
{$x$/1,$y'$/0.6,$y$/2}
{$0$,$\dfrac{1}{\mathrm{e}}$,$4$}
\tkzTabLine{ ,-,0,+, }
\tkzTabVar{+/,-/$-\dfrac{1}{\mathrm{e}}\log _4 \mathrm{e}$,+/$4 \log _4 4$}
\end{tikzpicture}
\end{center}
Suy ra $x\log _2 x \ge -\dfrac{1}{\mathrm{e}}\log _4 \mathrm{e} >-2$. Do đó:
\begin{itemize}
\item Hệ $(I)$ có tập nghiệm là $S=\left( 0;\dfrac{7}{2} \right)$;
\item Hệ $(II)$ vô nghiệm.
\end{itemize}
Vậy bất phương trình đã cho có tập nghiệm là $S=\left( 0;\dfrac{7}{2} \right)$ và $a=0$, $b=\dfrac{7}{2}$ nên được $a + b= \dfrac{7}{2}$.
}
\end{ex}
\begin{ex}%[Nguyễn Thành Nhân]%[Dự án 11 chương trình mới]%[1K6GK-5]
Bất phương trình $\left(x^2-3x\right) \log _3 x + 8 \sqrt{(x+3)(5-x)} < x\left( \sqrt{-x^2+2x+15} \cdot \log _3 x +8\right) -24$ có tập nghiệm là $S=(a;b]$. Khi đó $b-a$ bằng
\choice
{$2+\sqrt{7}$}
{$2-\sqrt{7}$}
{\True $3-\sqrt{7}$}
{$3+\sqrt{7}$}
\loigiai{
Điều kiện: $\heva{& x>0\\& -x^2+2x+15\ge 0} \Leftrightarrow \heva{& x>0\\& -3 \le x \le 5} \Leftrightarrow 0<x \le 5$.\\
Bất phương trình đã cho tương đương
\[\left( \sqrt{-x^2+2x+15}-x+3 \right) \left( x\log _3 x-8 \right) >0 \Leftrightarrow \hoac{& \heva{& \sqrt{-x^2+2x+15}-x+3>0 \\& x\log _3 x >8} & (I)\\& \heva{& \sqrt{-x^2+2x+15}-x+3<0 \\& x\log _3 x <8.} &(II)}\]
Ta có
{\allowdisplaybreaks
\begin{align*}
\sqrt{-x^2+2x+15}-x+3=0 & \Leftrightarrow \heva{& x-3 \ge 0\\ & -x^2+2x+15=(x-3)^2}\\
& \Leftrightarrow \heva{& x\ge 3 \\ & x=2-\sqrt{7}~\text{hoặc}~x=2+\sqrt{7}}\\
& \Leftrightarrow x=2+\sqrt{7}.
\end{align*}}
Bảng xét dấu của $\sqrt{-x^2+2x+15}-x+3$
\begin{center}
\begin{tikzpicture}[scale=1, font=\footnotesize,line join=round, line cap=round,>=stealth]
\tkzTabInit[nocadre=false,lgt=4.5,espcl=2.5,deltacl=0.6]
{$x$/1,$\sqrt{-x^2+2x+15}-x+3$/1}
{$0$,$2+\sqrt{7}$,$5$}
\tkzTabLine{ ,+,0,-, }
\end{tikzpicture}
\end{center}
Xét hàm số $y=x\log _3 x$ có $y'=\log _3 x + \dfrac{1}{\ln 3}=\log _3 \mathrm{e}x$, $y'=0 \Leftrightarrow x = \dfrac{1}{\mathrm{e}} \in (0;5]$.\\
Bảng biến thiên của $y$
\begin{center}
\begin{tikzpicture}[scale=1, font=\footnotesize,line join=round, line cap=round,>=stealth]
\tkzTabInit[nocadre=false,lgt=1.2,espcl=2.5,deltacl=0.8]
{$x$/1,$y'$/0.6,$y$/2}
{$0$,$\dfrac{1}{\mathrm{e}}$,$5$}
\tkzTabLine{ ,-,0,+, }
\tkzTabVar{+/,-/$-\dfrac{1}{\mathrm{e}}\log _3 \mathrm{e}$,+/$5 \log _3 5$}
\end{tikzpicture}
\end{center}
Suy ra $x\log _3 x \le 5 \log _3 5 < 8$. Do đó:
\begin{itemize}
\item Hệ $(I)$ vô nghiệm;
\item Hệ $(II)$ có tập nghiệm là $S=\left( 2+\sqrt{7};5 \right]$.
\end{itemize}
Vậy bất phương trình đã cho có tập nghiệm là $S=\left( 2+\sqrt{7};5 \right]$ và $a=2+\sqrt{7}$, $b=5$ nên được $b-a=3-\sqrt{7}$.
}
\end{ex}
\begin{ex}%[Nguyễn Thành Nhân]%[Dự án 11 chương trình mới]%[1K6GK-5]
Cho hai số thực $x$, $y$ thỏa mãn $xy=4$, $x \ge \dfrac{1}{2}$, $y \ge 1$. Gọi $M$, $m$ lần lượt là giá trị lớn nhất
và giá trị nhỏ nhất của biểu thức $P= \log_2^2 x + \left( \log_2y -1 \right)^2 $. Tính $S=M+m$.
\choice
{$\dfrac{1}{2}$}
{$-\dfrac{1}{2}$}
{$\dfrac{5}{2}$}
{\True $\dfrac{11}{2}$}
\loigiai{
Theo giả thiết ta có $y=\dfrac{4}{x} \ge 1 \Rightarrow x \le 4 \Rightarrow \dfrac{1}{2} \le x \le 4 \Rightarrow \log_2 x \in \left[-1;2\right] $. \\
Đặt $t = \log_2 x$, $\left( -1 \le t \le 2 \right) $, ta có
\begin{eqnarray*}
&&	P= \log_2^2 x + \left( \log_2y -1 \right)^2	= \log_2^2 x + \left( 1-\log_2 x \right)^2\\
&\Rightarrow& P(t) = t^2 + (1-t)^2 = 2t^2 -2t +1  \in \left[ \dfrac{1}{2}; 5 \right]	, \,\, \forall t \in [-1;2]\\
& \Rightarrow& m= \min_{t \in [-1;2]} P = \dfrac{1}{2}, \,\, M= \max_{t \in [-1;2]} P = 5.
\end{eqnarray*}
Do đó $S=M+m=\dfrac{11}{2}$.
}
\end{ex}

\begin{ex}%[Nguyễn Thành Nhân]%[Dự án 11 chương trình mới]%[1K6GK-5]
Cho $x$, $y$ là hai số thực dương thỏa mãn $y \le 4x$. Giá trị nhỏ nhất của biểu thức $P= \ln \dfrac{2x+5y}{y}- \dfrac{2y+5x}{x}$ có dạng $\ln \dfrac{m}{2} -n$. Tính tổng $m+n$.
\choice
{$25$}
{\True $24$}
{$29$}
{$4$}
\loigiai{
Theo giả thiết ta có $\dfrac{x}{y} \ge \dfrac{1}{4} $. Đặt $t= \dfrac{x}{y}$, $t \ge \dfrac{1}{4}$ ta có
\begin{eqnarray*}
&&	P= \ln \dfrac{2x+5y}{y}- \dfrac{2y+5x}{x} = \ln \left( \dfrac{2x}{y} +5 \right) - \dfrac{2y}{x}-5\\
& \Rightarrow& P(t) = \ln \left( 2t +5\right) -\dfrac{2}{t} -5\\
&\Rightarrow & P'(t) = \dfrac{2}{ 2t+5}  +\dfrac{2}{t^2} >0, \,\, \forall t \ge \dfrac{1}{4}.
\end{eqnarray*}
Suy ra hàm số $P(t)$ đồng biến trên $\left[ \dfrac{1}{4} ; + \infty \right) $. Do đó $ \min \limits_{t \in \left[ \frac{1}{4};+\infty \right)} P =P\left( \dfrac{1}{4}\right) = \ln \dfrac{11}{2}-13$. Từ đó $m=11$, $n=13$. Vậy $m+n=24$.
}
\end{ex}
\begin{ex}%[Nguyễn Thành Nhân]%[Dự án 11 chương trình mới]%[1K6GK-5]
Cho hai số thực dương $x$, $y$ thỏa mãn $xy \le 4y-1$. Giá trị nhỏ nhất của biểu thức $P= \dfrac{6(2x+y)}{x} + \ln \dfrac{x+2y}{y}$ được biểu diễn dưới dạng $a+ \ln b$ với $a \in \mathbb{Q}$, $b \in \mathbb{Z}^+$. Tính tích $ab$.
\choice
{\True $81$}
{$45$}
{$108$}
{$115$}
\loigiai{
Từ giả thiết ta có
$$ \heva{&xy \le 4y-1 \\&x>0, y>0 } \Leftrightarrow \dfrac{x}{y} \le \dfrac{4}{y} - \dfrac{1}{y^2} \Leftrightarrow \dfrac{x}{y} \le 4 - \left(2- \dfrac{1}{y} \right)^2 \le 4. $$
Đặt $t= \dfrac{x}{y}$, $0<t \le 4$ ta có
\begin{eqnarray*}
&& P= \dfrac{6(2x+y)}{x} + \ln \dfrac{x+2y}{y} \\
&\Rightarrow& P(t) = 12 +\dfrac{6}{t} + \ln(t+2)\\
&\Rightarrow& P'(t) = -\dfrac{6}{t^2} + \dfrac{1}{t+2} = \dfrac{t^2 -6t-12}{t^2(t+2)} <0,\, \, \forall t \in \left( 0;4\right].
\end{eqnarray*}
Suy ra hàm số $P(t)$ nghịch biến trên $\left( 0;4\right] $, do đó $ \min\limits_{t \in (0;4]}  = P(4) = 12+\dfrac{3}{2} + \ln 6 = \dfrac{27}{2} + \ln6$. \\
Từ đó $a= \dfrac{27}{2}$, $b=6$. Vậy $ab=81$.
}
\end{ex}

\begin{ex}%[Nguyễn Thành Nhân]%[Dự án 11 chương trình mới]%[1K6GK-5]
Tìm tất cả các giá trị thực của tham số $m$ để bất phương trình $\left[\log_2(2x)\right]^2+2\log_2 x+m-3 \ge 0$ có nghiệm đúng với mọi giá trị $x \in (1;32)$.
\choice
{$m\le 2$}
{$m<2$}
{\True $m\ge 2$}
{$m>2$}
\loigiai{
Đặt $t=\log_2 x$, $x\in (1;32) \Rightarrow t\in (0;5)$.\\
\begin{eqnarray*}
& & \left[\log_2(2x)\right]^2+2\log_2 x+m-3 \ge 0, \forall x\in (1;32)\\
& \Leftrightarrow & (1+\log_2 x)^2+2\log_2 x+m-3 \ge 0, \forall x\in (1;32)\\
& \Leftrightarrow & (1+t)^2+2t+m-3\ge 0, \forall t \in (0;5)\\
& \Leftrightarrow & t^2+4t-2+m\ge 0, \forall t\in (0;5)\\
& \Leftrightarrow &  m\ge -t^2-4t+2=g(t), \forall t\in (0;5)\\
& \Leftrightarrow &  m\ge \max\limits_{x \in [0;5]}g(t)=g(0)=2\\
& \Leftrightarrow &  m\ge 2.
\end{eqnarray*}
Vậy $m\ge 2$ thì bất phương trình $\left[\log_2(2x)\right]^2+2\log_2 x+m-3 \ge 0$ có nghiệm đúng với mọi giá trị $x \in (1;32)$.
}
\end{ex}
\Closesolutionfile{ans}
\begin{indapan}{10}
	{ans/ans-1K6-21-Dang6}
\end{indapan}



\begin{dang}{Bài toán thực thế, liên môn}
\end{dang}

\begin{vd}%[1K6KK-6]
	Một người gửi ngân hàng $150$ triệu đồng theo thể thức lãi kép, lãi suất $0{,}58 \%$ một tháng (kể từ tháng thứ $2$, tiền lãi được tính theo phần trăm của tổng tiền lãi tháng trước đó và tiền gốc của tháng trước đó). Sau ít nhất bao nhiêu tháng, người đó có $180$ triệu đồng?
	\loigiai{
		Sau $n$ tháng gửi tiết kiệm, ta có
		$M_n=M_0{(1+r)}^n$, trong đó $M_0=15\cdot10^7$, $r=0{,}0058$.\\
		Do đó $M_n=15\cdot10^7\cdot(1{,}0058)^n$.\\
		Theo giả thiết, ta có $M_n=18\cdot10^7$ (đồng).\\ 
		Do đó, ta có $18\cdot10^7=15\cdot10^7\cdot(1{,}0058)^n\Leftrightarrow(1{,}0058)^n=\dfrac{6}{5}$.
		Sử dụng máy tính cầm tay, ta tính được $n\approx \log (\dfrac{6}{5}): \log (1{,}0058)$ hay $n\approx 31{,}526$.
		Do đó $n=32$.
	}
\end{vd}

\begin{vd}%[1K6KK-6]
	Một người gửi $100$ triệu đồng vào ngân hàng với lãi suất $0,5\%/$ tháng và ông ta rút đều đặn mỗi tháng một triệu đồng kể từ sau ngày gửi một tháng cho đến khi hết tiền (tháng cuối cùng có thể không còn đủ một triệu đồng). Hỏi ông ta rút hết tiền sau bao nhiêu tháng?
	\loigiai{Gọi số tiền lúc đầu người đó gửi là $A$ triệu đồng, lãi suất gửi ngân hàng một tháng là $r$, $S_n$ là số tiền còn lại sau $n$ tháng.\\
		Sau $1$ tháng kể từ ngày gửi tiền, số tiền còn lại của người đó là: $S_1=A(1+r)-1$.\\
		Sau $2$ tháng kể từ ngày gửi tiền, số tiền còn lại của người đó là:
		\begin{align*}
			S_2=\left[A(1+r)-1\right](1+r)-1=A(1+r)^2-(1+r)-1
		\end{align*}
		Sau $n$ tháng kể từ ngày gửi tiền, số tiền còn lại của người đó là:
		\begin{align*}
			S_n=A(1+r)^n-(1+r)^{n-1}-(1+r)^{n-2}-\cdots - (1+r)-1=A(1+r)^n-\dfrac{(1+r)^n-1}{r}.
		\end{align*}
		Giả sử sau $n$ tháng người đó rút hết tiền. Khi đó ta có
		\begin{eqnarray}
			S_n=0&\Leftrightarrow & A(1+r)^n-\dfrac{(1+r)^n-1}{r}=0\Leftrightarrow (1+r)^n(Ar-1)+1=0\\
			&\Leftrightarrow &n=\log_{1+r}\dfrac{1}{1-Ar}\Leftrightarrow n= -\log_{1+r}(1-Ar). 
		\end{eqnarray}
		Với $A=100$ triệu đồng, $r=0{,}005$ ta có $n\approx 138{,}9757216$.}
\end{vd}

\begin{vd}%[1K6KK-6]
	Dân số thế giới cuối năm $2010$, ước tính khoảng $7$ tỉ người. Hỏi với mức tăng trưởng $1,5\%$ mỗi năm thì sau ít nhất bao nhiêu năm nữa dân số thế giới sẽ lên đến $10$ tỉ người?
	\loigiai{Áp dụng công thức: $S_n=A(1+r)^n$.\\
		Suy ra $n=\log_{(1+r)}\left(\dfrac{S_n}{A}\right)$.\\
		Trong đó $A=7$ (tỉ người), $S_n=10$ (tỉ người), $r=1,5\%=\dfrac{1,5}{100}$\\
		Ta được $n=23,95622454$.\\
		Vậy sau ít nhất $24$ năm nữa dân số thế giới sẽ lên đến $10$ tỉ người.}
\end{vd}

\begin{vd}%[1K6KK-6]
	Sự tăng trưởng của loại vi khuẩn tuân theo công thức $ S=A\cdot \mathrm{e}^{rt} $, trong đó $ A $ là số lượng vi khuẩn ban đầu, $ r $ là tỉ lệ tăng trưởng ($ r>0 $), $ t $ là thời gian tăng trưởng (tính theo đơn vị là giờ). Biết số vi khuẩn ban đầu là $ 100 $ con và sau $ 5 $ giờ có $ 300 $ con. Thời gian để số vi khuẩn tăng gấp đôi số vi khuẩn ban đầu là bao nhiêu giờ, bao nhiêu phút?
	\loigiai{Theo giả thiết, ta có hệ phương trình sau
		$$\heva{&3=\mathrm{e}^{5r}\\&2=\mathrm{e}^{rt}}\Leftrightarrow \heva{&5r=\ln 3\\&rt=\ln 2}\Rightarrow t=5\log_{3} 2\approx 3{,}15. $$
		Vậy thời gian để số vi khuẩn tăng gấp đôi số vi khuẩn ban đầu gần đúng nhất đối với kết quả là $ 3 $ giờ $ 9 $ phút.}
\end{vd}

\begin{vd}%[1K6KK-6]
	Dân số thế giới được dự đoán theo công thức $P(t)=ae^{ex}$, trong đó $a$, $b$ là các hằng số, $t$ là năm tính dân số. Theo số liệu thực tế, dân số thế giới năm $1950$ là $2560$ triệu người; dân số thế giới năm $1980$ là $3040$ triệu người. Hãy dự đoán dân số thế giới năm $2020$?
	\loigiai{Từ giả thiết ta có hệ phương trình $\heva{&a \cdot \mathrm{e}^{1950b}=2560&(1)\\&a \cdot \mathrm{e}^{1980b}=3040&(2).}$\\
		Chia $(2)$ cho $(1)$ ta được $\mathrm{e}^{30b}=\dfrac{19}{16} \Leftrightarrow 30b=\ln \dfrac{19}{16} \Leftrightarrow b=\dfrac{1}{30}\ln \dfrac{19}{16}$.\\
		Thay vào $(1)$ ta được $a=\dfrac{2560}{\left(\dfrac{19}{16}\right)^{65}}$.\\
		Vậy $P(2020)=\dfrac{2560}{\left(\dfrac{19}{16}\right)^{65}} \cdot \mathrm{e}^{2020 \cdot \dfrac{1}{30}\ln \dfrac{19}{16}} \approx 3823$ (triệu).}
\end{vd}

\begin{vd}%[1K6KK-6]
	Trong vật lí, sự phân rã của các chất phóng xạ được biểu diễn bởi công thức: $m\left( t \right) = {m_0}{\left( {\dfrac{1}{2}} \right)^{\frac{t}{T}}}$ trong đó ${m_0}$ là khối lượng ban đầu của chất phóng xạ (tại thời điểm $t = 0$); $T$ là chu kì bán rã (tức là khoảng thời gian để một nửa khối lượng chất phóng xạ bị biến thành chất khác). Chu kì bán rã của Cabon $^{14}C$ là khoảng $5730$ năm. Người ta tìm được trong một mẫu đồ cổ một lượng Cabon và xác định được nó đã mất khoảng $25\%$ lượng Cabon ban đầu của nó. Hỏi mẫu đồ cổ đó có tuổi là bao nhiêu?
	
	\loigiai{Giả sử khối lượng ban đầu của mẫu đồ cổ chứa Cabon là ${m_0}$, tại thời điểm $t$ tính từ thời điểm ban đầu ta~có 
		$$m(t) = m_0\mathrm{e}^{- \frac{\ln 2}{5370}t} \Leftrightarrow \dfrac{3m_0}{4} = m_0\mathrm{e}^{ - \frac{\ln 2}{5370}} \Leftrightarrow t = \dfrac{5370\ln\left(\dfrac34\right)}{-\ln2} \approx 2378\text{ (năm).}$$
	}
\end{vd}

\begin{vd}%[1K6KK-6]
	Áp suất không khí $P$ (đo bằng milimet thủy ngân, kí hiệu là mmHg) suy giảm theo hàm số mũ so với độ cao $x$ (đo bằng mét) tính từ mực nước biển, tức là $P$ giảm theo công thức $P=P_0\cdot\mathrm{e}^{xi}$. Trong đó $P_0=760$ mmHg là áp suất không khí ở mực nước biển $(x=0)$, $i$ là hệ số suy giảm. Biết rằng ở độ cao $1000$ m thì áp suất không khí là $672{,}71$ mmHg. Hỏi áp suất không khí ở độ cao $3000$ m là bao nhiêu?
	\loigiai{Ta có
		$$P_{1000}=760\cdot\mathrm{e}^{1000i}=672{,}71\Rightarrow i=\dfrac{\ln\left(\dfrac{672{,}71}{760}\right)}{1000}.$$
		Suy ra $P_{3000}=760\cdot\mathrm{e}^{3000 i}\approx 527{,}06$ mmHg.
	}
\end{vd}

\begin{vd}%[1K6KK-6]
	Một người thả một lá bèo vào một chậu nước. Sau $12$ giờ, bèo sinh sôi phủ kín mặt nước trong chậu. Biết rằng sau mỗi giờ lượng bèo tăng gấp $10$ lần lượng bèo trước đó và tốc độ tăng không đổi. Hỏi sau mấy giờ thì bèo phủ kín $\dfrac{1}{5}$ mặt nước trong chậu (kết quả làm tròn đến $1$ chữ số phần thập phân).
	
	\loigiai{Gọi $S$ là diện tích lá bèo thả ban đầu.\\
		Vì sau mỗi giờ, lượng bèo tăng gấp $10$ lần lượng bèo trước đó nên sau $12$ giờ, tổng diện tích các lá bèo trong chậu là $10^{12}S$.\\
		Sau $12$ giờ, bèo phủ kín mặt nước trong chậu nên diện tích mặt nước trong chậu là $10^{12}S$.\\
		Giả sử sau $x$ giờ thì bèo phủ kín $\dfrac{1}{5}$ mặt nước trong chậu.\\
		Ta có $10^xS=\dfrac{1}{5} \cdot 10^{12}S\Leftrightarrow 10^{12-x}=5\Leftrightarrow x=12-\log 5 \approx 11{,}3$.\\
		Vậy sau $11{,}3$ giờ thì bèo phủ kín $\dfrac{1}{5}$ mặt nước trong chậu.
	}
\end{vd}
\subsubsection{Bài tập vận dụng}
\begin{bt}%[1K6KK-6]
	Ông A gửi vào ngân hàng số tiền $100$ triệu đồng theo kỳ hạn $1$ tháng với lãi suất kép $0{,}5$\%/tháng. Hỏi sau tối thiểu bao nhiêu tháng gửi tiền, ông A thu được số tiền lãi là $10$ triệu đồng?
	\loigiai{Gọi $M$ là số tiền ông A đã gửi, $r$ là lãi suất mỗi tháng.\\
		Số tiền ông A nhận được sau $n$ tháng là $M(1+r)^n$.\\
		Số tiền lãi ông thu được là $L=M(1+r)^n-M$.\\
		Yêu cầu bài toán là $L=10\cdot10^6\Leftrightarrow M((1+r)^n-1)=10\cdot10^6\Leftrightarrow(1+r)^n=1+\dfrac{1}{10}$\\
		$\Leftrightarrow(1+r)^n=\dfrac{11}{10}\Leftrightarrow n=\log_{1+r}\dfrac{11}{10}\approx19{,}11$ tháng.\\
		Vậy sau tối thiểu $20$ tháng thì ông A thu được tiền lãi là $10$ triệu đồng.}
\end{bt}
\begin{bt}%[1K6KK-6]
	Đầu mỗi tháng anh A gửi vào ngân hàng $3$ triệu đồng với lãi suất kép là $0,6\%$ mỗi tháng. Hỏi sau ít nhất bao nhiêu tháng (khi ngân hàng đã tính lãi) thì anh A có được số tiền cả lãi và gốc nhiều hơn $100$ triệu đồng biết lãi suất không đổi trong quá trình gửi.
	\loigiai{Ta có $P=3$ triệu đồng, $r=0{,}006$. \\
		Sau tháng thứ nhất anh A có tổng số tiền là: $P_1=P\left(1+r\right)$.\\
		Sau tháng thứ hai anh A có tổng số tiền là:
		\begin{eqnarray*}
			P_2&=&\left[P+P\left(1+r\right)\right]\left(1+r\right)=P\left[1+\left(1+r\right)\right]\left(1+r\right)\\
			&=&P\left[\left(1+r\right)+\left(1+r\right)^2\right].
		\end{eqnarray*}
		Sau tháng thứ ba anh A có tổng số tiền là:
		\begin{eqnarray*}
			P_3&=&\left[P+P\left[\left(1+r\right)+\left(1+r\right)^2\right]\right]\left(1+r\right)=P\left[1+\left[\left(1+r\right)+\left(1+r\right)^2\right]\right]\left(1+r\right)\\
			&=&P\left[\left(1+r\right)+\left(1+r\right)^2+\left(1+r\right)^3\right].
		\end{eqnarray*}
		Sau tháng thứ tư anh A có tổng số tiền là:
		\begin{eqnarray*}
			P_4&=&\left[P+P\left[\left(1+r\right)+\left(1+r\right)^2+\left(1+r\right)^3\right]\right]\left(1+r\right)\cr
			&=&P\left[1+\left[\left({1+r}\right)+\left(1+r\right)^2+\left(1+r\right)^3\right]\right]\left({1+r}\right)\\
			&=&P\left[\left(1+r\right)+\left(1+r\right)^2+\left(1+r\right)^3+\left(1+r\right)^4\right].
		\end{eqnarray*}
		$\cdots$\\
		Sau tháng thứ $n$ anh A có tổng số tiền là: $$P_n=P\left[\left(1+r\right)+\left(1+r\right)^2+\left(1+r\right)^3+\cdots +\left(1+r\right)^n\right].$$
		Để có được số tiền cả lãi và gốc nhiều hơn $100$ triệu thì
		$$P_n>100\Leftrightarrow P\left[\left(1+r\right)+\left(1+r\right)^2+\left(1+r\right)^3+\cdots +{\left(1+r\right)}^n\right]>100.$$
		Với $P=3$ triệu đồng, $r=0{,}006$ ta có $$\left(1+r\right)+\left(1+r\right)^2+\left(1+r\right)^3+\cdots +\left(1+r\right)^n>\dfrac{100}{3}\left(*\right).$$
		Xem vế trái là cấp số nhân với $u_1=1+r,q=1+r$ khi đó
		$$\left(*\right)\Leftrightarrow \dfrac{\left(1+r\right)\left[{\left(1+r\right)^n-1}\right]}{\left(1+r\right)-1}>\dfrac{100}{3}\Leftrightarrow \left({1{,}006}\right)^n>\dfrac{603}{503}\Leftrightarrow n>\log_{1{,}006}\dfrac{603}{503}\approx 30{,}31.$$
		Vì $n$ nguyên dương nên ta chọn $n=31$.}
\end{bt}

\begin{bt}%[1K6KK-6]
	Cho biết sự tăng dân số được ước tính theo công thức $S=A \cdot \mathrm{e}^{Nr}$ (trong đó $A$ là dân số của năm lấy làm mốc tính, $S$ là dân số sau $N$ năm, $r$ là tỉ lệ tăng dân số hằng năm). Đầu năm 2010 dân số của một tỉnh là $1~038~229$ người tính đến đầu năm 2015 dân số của tỉnh là $1~153~600$ người. Hỏi nếu tỉ lệ tăng dân số hằng năm giữ nguyên thì đầu năm 2020 dân số của tỉnh nằm trong khoảng nào?
	
	\loigiai{Áp dụng công thức $S=A \cdot \mathrm{e}^{Nr}$ từ đầu năm 2010 đến đầu năm 2015 ta có
		$$1153600=1038229 \cdot \mathrm{e}^{5r} \Leftrightarrow r=\dfrac{1}{5}\ln \dfrac{1153600}{1038229}.$$
		Đầu năm 2020 dân số của tỉnh Bắc Ninh là $S=1038229 \cdot \mathrm{e}^{10 \cdot \tfrac{1}{5}\ln \tfrac{1153600}{1038229}} \approx 1281792$ người.}
\end{bt}
\begin{bt}%[1K6KK-6]
	Theo dự báo với mức tiêu thụ dầu không đổi như hiện nay thì trữ lượng dầu của nước A sẽ hết sau 100 năm nữa. Nhưng do nhu cầu thực tế, mức tiêu thụ tăng lên 4\% mỗi năm. Hỏi sau bao nhiêu năm số dầu dữ trữ của nước A sẽ hết.
	\loigiai{Mức tiêu thụ dầu hàng năm của nước A theo dự báo là M thì lượng dầu của nước A là 100M.\\
		Mức tiêu thụ dầu theo thực tế là: Gọi $x_n$ là mức tiêu thụ dầu ở năm thứ $n$ thì ta có:
		\begin{itemize}
			\item Năm thứ nhất: $x_1=M$.
			\item Năm thứ hai: $x_2=M+4\%M=1{,}04M$.
			\item Năm thứ $n$: $x_n=1{,}04^{n-1}M$.
		\end{itemize}
		Tổng tiêu thụ trong $n$ năm là: $$x_1+x_2+x_3+\cdots+x_n=M+1{,}04M+1{,}04^2M+\cdots+1{,}04^{n-1}M=\dfrac{1{,}04^n-1}{1{,}04-1}M.$$
		Theo bài ra ta có phương trình: $\dfrac{1{,}04^n-1}{0{,}04}M=100M$. Giải phương trình bằng lệnh Solve được $n=41$.\\}
\end{bt}

\begin{bt}%[1K6KK-6]
	Dân số thế giới được tính theo công thức $S=A\mathrm{e}^{nr}$, trong đó $A$ là dân số của số năm làm mốc tính, $S$ là dân số sau $n$ năm, $r$ là tỉ lệ tăng dân số hàng năm. Biết rằng dân số Việt Nam vào thời điểm giữa năm 2016 là $90,5$ triệu người và tỉ lệ tăng dân số là $1{,}06\%$ năm. Nếu tỉ lệ tăng dân số hàng năm không đổi thì sau bao nhiêu năm dân số Việt Nam có khoảng $100$ triệu người?
	\loigiai{Ta có $100\cdot10^6=90,5\cdot10^6\mathrm{e}^{1{,}06\%n}\Rightarrow \mathrm{e}^{1{,}06\%n}=\dfrac{200}{181}\Rightarrow 1{,}06\%n=\ln\left(\dfrac{200}{181}\right)\Rightarrow n\approx 9{,}4$.}
\end{bt}

\begin{bt}%[1K6KK-6]
	Chu kì bán rã của nguyên tố phóng xạ ponoli 210 là $138$ ngày (nghĩa là sau $138$ ngày khối lượng của nguyên tố đó chỉ còn một nửa). Thời gian phân rã phóng xạ ponoli 210 để từ $20$ gam còn lại $2{,}22 \cdot 10^{-15}$ gam là bao nhiêu năm?
	
	\loigiai{Gọi $t$ (ngày) là số chu kì bán rã. Khi đó ta có phương trình
		$$20 \cdot \left(\dfrac{1}{2}\right)^t=2{,}22 \cdot 10^{-15} \Rightarrow t \approx 53.$$
		Thời gian phân rã gần bằng $53\cdot138 \colon 365 \approx 20$ (năm).}
\end{bt}
\subsubsection{Bài tập trắc nghiệm}
\Opensolutionfile{ans}[ans/ans-1K6-21-Dang7]
\begin{ex}%[1K6KK-6]
	Một người gửi tiết kiệm $200$ triệu đồng với lãi suất $5\%$ một năm và lãi hàng năm được nhập vào vốn. Sau ít nhất bao nhiêu năm người đó nhận được số tiền nhiều hơn $300$ triệu đồng?
	\choice
	{$11$ năm}
	{$10$ năm}
	{$8$ năm}
	{\True $9$ năm}
	
	\loigiai{Số tiền cả vốn lẫn lãi người đó có được sau $n$ kỳ hạn là $S_{n} = A (1+r)^n = 200 (1+0{,}05)^n=200\cdot 1{,}05^n$.\\
		Ta có $S_{n} >300 \Leftrightarrow 1{,}05^n > 1{,}5 \Leftrightarrow n > \log_{1{,}05} 1{,}5 \approx 8{,}31$ năm.\\
		Vậy người này phải gửi ít nhất $9$ năm để có số tiền hơn $300$ triệu.}
\end{ex}

\begin{ex}%[1K6KK-6]
	Đầu mỗi tháng anh A gửi vào ngân hàng $3$ triệu đồng với lãi suất $0{,}7$\% mỗi tháng. Biết không rút tiền ra khỏi ngân hàng thì cứ sau mỗi tháng số tiền lãi sẽ được nhập vào gốc để tính lãi cho tháng tiếp theo. Hỏi sau ít nhất bao nhiêu tháng (khi ngân hàng đã tính lãi) thì anh $A$ có được số tiền cả gốc lẫn lãi nhiều hơn $100$ triệu đồng? Giả định trong suốt thời gian gửi lãi suất không đổi và anh $A$ không rút tiền ra.
	\choice
	{$29$ tháng}
	{$33$ tháng}
	{$28$ tháng}
	{\True $30$ tháng}
	
	\loigiai{Gọi $P_0$, $P_n$ (triệu đồng) là số tiền ban đầu và số tiên sau $n$ tháng của anh $A$. Số tiền cuối tháng thứ $n$ của anh $A$ là
		$$P_n=P_0\cdot\dfrac{(1+r)\left[(1+r)^n-1\right]}{r}.$$
		Do đó, để được số tiền hơn $100$ triệu đồng thì số tháng cần gửi là$$n=\dfrac{\ln\left(\dfrac{P_n\cdot r}{P_0}+1+r\right)}{\ln(1+r)}-1\approx 29{,}88\,\,\text{(tháng)}.$$}
\end{ex}

\begin{ex}%[1K6KK-6]
	Cho biết năm $2016$, dân số Việt Nam có $94~444~200$ người và tỉ lệ tăng dân số là $1{,}06$\%. Nếu tỉ lệ tăng dân số hàng năm không đổi thì vào năm nào dân số Việt Nam sẽ đạt $100~000~000$ người?
	\choice
	{$2023$}
	{$2021$}
	{$2020$}
	{\True $2022$}
	
	\loigiai{Gọi $T_0=94444200$ là số dân năm $2016$, $r=1{,}06$\% là tỉ lệ tăng dân số, $n$ là số năm và $T_n$ là số dân sau $n$ năm.\\
		Ta có $T_n=T_0(1+r)^n\Leftrightarrow (1+r)^n=\dfrac{T_n}{T_0}\Leftrightarrow n=\log_{1+r} \dfrac{T_n}{T_0}$.\\
		Theo giả thiết suy ra $n=\log_{1+1{,}06\%} \dfrac{100 000 000}{94 444 200}\approx 5{,}42$.\\
		Vậy năm $2022$ dân số Việt Nam sẽ đạt $100000000$ người.}
\end{ex}

\begin{ex}%[1K6KK-6]
	Sự tăng dân số được ước tính theo công thức $P_n=P_0\mathrm{e}^{n r}$, trong đó $P_0$ là dân số của năm lấy làm mốc tính, $P_n$ là dân số sau $n$ năm, $r$ là tỉ lệ tăng dân số hàng năm. Biết rằng năm $2001$, dân số Việt Nam là $78~685~800$ triệu và tỉ lệ tăng dân số năm đó là $1{,}7\%$. Hỏi cứ tăng dân số với tỉ lệ như vậy thì đến năm nào dân số nước ta ở mức $100$ triệu người?
	\choice
	{$2018$}
	{$2017$}
	{$2015$}
	{\True $2016$}
	
	\loigiai{$P_n=P_0\mathrm{e}^{nr}\Leftrightarrow 100000000=78685800\mathrm{e}^{n \cdot 1{,}7\%}\Leftrightarrow n=\dfrac{\ln \dfrac{1000000}{786858}}{1{,}7\%} \approx 14{,}1$.\\
		Sau 15 năm thì dân số nước ta ở mức 100 triệu người.\\
		Do đó năm 2016 dân số nước ta ở mức $100$ triệu người.}
\end{ex}

\begin{ex}%[1K6KK-6]
	Sự tăng trưởng của một loài vi khuẩn theo công thức $S=A \cdot \mathrm{e}^{rt}$ ($A$ là lượng vi khuẩn ban đầu, $r$ là tỉ lệ tăng trưởng $r>0$, $t$ là thời gian tăng trưởng, $S$ là lượng vi khuẩn có được sau thời gian tăng trưởng). Giả sử số lượng vi khuẩn ban đầu là $100$ con và sau 5 giờ có $300$ con. Thời gian để vi khuẩn tăng gấp rưỡi là
	\choice
	{\True $5\log_3\dfrac{3}{2}$}
	{$5\log_32$}
	{$\log_3\dfrac{3}{2}$}
	{$\log_32$}
	
	\loigiai{Theo đề bài ta có phương trình: $300=100\cdot\mathrm{e}^{r \cdot 5} \Rightarrow r=\dfrac{1}{5}\ln 3$.\\
		Suy ra, thời gian để số lượng vi khuẩn tăng gấp rưỡi là
		$$100\mathrm{e}^{\tfrac{1}{5}\ln 3 \cdot t}=150 \Leftrightarrow \left(\mathrm{e}^{\ln 3}\right)^{\tfrac{t}{5}}=\dfrac{3}{2} \Leftrightarrow 3^{\tfrac{t}{5}}=\dfrac{3}{2} \Leftrightarrow t=5\log_3\dfrac{3}{2}.$$}
\end{ex}

\begin{ex}%[1K6KK-6]
	Số lượng của loại vi khuẩn A trong một phòng thí nghiệm được tính theo công thức $S(t)=S(0) \cdot 2^t$, trong đó $S(0)$ là số lượng vi khuẩn A lúc ban đầu, $S(t)$ là số lượng vi khuẩn A có sau $t$ phút. Biết sau $3$ phút thì số lượng vi khuẩn A là $625$ nghìn con. Hỏi sau bao lâu, kể từ lúc ban đầu, số lượng vi khuẩn A là $10$ triệu con?
	\choice
	{$19$ phút}
	{$48$ phút}
	{$12$ phút}
	{\True $7$ phút}
	
	\loigiai{Từ giả thiết ta có $S(3)=S(0) \cdot 2^3$ hay $S(0)=\dfrac{S(3)}{2^3}=\dfrac{625000}{8}=78125$.\\
		Vậy sau $t$ phút, kể từ lúc ban đầu, số lượng vi khuẩn A là 10 triệu con thì $$10000000=78125 \cdot 2^t\Leftrightarrow 2^t=\dfrac{10000000}{78125}=128.$$
		Từ đây suy ra $t=7$ (phút).}
\end{ex}

\begin{ex}%[1K6KK-6]
	Trên một chiếc đài Radio FM có vạch chia để người dùng có thể dò sóng cần tìm. Vạch ngoài cùng bên trái và vạch ngoài cùng bên phải tương ứng với $88$Mhz và $108$Mhz. Hai vạch này cách nhau $10$ cm. Biết vị trí của vạch cách vạch ngoài cùng bên trái $d$ cm thì có tần số bằng $k \cdot a^d$(Mhz) với $k$ và $a$ là hai hằng số. Tìm vị trí tốt nhất của vạch để bắt sóng VOV1 với tần số $102,7$Mhz.
	\choice
	{Cách vạch ngoài cùng bên phải $1,98$ cm}
	{Cách vạch ngoài cùng bên phải $2,46$ cm}
	{\True Cách vạch ngoài cùng bên trái $7,35$ cm}
	{Cách vạch ngoài cùng bên trái $8,23$ cm}
	
	\loigiai{Ta có
		\begin{itemize}
			\item $d=0 \Rightarrow k \cdot a^0=88 \Rightarrow k=88$.\\
			\item $d=10 \Rightarrow k \cdot a^{10}=108\Rightarrow 88 \cdot a^{10}=108\Rightarrow a^{10}=\dfrac{108}{88}\Rightarrow a=\sqrt[10]{\dfrac{108}{88}}$.\\
			\item Gọi $d_1$ là vị trí để vạch có tần số $102,7$Mhz khi đó ta có
			$$88 \left(\sqrt[10]{\dfrac{108}{88}}\right)^{d_1}=102,7\Leftrightarrow \left(\sqrt[10]{\dfrac{108}{88}}\right)^{d_1}=\dfrac{102,7}{88}\Leftrightarrow d_1=\log_{\sqrt[10]{\dfrac{108}{88}}}\dfrac{102,7}{88}=7,54.$$
		\end{itemize}		
		Vậy vị trí tốt nhất của vạch để bắt sóng VOV1 với tần số $102,7$Mhz là $7,35$ cm.}
\end{ex}

\begin{ex}%[1K6KK-6]
	Một điện thoại đang nạp pin, dung lượng pin nạp được tính theo công thức $Q(t)=Q_0\cdot(1-\mathrm{e}^{-t\sqrt2})$ với $t$ là khoảng thời gian tính bằng giờ và $Q_0$ là dung lượng nạp tối đa (pin đầy). Hãy tính thời gian nạp pin của điện thoại tính từ lúc cạn hết pin cho đến khi điện thoại đạt được $90\%$ dung lượng pin tối đa (kết quả được làm tròn đến hàng phần trăm).
	\choice
	{$1{,}65$ giờ}
	{$1{,}61$ giờ}
	{\True $1{,}63$ giờ}
	{$1{,}50$ giờ}
	
	\loigiai{Ta có phương trình: 
		$$90\%Q_0=Q_0(1-\mathrm{e}^{-t\sqrt2})\Leftrightarrow-\mathrm{e}^{-t\sqrt2}=0,9\Leftrightarrow \mathrm{e}^{-t\sqrt2}=0,1\Leftrightarrow -t\sqrt2=\ln0{,}1\approx-2{,}30259\Leftrightarrow  t\approx1{,}62818.$$
	}
\end{ex}

\begin{ex}%[1K6KK-6]
	Áp suất không khí $P$ (đo bằng milimet thủy ngân, kí hiệu là mmHg) suy giảm mũ so với độ cao $x$ (so với mực nước biển) (đo bằng mét) theo công thức $P=P_0 \cdot \mathrm{e}^{xi}$, trong đó $P_0=760$mmHg là áp suất ở mực nước biển ($x=0$), $i$ là hệ số suy giảm. Biết rằng ở độ cao $1000$m thì áp suất của không khí là $672,71$mmHg. Hỏi áp suất không khí ở độ cao $3343$m là bao nhiêu (làm tròn đến hàng phần trăm)?
	\choice
	{\True $505,45$mmHg}
	{$530,23$mmHg}
	{$485,36$mmHg}
	{$495,34$mmHg}
	
	\loigiai{Theo giả thiết ta có $x=1000$m, $P=672,71$mmHg, $P_0=760$mmHg. Khi đó ta được
		$$P=P_0 \cdot \mathrm{e}^{xi}\Leftrightarrow 672,71=760 \cdot \mathrm{e}^{1000i}\Leftrightarrow \mathrm{e}^{1000i}=\dfrac{672,71}{760}\Leftrightarrow i=\dfrac{1}{1000}\ln \dfrac{672,71}{760}.$$
		Với $x=3343$m, $i=\dfrac{1}{1000}\ln \dfrac{672,71}{760}$, $P_0=760$mmHg thay vào công thức $P=P_0 \cdot \mathrm{e}^{xi}$ ta được
		$$P=760 \cdot \mathrm{e}^{3343 \cdot \tfrac{1}{1000}\ln \tfrac{672,71}{760}} \approx 505,4548502.$$}
\end{ex}

\begin{ex}%[1K6KK-6]
	Kết thúc năm 2017, thu nhập bình quân đầu người của Việt Nam đạt $2300$ USD/người/năm. Trong hội nghị mới đây bàn về “Tầm nhìn mới, động lực mới cho tăng trưởng kinh tế”, đại diện chính phủ Việt Nam đặt mục tiêu thu nhập bình quân đầu người của nước ta vào cuối năm 2035 sẽ đạt mức $10~000$ USD/người/năm (theo giá hiện hành). Hỏi để đạt được mục tiêu đó, trung bình mỗi năm thu nhập bình quân đầu người của nước ta tăng bao nhiêu $\%$ (tính gần đúng).
	\choice
	{$8{,}2$}
	{$8{,}7$}
	{$7{,}5$}
	{\True $8{,}5$}
	
	\loigiai{Giả sử để đạt được mục tiêu đề ra, trung bình mỗi năm thu nhập bình quân đầu người của nước ta tăng $x\left(\%\right)$. \\
		Cuối năm 2018, thu nhập bình quân đầu người của Việt Nam là
		$$S_1=2300+\dfrac{x}{100} \cdot 2300=2300 \left(1+\dfrac{x}{100}\right)\text{ (USD)}.$$
		Cuối năm 2019, thu nhập bình quân đầu người của Việt Nam là
		$$S_2=S_1+\dfrac{x}{100} \cdot S_1=2300 \left(1+\dfrac{x}{100}\right)^2\text{ (USD)}.$$
		Cuối năm 2035, thu nhập bình quân đầu người của Việt Nam là
		$$S_{18}=2300 \left(1+\dfrac{x}{100}\right)^{18}\text{ (USD)}.$$
		Ta có 
		\allowdisplaybreaks
		\begin{eqnarray*}
			S_{18}=10000 &\Leftrightarrow& 2300 \left(1+\dfrac{x}{100}\right)^{18}=10000 \\
			&\Leftrightarrow& \left(1+\dfrac{x}{100}\right)^{18}=\dfrac{100}{23}\\
			&\Leftrightarrow& 18\log \left(1+\dfrac{x}{100}\right)=\log \dfrac{100}{23} \\
			&\Leftrightarrow& x \approx 8{,}5.
	\end{eqnarray*}}
\end{ex}
\Closesolutionfile{ans}
\begin{indapan}{10}
	{ans/ans-1K6-21-Dang7}
\end{indapan}


\begin{dang}{Phương pháp đặt ẩn phụ cho PT mũ, logarit}
	
\end{dang}

\begin{vd}%[1K6BK-7]
	Giải phương trình $4\cdot 4^{x}-9 \cdot 2^{x+1}+8=0$.
	\loigiai{Đặt $ t={2^x},t>0$ khi đó phương trình ${{4\cdot 4}^x}-{{9\cdot 2}^{x+1}}+8=0$ tương đương với
		$$4{t^2}-18t+8=0\Leftrightarrow \hoac{& t=4 \\ & t=\dfrac{1}{2} }\Leftrightarrow \hoac{& {x}=2 \\ & {x}=-1.}$$
		Vậy tập nghiệm $S=\{ 2;-1 \}$.}
\end{vd}

\begin{vd}%[1K6KK-7]
	Giải phương trình $(7+4 \sqrt{3})^{x}+(2+\sqrt{3})^{x}=6$.
	\loigiai{Đặt $ t={{\left( 2+\sqrt {3} \right)}^x},t>0$. Khi đó phương trình ${{\left( 7+4\sqrt {3} \right)}^x}+{{\left( 2+\sqrt {3} \right)}^x}=6$ tương đương với
		$${t^2}+t-6=0\Leftrightarrow \hoac{& t=2 \\ & t=-3 \quad (\text{loại}).}$$
		Với $t=2\Leftrightarrow {{\left( 2+\sqrt {3} \right)}^x}=2\Leftrightarrow x={{\log }_{\left( 2+\sqrt {3} \right)}}2$.\\
		Vậy $S=\left\{ {{\log }_{\left( 2+\sqrt {3} \right)}}2 \right\}$.}
\end{vd}

\begin{vd}%[1K6BK-7]
	Giải phương trình $\log_3^2x-2\log_3^{}x-7=0$.
	\loigiai{ Điều kiện: $x>0$. Đặt $\log_3x=t,\left(t\in\mathbb{R}\right)$ \\
		Phương trình $\log_3^2x-2\log_3^{}x-7=0$ trở thành $t^2-2t-7=0\Leftrightarrow\hoac{
			& t=1-2\sqrt{2}\\
			& t=1+2\sqrt{2}.
		}$\begin{itemize}
			\item Với $t=1-2\sqrt{2}\Rightarrow{\log_3}x=1-2\sqrt{2}\Leftrightarrow x=3^{1-2\sqrt{2}}$.
			\item Với $t=1+2\sqrt{2}\Rightarrow{\log_3}x=1+2\sqrt{2}\Leftrightarrow x=3^{1+2\sqrt{2}}$.
	\end{itemize}}
\end{vd}
\begin{vd}%[1K6BK-7]
	Giải phương trình $\log_{\dfrac{1}{3}}^2x - 5 \log_{3}x + 6 = 0$.
	\loigiai{
		Điều kiện xác định $x > 0$. \\
		Phương trình tương đương với $(- \log_{3}x)^2 - 5 \log_{3}x + 6 = 0$. \\
		Đặt $t = \log_{3}x$, phương trình trở thành $t^2 - 5t + 6 = 0 \Leftrightarrow \hoac{& t = 2 \\& t = 3} \Rightarrow \hoac{& x = 9 \\& x = 27.}$}
\end{vd}

\begin{vd}%[1K6KK-7]
	Giải phương trình $2\log_2x+3\log_x2=7$.
	\loigiai{Điều kiện $\heva{
			& x>0\\
			& x\ne 1
		}$. Đặt $\log_2x=t,t\in\mathbb{R}$.\\
		Phương trình $2\log_2x+3\log_x2=7$ trở thành $2t+3\cdot\dfrac{1}{t}=7\Leftrightarrow 2t^2-7t+3=0\Leftrightarrow\hoac{
			& t=3\\
			& t=\dfrac{1}{2}.
		}$\begin{itemize}
			\item Với $t=3\Rightarrow{\log_2}x=3\Leftrightarrow x=8$.
			\item Với $t=\dfrac{1}{2}\Rightarrow{\log_2}x=\dfrac{1}{2}\Leftrightarrow x=\sqrt{2}$.
	\end{itemize}}
\end{vd}
\Opensolutionfile{ans}[ans/ans-1K6-21-Dang8]
\subsection{Bài tập trắc nghiệm}
\begin{ex}%[1K6BK-7]
	Phương trình $3^{2x+1}-4\cdot 3^x+1=0$ có hai nghiệm $x_1$, $x_2$ ($x_1<x_2$). Trong các mệnh đề sau, mệnh đề nào đúng?
	\choice
	{$x_1x_2=-1$}
	{\True $x_1+2x_2=-1$}
	{$2x_1+x_2=0$}
	{$x_1+x_2=-2$}
	\loigiai{
		Ta viết phương trình đã cho thành $3\cdot 3^{2x}-4\cdot 3^x+1=0$.\\
		Đặt $t=3^x\,\,(t>0)$ ta được phương trình $3t^2-4t+1=0\Leftrightarrow \hoac{&t=1\\&t=\dfrac{1}{3}.}$\\
		Với $t=1$ thì $x=0$, với $t=\dfrac{1}{3}$ thì $x=-1$.\\
		Vậy $x_1=-1$ và $x_2=0$. Do đó $x_1+2x_2=-1$.
	}
\end{ex}

\begin{ex}%[1K6BK-7]
	Gọi $S$ là tổng các nghiệm của phương trình $4^{x-1}-3\cdot 2^x+7=0$. Giá trị của $S$ bằng
	\choice
	{$\log_2 7$}
	{$12$}
	{$28$}
	{\True $\log_2 28$}
	\loigiai{
		Đặt $t=2^x\,\,(t>0)$, phương trình trở thành $\dfrac{1}{4}t^2-3t+7=0$.\\
		Ta thấy phương trình ẩn $t$ này có $2$ nghiệm dương phân biệt, đặt là $t_1$, $t_2$. Khi đó, phương trình ẩn $x$ ban đầu có $2$ nghiệm là $x_1$, $x_2$. \\
		Theo định lí Vi-ét ta có $t_1t_2=28\Leftrightarrow 2^{x_1+x_2}=28\Leftrightarrow x_1+x_2=\log_2 28$.
	}
\end{ex}


\begin{ex}%[1K6BK-7]
	Tìm tập nghiệm $S$ của phương trình $7^x+2\cdot 7^{1-x}-9=0$
	\choice
	{$S=\left\{1;\log_2 7\right\}$}
	{\True $S=\left\{1;\log_7 2\right\}$}
	{$S=\left\{\log_2 7\right\}$}
	{$S=\left\{1\right\}$}
	\loigiai{
		Phương trình đã cho tương đương với
		\begin{eqnarray*}
			&&7^x+\dfrac{14}{7^x}-9=0\\
			&\Leftrightarrow& 7^{2x}-9\cdot7^x+14=0\\
			&\Leftrightarrow& \hoac{&7^x=2\\&7^x=7}\Leftrightarrow \hoac{&x=\log_7 2\\&x=1.}
		\end{eqnarray*}
	}
\end{ex}
\begin{ex}%[1K6BK-7]
	Nghiệm của phương trình $25^x+15^x=2\cdot 9^x$ là
	\choice
	{$x=0$, $x=\dfrac{5}{2}$}
	{\True $x=0$}
	{$x=\dfrac{2}{5}$, $x=\dfrac{5}{2}$}
	{$x=1$}
	\loigiai{
		Chia hai vế của phương trình cho $9^x$ ta được
		\begin{eqnarray*}
			\left(\dfrac{5}{3}\right)^{2x}+\left(\dfrac{5}{3}\right)^x-2=0\Leftrightarrow\hoac{&\left(\dfrac{5}{3}\right)^x=1\\&\left(\dfrac{5}{3}\right)^x=-2 &\text{(loại)}}\Leftrightarrow x=0.
		\end{eqnarray*}
	}
\end{ex}



\begin{ex}%[1K6BK-7]
	Nghiệm của phương trình $2^{2x^2+1}-9\cdot 2^{x^2+x}+2^{2x+2}=0$ là
	\choice
	{$x=-1$, $x=\dfrac{1}{2}$}
	{$x=1$, $x=2$}
	{$x=1$, $x=\dfrac{1}{2}$}
	{\True $x=-1$, $x=2$}
	\loigiai{
		Ta có \[2^{2x^2+1}-9\cdot 2^{x^2+x}+2^{2x+2}=0\Leftrightarrow 2\cdot 2^{2x^2}-9\cdot 2^{x^2+x}+4\cdot 2^{2x}=0.\]
		Chia hai vế của phương trình cho $2^{2x}$ ta được
		\begin{eqnarray*}
			&& 2\cdot 2^{2(x^2-x)}-9\cdot 2^{x^2-x}+4=0\\
			&\Leftrightarrow& \hoac{&2^{x^2-x}=4\\&2^{x^2-x}=\dfrac{1}{2}} \\
			&\Leftrightarrow& \hoac{&x^2-x=2\\&x^2-x=-1}\\
			&\Leftrightarrow& \hoac{&x=2\\&x=-1.}
		\end{eqnarray*}
	}
\end{ex}








\begin{ex}%[1K6BK-7]
	Tất cả các nghiệm của phương trình $\log_3^2 x-4\log_3(3x)+7=0$ là
	\choice
	{$x=\dfrac{1}{3}$}
	{\True $x=3$, $x=27$}
	{$x=27$}
	{$x=\dfrac{1}{3}$, $x=\dfrac{1}{27}$}
	\loigiai{
		Điều kiện $x>0$. Phương trình đã cho tương đương với
		\begin{eqnarray*}
			&& \log_3^2 x-4\left(1+\log_3 x\right)+7=0\\
			&\Leftrightarrow& \log_3^2 x-4\log_3 x+3=0\\
			&\Leftrightarrow& \hoac{&\log_3 x=1\\&\log_3 x=3}\\
			&\Leftrightarrow& \hoac{&x=3\\&x=27} \text{(nhận).}
		\end{eqnarray*}
	}
\end{ex}


\begin{ex}%[1K6BK-7]
	Tích tất cả các nghiệm của phương trình $\log_3^2 x-\log_9(9x^2)-1=0$ bằng
	\choice
	{$\dfrac{1}{3}$}
	{\True $3$}
	{$\dfrac{1}{27}$}
	{$27$}
	\loigiai{
		Điều kiện $x>0$. Phương trình đã cho tương đương với
		\begin{eqnarray*}
			&& \log_3^2 x-\left(1+\log_3 |x|\right)-1=0\\
			&\Leftrightarrow& \log_3^2 x-\log_3 x-2=0\\
			&\Leftrightarrow& \hoac{&\log_3 x=-1\\&\log_3 x=2}\\
			&\Leftrightarrow& \hoac{&x=\dfrac{1}{3}\\&x=9} \text{(nhận).}
		\end{eqnarray*}
		Vậy tích các nghiệm bằng $3$.
	}
\end{ex}

\begin{ex}%[1K6BK-7]
	Tích tất cả các nghiệm của phương trình $\log_3 x\cdot\log_3(27x)-4=0$ bằng
	\choice
	{\True $\dfrac{1}{27}$}
	{$\dfrac{244}{81}$}
	{$3$}
	{$\dfrac{1}{3}$}
	\loigiai{
		Điều kiện $x>0$. Phương trình đã cho tương đương với
		\begin{eqnarray*}
			&& \log_3 x\left(3+\log_3 x\right)-4=0\\
			&\Leftrightarrow& \log_3^2 x+3\log_3 x-4=0\\
			&\Leftrightarrow& \hoac{&\log_3 x=1\\&\log_3 x=-4}\\
			&\Leftrightarrow& \hoac{&x=3\\&x=\dfrac{1}{81}} \text{(nhận).}
		\end{eqnarray*}
		Vậy tích các nghiệm bằng $\dfrac{1}{27}$.	
	}
\end{ex}

\begin{ex}%[1K6BK-7]
	Tích tất cả các nghiệm của phương trình $\log_2^2(4x)-3\log_{\sqrt{2}}x-7=0$ bằng
	\choice
	{$\sqrt{2}$}
	{$8\sqrt{2}$}
	{$\dfrac{\sqrt{2}}{2}$}
	{\True $4$}
	\loigiai{
		Điều kiện $x>0$. Phương trình đã cho tương đương với
		\begin{eqnarray*}
			&& \left(2+\log_2 x\right)^2-6\log_2 x-7=0\\
			&\Leftrightarrow& \log_2^2 x-2\log_2 x-3=0\\
			&\Leftrightarrow& \hoac{&\log_2 x=-1\\&\log_2 x=3}\\
			&\Leftrightarrow& \hoac{&x=\dfrac{1}{2}\\&x=8} \text{(nhận).}
		\end{eqnarray*}
		Vậy tích các nghiệm bằng $4$.
	}
\end{ex}

\begin{ex}%[1K6BK-7]
	Tích tất cả các nghiệm của phương trình $\log^2 x^3-20\log \sqrt{x}+1=0$ bằng
	\choice
	{$1$}
	{\True $10\sqrt[9]{10}$}
	{$10$}
	{$\sqrt[10]{10}$}
	\loigiai{
		Điều kiện $x>0$. Phương trình đã cho tương đương với
		\begin{eqnarray*}
			&& 9\log^2 x-10\log x+1=0\\
			&\Leftrightarrow& \hoac{&\log x=1\\&\log x=\dfrac{1}{9}}\\
			&\Leftrightarrow& \hoac{&x=10\\&x=\sqrt[9]{10}} \text{(nhận).}
		\end{eqnarray*}
		Vậy tích các nghiệm bằng $10\sqrt[9]{10}$.
	}
\end{ex}

\Closesolutionfile{ans}
\begin{indapan}{10}
	{ans/ans-1K6-21-Dang8}
\end{indapan}


\begin{dang}{Phương pháp logarit hóa, mũ hóa cho PT mũ, logarit}
	
\end{dang}

\begin{dang}{Phương pháp hàm số, đánh giá cho PT mũ, logarit}
	
\end{dang}
\Opensolutionfile{ans}[ans/ans-1K6-21-Dang10]
\begin{ex}%[1K6BK-9]%[Câu 1]%[Dự án 2 - KNTT- Loc Do]
	Tìm tất cả các giá trị thực của tham số $m$ để phương trình $2^{x^2+1}-2m^2=0$ có nghiệm.
	\choice
	{\True $\hoac{&m\ge 1 \\&m\le -1}$}
	{$m>0$}
	{$\left(-\infty;0\right)$}
	{$\left(0;1\right)$}
	\loigiai{
		Phương trình đưa về $2^{x^2+1}=2m^2$.\\
		Nhận xét $2^{x^2+1}\ge 2$, điều kiện có nghiệm của phương trình $2^{x^2+1}=2m^2$ là \[2m^2\ge 2\Leftrightarrow m^2\ge 1\Leftrightarrow \hoac{&m\ge 1 \\&m\le -1.}\]}
\end{ex}
\begin{ex}%[1K6BK-9]%[Câu 2]%[Dự án 2 - KNTT- Loc Do]
	Có bao nhiêu số nguyên âm $m$ để phương trình $3^x=m+4$ có nghiệm thuộc khoảng $\left(0;+\infty\right)$.
	\choice
	{$3$}
	{\True $2$}
	{$4$}
	{$5$}
	\loigiai{
		Với mọi $x>0$ thì $3^x>3^0=1$. \\
		Do đó để phương trình đã cho có nghiệm thuộc khoảng $\left(0;+\infty\right)$ thì $m+4>1\Leftrightarrow m>-3$. \\
		Vì $m$ nguyên âm nên $m\in \left\{-2;-1\right\}$. \\
		Vậy có 2 giá trị $m$ cần tìm.
	}
\end{ex}
\begin{ex}%[1K6BK-9]%[Câu 3]%[Dự án 2 - KNTT- Loc Do]
	Cho phương trình:$5^{x-2}+x=3$. Nghiệm của phương trình là
	\choice
	{$x=0$}
	{\True $x=2$}
	{$x=4$}
	{$x=1$}
	\loigiai{
		Tập xác định $\mathscr{D}=\mathbb{R}$.\\
		Ta có $5^{x-2}+x=3\Leftrightarrow{5^{x-2}}+x-2=1$.\\
		Xét hàm số $f(t)=5^t+t$ có $f'(t)=5^t\ln 5+1>0$ $\forall t\in \mathbb{R}$.\\
		Vậy hàm số $f(t)=5^t+t$ luôn đồng biến trên $\mathbb{R}$ suy ra phương trình $5^{x-2}+x-2=1$ nếu có nghiệm thì là duy nhất. \\
		Thay $x=2$ vào thỏa mãn, vậy phương trình có nghiệm $x=2$.}
\end{ex}
\begin{ex}%[1K6BK-9]%[Câu 4]%[Dự án 2 - KNTT- Loc Do]
	Giải phương trình $1+3^{\frac{x}{2}}=2^x$.
	\choice
	{$\{2;\log_23\}$}
	{$\{2;\log_32\}$}
	{$\{2;-\log_23\}$}
	{\True $\{2\}$}
	\loigiai{
		Xét phương trình: $1+3^{\tfrac{x}{2}}=2^x\Leftrightarrow{\left(\dfrac{1}{2}\right)^x}+\left(\dfrac{\sqrt{3}}{2}\right)^x=1$.\\
		Ta thấy vế phải là hàm hằng, vế trái là hàm số $f(x)=\left(\dfrac{1}{2}\right)^x+\left(\dfrac{\sqrt{3}}{2}\right)^x$ là hàm nghịch biến trên $\mathbb{R}$. \\
		Nên phương trình đã cho có nghiệm duy nhất.	Ta có $1+3^{\tfrac{2}{2}}=2^2=4$. \\
		Vậy $x=2$ là nghiệm duy nhất của phương trình.}
\end{ex}
\begin{ex}%[1K6BK-9]%[Câu 5]%[Dự án 2 - KNTT- Loc Do]
	Có bao nhiêu số nguyên dương $a$ nhỏ hơn $2021$ sao cho tồn tại số nguyên $x$ thỏa mãn $2^a\left(2^{x+3}-1\right)+a^2=a\left(2^a+2^{x+3}-1\right)?$
	\choice
	{$12$}
	{$15$}
	{\True $10$}
	{$14$}
	\loigiai{
		Ta có 
		\begin{eqnarray*}
			&& 2^a\left(2^{x+3}-1\right)+a^2=a\left(2^a+2^{x+3}-1\right)\\
			&\Leftrightarrow& 2^{x+3} \left(2^a-a\right)-2^a+a^2=a\cdot 2^a-a \\
			&\Leftrightarrow& 2^{x+3}\left(2^a-a\right)=a\left(2^a-a\right)+2^a-a\\
			&\Leftrightarrow& 2^{x+3}=a+1.\qquad(*)
		\end{eqnarray*}
		Để có nghiệm nguyên thì $a+1=2^k\left(k\in \mathbb{Z}\right)\Leftrightarrow a=2^k-1$.\\
		Lại có $1\le a\le 2020\Leftrightarrow 1\le{2^k}-1\le 2020\Leftrightarrow 1\le k\le{\log_2}(2021)\approx 10{,}98$.\\
		Vậy có $10$ giá trị $a$ thỏa mãn.}
\end{ex}
\begin{ex}%[1K6BK-9]%[Câu 6]%[Dự án 2 - KNTT- Loc Do]
	Có bao nhiêu giá trị nguyên của tham số $m$ để phương trình $\left(\dfrac{1}{7}\right)^x-\log_7(m-1)=0$ có nghiệm dương?
	\choice
	{\True $7$}
	{$4$}
	{$5$}
	{$6$}
	\loigiai{
		Với $x>0$ ta có $0<\left(\dfrac{1}{7}\right)^x<1$. \\
		Do đó để phương trình $\left(\dfrac{1}{7}\right)^x-\log_7(m-1)=0$ có nghiệm dương thì \[0<\log_7\left(m-1\right)<1\Leftrightarrow{7^\circ}<m-1<7^1\Leftrightarrow 0<m<8.\]
		Vì $m\in \mathbb{Z}$ và $0<m<8$ nên $m\in \left\{1;2;3;4;5;6;7\right\}$.\\
		Vậy có $7$ giá trị $m$ thoả mãn đề bài.}
\end{ex}
\begin{ex}%[1K6KK-9]%[Câu 7]%[Dự án 2 - KNTT- Loc Do]
	Có bao nhiêu cặp số $(x;y)$ thuộc đoạn $[1;2020]$ thỏa mãn $y$ là số nguyên và $x+\ln x=y+\mathrm{e}^y$?
	\choice
	{$6$}
	{\True $7$}
	{$2021$}
	{$2020$}
	\loigiai{
		Ta có $x+\ln x=y+\mathrm{e}^y\Leftrightarrow x+\ln x=\mathrm{e}^y+\ln{\mathrm{e}^y}$.\\
		Xét hàm số $f(t)=t+\ln t$ trên $\left(0;+\infty\right)$, ta có $f'(t)=1+\dfrac{1}{t}>0,\forall t>0\Rightarrow f(t)$ đồng biến trên $(0;+\infty)$ \break $\Rightarrow $ $f(x)=f\left(\mathrm{e}^y\right)\Rightarrow x=\mathrm{e}^y$.\\
		Do $x\in [1;2020]$ nên $1\le\mathrm{e}^y\le 2020 \Leftrightarrow 0\le y\le \ln 2020$.\\
		Mà $y\in \mathbb{Z}$ và $y\in [1;2020]$ nên $y\in \left\{1;2;3;4;5;6;7\right\}$.\\
		Vậy có $7$cặp số $(x;y)$ thỏa mãn bài toán.}
\end{ex}
\begin{ex}%[1K6KK-9]%[Câu 8]%[Dự án 2 - KNTT- Loc Do]
	Phương trình $3^x=1+2x+2\log_3\left(1+4x\right)$ có 2 nghiệm phân biệt lần lượt là $x_1,x_2\left(x_1<x_2\right)$ khi đó $x_1^3-x_2^3$ bằng
	\choice
	{$-4$}
	{\True $-8$}
	{$-5$}
	{$-7$}
	\loigiai{
		Điều kiện $x>-\dfrac{1}{4}$ và đặt $\log_3(1+4x)=y$.\\
		Suy ra $\heva{&1+4x=3^y \\&1+2x+2y=3^x} \Rightarrow 2x-2y=3^y-3^x\Leftrightarrow{3^x}+2x=3^y+2y\Leftrightarrow f(x)=f(y)$.\\
		Do hàm số $f(t)=3^t+2t$ luôn đơn điệu tăng nên $x=y$. \\
		Từ đó suy ra
		$3^x=4x+1\Leftrightarrow f(x)=3^x-4x-1=0$ và có $f'(x)=3^x\ln 3-4=0$
		\break $\Leftrightarrow x=\log_3\dfrac{4}{\ln 3}>-\dfrac{1}{4}$ nên phương trình $f(x)=0$ có tối đa hai nghiệm và có	$f(2)=f(0)=0$ nên là hai nghiệm cần tìm.
	}
\end{ex}
\begin{ex}%[1K6KK-9]%[Câu 9]%[Dự án 2 - KNTT- Loc Do]
	Phương trình $11^x=2\log_{11}(10x+1)^5+1$ có $2$ nghiệm phân biệt lần lượt là $x_1,x_2\,(x_1<x_2)$ khi đó $x_1^3+x_2^3$ bằng
	\choice
	{$0$}
	{\True $1$}
	{$5$}
	{$2$}
	\loigiai{
		Điều kiện $x>-\dfrac{1}{10}$ thì phương trình $11^x=2\log_{11}(10x+1)^5+1 \Leftrightarrow{11^x}=10\log_{11}(10x+1)+1$.\\
		Đặt $\log_{11}(10x+1)=y\Rightarrow \heva{&10x+1=11^y \\&10y+1=11^x}\Rightarrow 10x+11^x=10y+11^y$
		$\Leftrightarrow f(x)=f(y)$. \\
		Do hàm số $f(t)=10t+11^t$ đơn điệu tăng trên $\mathbb{R}$ nên sẽ có
		$x=y\Rightarrow {11^x}-10x-1=0$ và hàm số $f(x)=11^x-10x-1$ có
		$f'(x)=11^x\ln 11-10=0\Leftrightarrow{11^x}=\dfrac{10}{\ln 11}\Leftrightarrow x=\log_{11}\dfrac{10}{\ln 11}$ nên phương trình $f(x)=0$ không quá hai nghiệm và có $f(1)=f(0)=0$ nên phương trình đã cho có hai nghiệm $x=0 \vee x=1$.
	}
\end{ex}
\begin{ex}%[1K6KK-9]%[Câu 10]%[Dự án 2 - KNTT- Loc Do]
	Phương trình $7^{x-1}=6\log_7(6x-5)+1$ có 2 nghiệm phân biệt lần lượt là $x_1,x_2(x_1<x_2)$ khi đó $x_1-3x_2$ bằng
	\choice
	{\True $-5$}
	{$4$}
	{$5$}
	{$-4$}
	\loigiai{
		Điều kiện $x>\dfrac{5}{6}$. \\
		Đặt $\log_7(6x-5)=y-1\Rightarrow \heva{&6x-5=7^{y-1} \\&6(y-1)+1=7^{x-1}}$
		$\Rightarrow 7^{x-1}+6x=7^{y-1}+6y\Leftrightarrow f(x)=f(y)$. \\
		Xét hàm số $f(t)=7^{t-1}+6t$ có
		$f'(t)=7^{t-1}\ln 7+6>0,\forall t\Rightarrow $ $f(t)$ đơn điệu tăng trên $\mathbb{R}$ và $f(x)=f(y)$
		$\Leftrightarrow x=y$ hay $7^{x-1}=6x-5\Leftrightarrow f(x)=7^{x-1}-6x+5=0$.\\
		Xét hàm số $f(x)=7^{x-1}-6x+5$ trên $\left(\dfrac{5}{6};+\infty\right)$ có $f'(x)=7^{x-1}\ln 7-6=0$ \break
		$\Leftrightarrow x=1+\log_7\dfrac{6}{\ln 7}>\dfrac{5}{6}\Rightarrow $ phương trình không quá hai nghiệm và nhận thấy rằng \break $f(1)=f(2)=0$ nên $x=1\text{;}x=2$ là hai nghiệm cần tìm.}
\end{ex}
\begin{ex}%[1K6KK-9]%[Câu 11]%[Dự án 2 - KNTT- Loc Do]
	Có bao nhiêu cặp số nguyên $(x;y)$ thỏa mãn $0\le x\le 2020$ và \break $\log_3(3x+3)+x=2y+9^y?$
	\choice
	{$2019$}
	{$6$}
	{$2020$}
	{\True $4$}
	\loigiai{
		Điều kiện $x>-1$.\\
		Ta có $\log_3(3x+3)+x=2y+9^y$
		$\Leftrightarrow 3\log_3(3x+3)+3^{\log_3(3x+3)}=3(2y+1)+3^{2y+1}(*)$.\\
		Xét hàm số $f(t)=3t+3^t$ trên $\mathbb{R}$, vì $f'(t)=3+3^t.\ln 3>0,\forall t>0$ nên hàm số $f(t)$ đồng biến trên $\mathbb{R}$.\\
		Từ đó $(*)\Leftrightarrow f(\log_3(3x+3))=f(2y+1)$ $\Leftrightarrow{\log_3}(3x+3)=2y+1$.\\
		Mặt khác $0\le x\le 2020\Leftrightarrow{\log_3}(3x+3)\in \left[1;\log_3(6063)\right]\Rightarrow 2y+1\in \left[1;\log_3(6063)\right]$ \break
		$\heva{&1\le 2y+1\le{\log_3}(6063) \\&y\in Z}\Leftrightarrow 0\le y\le 3$. Vậy có $4$ cặp $(x;y)$ thỏa mãn.}
\end{ex}
\begin{ex}%[1K6GK-9]%[Câu 12]%[Dự án 2 - KNTT- Loc Do]
	Có bao nhiêu giá trị nguyên của tham số $a$ trên đoạn $\left[-10;10\right]$ để phương trình $\mathrm{e}^{x+a}-\mathrm{e}^x=\ln (1+x+a)-\ln (1+x)$ có nghiệm duy nhất.
	\choice
	{$2$}
	{$10$}
	{$1$}
	{\True $20$}
	\loigiai{
		Xét hàm số $f(x)=\mathrm{e}^{x+a}-\mathrm{e}^x-\ln (1+x+a)+\ln (1+x)$ với $\heva{&x>-1 \\&x>-a-1}$. \\
		Ta xét các trường hợp sau
		Nếu $a=0\Rightarrow f(x)=0$ có vô số nghiệm, nên $a=0$ không thỏa.\\
		Ta có $f'(x)=\mathrm{e}^{x+a}-\mathrm{e}^x-\dfrac{1}{1+x+a}+\dfrac{1}{1+x}$, với $\heva{&x>-1 \\&x>-a-1}$.\\
		Nếu $a\ge 1\Rightarrow f'(x)=\left(\mathrm{e}^{x+a}-\mathrm{e}^x\right)+\left(\dfrac{1}{1+x}-\dfrac{1}{1+x+a}\right)>0,\forall x>-1$, đo đó $f(x)$ đồng biến với $\forall x>-1$.\\
		Lại có $\heva{&\underset{x\to{\left(-1\right)^{+}}}{\mathop{\lim}}f(x)=-\infty \\&f\left(1\right)=\mathrm{e}\left(\mathrm{e}^a-1\right)-\ln \dfrac{a+2}{2}>0}$.
		Suy ra phương trình $f(x)=0$ có nghiệm duy nhất.\\
		Nếu $a\le -1\Rightarrow f'(x)=\left(\mathrm{e}^{x+a}-\mathrm{e}^x \right)+\left(\dfrac{1}{1+x}-\dfrac{1}{1+x+a}\right)<0,\forall x>-a-1$, đo đó $f(x)$ nghịch biến với $\forall x>-a-1$, lại có $\heva{&\underset{x\to{\left(-a-1\right)^{+}}}{\mathop{\lim}}f(x)=+\infty \\&f(10)=\mathrm{e}^{10}\left(\mathrm{e}^a-1\right)-\ln \dfrac{a+11}{10}<0}$. \\
		Suy ra phương trình $f(x)=0$ có nghiệm duy nhất.\\
		Vậy ta có $\hoac{&-10\le a\le -1 \\&1\le a\le 10}\Rightarrow $ có $20$ giá trị nguyên thỏa yêu cầu đề bài.}
\end{ex}
\begin{ex}%[1K6GK-9]%[Câu 13]%[Dự án 2 - KNTT- Loc Do]
	Tổng các nghiệm của phương trình $$\dfrac{1}{5^{x+8}}+\dfrac{3^{\left(x+2\right)^2}}{27}+x+1=\dfrac{1}{5\cdot 5^{x^2+4x}}+9\cdot 3^{x+6}+(x+4)(2-x)$$ bằng
	\choice
	{$\sqrt{37}$}
	{$-6$}
	{$3$}
	{\True $-3$}
	\loigiai{
		Ta có
		\begin{eqnarray*}
			&&	\dfrac{1}{5^{x+8}}+\dfrac{3^{(x+2)^2}}{27}+x+1=\dfrac{1}{5\cdot 5^{x^2+4x}}+9\cdot 3^{x+6}+(x+4)(2-x)\\
			& \Leftrightarrow & \dfrac{1}{5^{x+8}}-3^{x+8}-(x+8)=\dfrac{1}{5^{x^2+4x+1}}-3^{x^2+4x+1}-\left(x^2+4x+1\right)\qquad (1).
		\end{eqnarray*}
		Xét hàm số $f(t)=\dfrac{1}{5^t}-3^t-t$,
		có $f'(t)=\dfrac{-\ln 5}{5^t}-3^t\ln 3-1<0,\forall t$ nên hàm số $y=f(t)$ luôn nghịch biến.\\
		Do đó, phương trình $(1)$ có nghiệm $\Leftrightarrow f(x+8)=f(x^2+4x+1)$ có nghiệm \break 
		$\Leftrightarrow x+8=x^2+4x+1$ $\Leftrightarrow x^2+3x-7=0\Leftrightarrow \hoac{&x=\dfrac{-3+\sqrt{37}}{2} \\&x=\dfrac{-3-\sqrt{37}}{2}.}$\\
		Vậy tổng các nghiệm của phương trình trên là $-3$.}
\end{ex}
\begin{ex}%[1K6GK-9]%[Câu 14]%[Dự án 2 - KNTT- Loc Do]
	Số nghiệm của phương trình $\left(\dfrac{1}{16}\right)^{\cos^3x}-\left(\dfrac{1}{8}\right)^{\cos x}=\cos 3x$ trên $\left[0;2021\right]$ là
	\choice
	{$1932$}
	{\True $1930$}
	{$1925$}
	{$1927$}
	\loigiai{
		Ta có $\left(\dfrac{1}{16}\right)^{\cos^3x}-\left(\dfrac{1}{8}\right)^{\cos x}=\cos 3x\Leftrightarrow 3\cos x-\left(\dfrac{1}{2}\right)^{3\cos x}=4\cos^3x-\left(\dfrac{1}{2}\right)^{4\cos^3x}$ $\left(*\right)$.\\
		Xét hàm số $f(t)=t-\left(\dfrac{1}{2}\right)^t$ trên $\mathbb{R}$, ta có:$f'(t)=1+\left(\dfrac{1}{2}\right)^t\ln 2>0$,$\forall t\in \mathbb{R}$.\\
		Suy ra hàm số $f(t)$ đồng biến trên $\mathbb{R}$. \\
		Khi đó
		\begin{eqnarray*}
			(*) &\Leftrightarrow& f(3\cos x)=f(4\cos^3x)\\
			&\Leftrightarrow&  3\cos x=4\cos^3x\\
			&\Leftrightarrow&  4\cos^3x-3\cos x=0 \\
			&\Leftrightarrow& \cos 3x=0\Leftrightarrow 3x=\dfrac{\pi}{2}+k\pi \Leftrightarrow x=\dfrac{\pi}{6}+k\dfrac{\pi}{3},\, (k\in \mathbb{Z}).
		\end{eqnarray*}
		Với $x\in \left[0;2021\right]$, suy ra:$0\le \dfrac{\pi}{6}+k\dfrac{\pi}{3}\le 2021-\dfrac{1}{2}\le k\le \dfrac{6063}{\pi}-\dfrac{1}{2}\approx 1929,41$.\\
		Vậy có $1930$ số $k\in \left\{0;1;2;\cdots;1929\right\}$ hay phương trình có $1930$ nghiệm trên đoạn $[0;2021]$.}
\end{ex}
\begin{ex}%[1K6KK-9]%[Câu 15]%[Dự án 2 - KNTT- Loc Do]
	Số nghiệm $(x;y)$ của phương trình $4^x-2^{x+1}+2(2^x-1)\sin (2^x+y-1)+2=0$ thoả mãn $x,y\in (0;2\pi)$ là
	\choice
	{\True 1}
	{2}
	{3}
	4
	\loigiai{
		Ta có 
		\begin{eqnarray*}
			&& 4^x-2^{x+1}+2(2^x-1)\sin (2^x+y-1)+2=0\\
			&\Leftrightarrow & (2^x-1)^2+2(2^x-1)\sin (2^x+y-1)+\sin^2(2^x+y-1)+\cos^2(2^x+y-1)=0\\
			&\Leftrightarrow & (2^x-1+\sin (2^x+y-1))^2+\cos^2(2^x+y-1)=0\\
			&\Leftrightarrow & \heva{&2^x-1+\sin (2^x+y-1)=0 \\&\cos^2(2^x+y-1)=0}\\
			&\Leftrightarrow & \hoac{&\heva{&2^x=0 \\&\sin (2^x+y-1)=1}(\text{vô nghiệm}) \\&\heva{&2^x=2 \\&\sin (2^x+y-1)=-1}}\\
			&\Leftrightarrow & \heva{&x=1 \\&y=-1-\dfrac{\pi}{2}+k2\pi,k\in \mathbb{Z}.}			
		\end{eqnarray*}
		Đối chiếu điều kiện $x,y\in (0;2\pi)$ ta được $(x;y)=\left(1;\dfrac{3\pi}{2}-1\right)$.}
\end{ex}
\begin{ex}%[1K6KK-9]%[Câu 16]%[Dự án 2 - KNTT- Loc Do]
	Cho hai số thực $x,y$ thoả mãn $0 \leq x \leq 2020$ và $\log _2(2 x+2)+x-3 y=8^y$. Có tất cả bao nhiêu cặp số nguyên $(x; y)$ thỏa mãn các điều kiện đã cho?
	\choice
	{$2018$}
	{\True $4$}
	{$2019$}
	{$1$}
	\loigiai{
		Do $0 \leq x \leq 2020$ nên $x+1>0$.\\
		Ta có $\log_2(2x+2)+x-3y=8^y\Leftrightarrow{\log_2}(x+1)+x+1=\log_22^{3y}+2^{3y}$.\\
		Đặt $y=\log_2t+t(t\ge 1)\Rightarrow y'\text{=}\dfrac{1}{t.\ln 2}+1>0\forall >1\Rightarrow $hàm số $y=\log_2t+t(t\ge 1)$đồng biến.\\
		Từ và ta được $x+1=2^{3y}\Rightarrow x=2^{3y}-1$.\\
		Mà $0\le x\le 2020\Rightarrow 1\le{2^{3y}}\le 2021\Leftrightarrow 1\le{8^y}\le 2021\Leftrightarrow 0\le y\le{\log_8}2021\Leftrightarrow 0\le y\le 3,6604$.\\
		Mà $y$ nguyên nên $y=0, y=1, y=2, y=3$. \\
		Vậy có bốn cặp số nguyên $(x; y)$ thỏa mãn.}
\end{ex}
\begin{ex}%[1K6KK-9]%[Câu 17]%[Dự án 2 - KNTT- Loc Do]
	Có bao nhiêu số nguyên $a$ sao cho ứng với mỗi $a$, tồn tại số thực $b\ge a$ thỏa mãn $4^a=2^b+b$ và đoạn $[a;b]$ chứa không quá $5$ số nguyên?
	\choice
	{$5$}
	{$10$}
	{$6$}
	{\True $11$}
	\loigiai{
		Xét hàm số $f(b)=2^b+b$ có $f'(b)=2^b.\ln 2+1>0,\forall b$.\\
		Suy ra hàm số $f(b)=2^b+b$ đồng biến trên $\mathbb{R}$.\\
		Phương trình $4^a=2^b+b$ có nghiệm $b\ge a$ và đoạn $\left[a;b\right]$ chứa không quá 5 số nguyên khi và chỉ khi $\heva{&2^a+a\le{4^a} \\&4^a<2^{a+5}+a+5}\Leftrightarrow \heva{&4^a-2^a-a\ge 0 \\&4^a-32.2^a-a-5<0.}$\\
		Mà $a\in \mathbb{Z}\Rightarrow a\in \left\{-5;-4;\cdots;5\right\}$.\\
		Vậy có $11$ giá trị của $a$ thỏa mãn yêu cầu đề bài.}
\end{ex}
\begin{ex}%[1K6KK-9]%[Câu 18]%[Dự án 2 - KNTT- Loc Do]
	Có bao nhiêu cặp số nguyên dương $(x;y)$ thoả mãn $0<y<2020$ và \break $4^x+4x-8=16y+\log_4y^4$?
	\choice
	{$9$}
	{$8$}
	{$7$}
	{\True $6$}
	\loigiai{
		Ta có 
		\begin{eqnarray*}
			&& 4^x+4x-8=16y+\log_4y^4\\
			& \Leftrightarrow& 4^x+4x-8=16y+4\log_4y\\
			& \Leftrightarrow& 4^{x-1}+x-2=4y+\log_4y\\
			& \Leftrightarrow& 4^{x-1}+x-1=4^{\log_4(4y)}+\log_44y \qquad (*)
		\end{eqnarray*}
		Xét hàm số $f(t)=4^t+t$. Ta có $f'(t)=1+4^t.\ln 4>0,\forall t$.\\
		Suy ra hàm số $f(t)$ liên tục và đồng biến trên $\mathbb{R}$.\\
		Do đó $(*) \Leftrightarrow f(x-1)=f(\log_4(4y))\Leftrightarrow x-1=\log_4(4y)\Leftrightarrow x-2=\log_4y\Leftrightarrow y=4^{x-2}$.\\
		Vì $y\in (0;2020)$ nên $4^{x-2}<2020\Leftrightarrow x-2<\log_42020\Leftrightarrow x<2+\log_42020$.\\
		Do $x;y\in{\mathbb{Z}^{+}}$ nên $x\in \left\{2;3;4;5;6;7\right\}$.\\
		Ứng với mỗi giá trị nguyên của $x$ cho ta $1$ giá trị nguyên của $y$.\\
		Vậy có $6$ cặp số nguyên $(x;y)$ thoả mãn yêu cầu bài toán.}
\end{ex}
\begin{ex}%[1K6GK-9]%[Câu 19]%[Dự án 2 - KNTT- Loc Do]
	Có bao nhiêu số nguyên $a>2$ để phương trình sau có nghiệm $x>81$. \[\log \left[\left(\log_3x\right)^{\log a}+3\right]=\log_a\left(\log_3x-3\right)\qquad (1).\]
	\choice
	{$2$}
	{$4$}
	{$3$}
	{\True $7$}
	\loigiai{
		Đặt $t=\left(\log_3x\right)^{\log a}+3\Leftrightarrow t-3=\left(\log_3x\right)^{\log a}\Leftrightarrow{\log_a}\left(t-3\right)=\log_a\left(\left(\log_3x\right)^{\log a}\right)$.\\
		Ta có hệ $\heva{&\log t=\log_a\left(\log_3x-3\right) \\&\log_a\left(t-3\right)=\log \left(\log_3x\right)}$. \\
		Suy ra $\log_a\left(t-3\right)+\log t=\log_a\left(\log_3x-3\right)+\log \left(\log_3x\right)$.\\
		Hàm số $f\left(u\right)=\log_a\left(u-3\right)+\log u$ luôn đồng biến trên tập $\left(3;+\infty\right)$.\\
		Khi đó $\log_a\left(t-3\right)+\log t=\log_a\left(\log_3x-3\right)+\log \left(\log_3x\right)\Leftrightarrow t=\log_3x$.\\
		Ta có 
		\begin{eqnarray*}
			t-3=\left(\log_3x\right)^{\log a} &\Leftrightarrow& t-3=t^{\log a}\\
			&\Leftrightarrow& \log_3 (t-3)=\log_3t^{\log a}\\
			&\Leftrightarrow& \log_3(t-3)=\log a\log_3t\\
			&\Leftrightarrow& \log a=\dfrac{\log_3 (t-3)}{\log_3t}\\
			&\Leftrightarrow& \log a=\log_t\left(t-3\right)<1\Leftrightarrow 2<a<10.			
		\end{eqnarray*}
		Vậy có $7$ giá trị nguyên của $a$ cần tìm $a\in \left\{3;4;5;6;7;8;9\right\}$. }
\end{ex}
\begin{ex}%[1K6GK-9]%[Câu 20]%[Dự án 2 - KNTT- Loc Do]
	Có bao nhiêu giá trị nguyên của $m$ để phương trình $\log_3\left(\log_3x-m\right)=x+m$ có nghiệm thuộc khoảng $\left(\dfrac{1}{3};9\right)$?
	\choice
	{$5$}
	{$4$}
	{\True $6$}
	{$7$}
	\loigiai{
		Điều kiện $\heva{&x>0\\ &\log_3x -m>0.}$\\
		Đặt $t=\log_3x-m\Rightarrow {\log_3}x=t+m. \qquad (1)$\\
		Thay vào phương trình ta được $\log_3t=x+m. \qquad (2)$\\
		Từ $(1)$ và $(2)$ suy ra $\log_3x-t=\log_3t-x\Leftrightarrow{\log_3}x+x=\log_3t+t$ $(*)$.\\
		Xét hàm số $f(u)=\log_3u+u,u>0\Rightarrow f'(u)=\dfrac{1}{u\ln 3}+1>0,\forall u>0\Rightarrow $ hàm số $f(u)$ đồng biến và liên tục trên $\left(0;+\infty\right)$.\\
		Phương trình $(*) \Leftrightarrow f(x)=f(t)\Leftrightarrow x=t\Leftrightarrow x=\log_3x-m\Leftrightarrow m=\log_3x-x$.\\
		Xét hàm số $g(x)=\log_3x-x,$ với $x\in \left(\dfrac{1}{3};9\right)$.\\
		Khi đó $g'(x)=\dfrac{1}{x\ln 3}-1=0\Leftrightarrow x=\dfrac{1}{\ln 3}\in \left(\dfrac{1}{3};9\right)$.\\
		Lập bảng biến thiên suy ra $g(9)<m<g\left(\dfrac{1}{\ln 3}\right)\Leftrightarrow -7<m<-\log_3(\ln 3)-\dfrac{1}{\ln 3}$.\\
		Do $m\in \mathbb{Z}\Rightarrow m\in \left\{-6;-5;..;-1\right\}$.}
\end{ex}

\Closesolutionfile{ans}
\begin{indapan}{10}
	{ans/ans-1K6-21-Dang10}
\end{indapan}

\begin{dang}{Giải hệ phương trình mũ và lôgarit}
	Một số phương pháp giải hệ phương trình mũ và lôgarit thường gặp:
	\begin{enumerate}
		\item Phương pháp biến đổi tương đương;
		\item Phương pháp đặt ẩn phụ;
		\item Phương pháp hàm số;
		\item Phương pháp điều kiện cần và đủ.
	\end{enumerate}
\end{dang}
\subsubsection{Ví dụ mẫu}
%%=====Ví dụ 1
\begin{vd}%[1K6BK-0]
	Giải hệ phương trình:
	$ \left\{\begin{aligned}
		&x+y=1\\
		&2^x-2^y=2.\\
	\end{aligned}\right.$
	\loigiai{
		Ta có: \\
		$ \left\{\begin{aligned}
			&x+y=1\\
			&2^x-2^y=2\\
		\end{aligned}\right.\Leftrightarrow \left\{\begin{aligned}
			&2^{x+y}=2\\
			&2^x-2^{y}=2\\
		\end{aligned}\right.\Leftrightarrow \left\{\begin{aligned}
			&2^x\cdot\left(2^x-2\right)=2\\
			&2^y=2^x-2\\
		\end{aligned}\right.\Leftrightarrow \left\{\begin{aligned}
			&2^x=1+\sqrt{3}\\
			&2^y=-1+\sqrt{3}\\
		\end{aligned}\right.\Leftrightarrow \left\{\begin{aligned}
			&x=\log_2\left(1+\sqrt{3}\right)\\
			&y=\log_2\left(-1+\sqrt{3}\right).\\
		\end{aligned}\right.$\\
		Vậy nghiệm của hệ phương trình là $\left\{\begin{aligned}
			&x=\log_2\left(1+\sqrt{3}\right)\\
			&y=\log_2\left(-1+\sqrt{3}\right).\\
		\end{aligned}\right.$ 
	}
\end{vd}
%%=====Ví dụ 2
\begin{vd}%[1K6BK-0]
	Giải hệ phương trình: 
	$ \left\{\begin{aligned}
		&x^2+y^2=17\\
		&\log_2 x+\log_2 y=2.\\
	\end{aligned}\right.$
	\loigiai{
		Điều kiện: $x>0$, $y>0$.\\
		Ta có: \\
		$ \left\{\begin{aligned}
			&x^2+y^2=17\\
			&\log_2 x+\log_2 y=2.\\
		\end{aligned}\right.\Leftrightarrow \left\{\begin{aligned}
			&x^2+y^2=1 7\\
			&\log_2 \left(x\cdot y\right)=2\\
		\end{aligned}\right.\Leftrightarrow \left\{\begin{aligned}
			&x^2+y^2 =17\\
			&x\cdot y=4.\\
		\end{aligned}\right.$\\
		Giải hệ phương trình trên ta tìm được hai nghiệm $ \left\{\begin{aligned}
			&x=1\\
			&y=4\\
		\end{aligned}\right.$ và $ \left\{\begin{aligned}
			&x=4\\
			&y=1.\\
		\end{aligned}\right.$
	}
\end{vd}
%%=====Ví dụ 3
\begin{vd}%[1K6KK-0]
	Giải hệ phương trình: $ \left\{\begin{aligned}
		&3^{2 x+2}+2^{2 y+2}=17\\
		&2\cdot 3^{x+1}+3\cdot 2^y=8.\\
	\end{aligned}\right.$
	\loigiai{
		Đặt $u=3^x$, $v=2^y$ với $u>0$, $v>0$.\\ Khi đó ta được hệ phương trình: $ \left\{\begin{aligned}
			&9 u^2+4 v^2=17\\
			&6 u+3 v=8\\
		\end{aligned}\right.\Leftrightarrow \left\{\begin{aligned}
			&u=\dfrac{1}{3}\\
			&v=2\\
		\end{aligned}\right.\Rightarrow \left\{\begin{aligned}
			&x=-1\\
			&y=1.\\
		\end{aligned}\right.$
	}
\end{vd}
%%=====Ví dụ 4
\begin{vd}%[1K6KK-0]
	Giải hệ phương trình: $ \left\{\begin{aligned}
		&2^{2|x|+1}-3.2^{|x|}=y^2-2\\
		&2 y^2-3 y=2^{2|x|}-2.\\
	\end{aligned}\right.$
	\loigiai{
		Đặt $u=2^{|x|}$ với $u\geq 1$. \\Ta có hệ phương trình: $ \left\{\begin{aligned}
			&2 u^2-3 u=y^2-2\\
			&2 y^2-3 y=u^2-2\\
		\end{aligned}\right.\Leftrightarrow \left\{\begin{aligned}
			&u=2\\
			&y=2.\\
		\end{aligned}\right.$\\
		Ta tìm được các nghiệm của hệ phương trình: $\left(0;1\right)$, $\left(1;2\right)$, $\left(-1;2\right)$.
	}
\end{vd}
%%=====Ví dụ 5
\begin{vd}%[1K6KK-0]
	Giải hệ phương trình: $ \left\{\begin{aligned}
		&2^x+2 x=3+y\\
		&2^y+2 y=3+x.\\
	\end{aligned}\right.$
	\loigiai{
		Ta có: 
		$ \left\{\begin{aligned}
			&2^x+2 x=3+y\\
			&2^y+2 y=3+x.\\
		\end{aligned}\right.\Leftrightarrow \left\{\begin{aligned}
			&2^x+2 x=3+y\\
			&2^x-2^y+2 x-2 y=-x+y\\
		\end{aligned}\right.\Leftrightarrow \left\{\begin{aligned}
			&2^x+2 x=3+y\\
			&2^x+3 x=2^y+3 y\quad \left(*\right).\\
		\end{aligned}\right.$\\
		Xét hàm số $f\left(t\right)=2^t+3t$ đồng biến trên $\mathbb{R}$. \\Khi đó $\left(*\right)\Leftrightarrow f\left(x\right)=f\left(y\right)\Leftrightarrow x=y$.\\
		Hệ phương trình thành $ \left\{\begin{aligned}
			&2^x+2 x=3+y\\
			&x=y\\
		\end{aligned}\right.\Leftrightarrow \left\{\begin{aligned}
			&2^x+2 x=3+x\\
			&x=y\\
		\end{aligned}\right.\Leftrightarrow \left\{\begin{aligned}
			&2^x=-x+3\quad \left(**\right)\\
			&x=y.\\
		\end{aligned}\right.$\\
		Giải phương trình $\left(**\right)$.
		\begin{itemize}
			\item Dễ thấy $x=1$ là nghiệm của phương trình $\left(**\right)$.
			\item Với $x>1$ thì phương trình $\left(**\right)$ có $ \left\{\begin{aligned}
				&VT>2\\
				&VP<2\\
			\end{aligned}\right.$ nên $\left(**\right)$ vô nghiệm khi $x>1$.
			\item Với $x<1$ thì phương trình $\left(**\right)$ có $ \left\{\begin{aligned}
				&VT<2\\
				&VP>2\\
			\end{aligned}\right.$ nên $\left(**\right)$ vô nghiệm khi $x<1$.
		\end{itemize}
		Vậy phương trình $\left(**\right)$ có nghiệm duy nhất nên hệ phương trình có nghiệm $x=y=1$.
	}
\end{vd}
%%=====Ví dụ 6
\begin{vd}%[1K6KK-0]
	Giải hệ phương trình: $ \left\{\begin{aligned}
		&\ln x-\ln y=y-x\\
		&x^2+y^2-6 x-2 y+6=0.\\
	\end{aligned}\right.$
	\loigiai{
		Điều kiện: $x>0$, $y>0$.\\
		Ta có: $\ln x-\ln y=y-x\Leftrightarrow \ln x+x=\ln y+y\quad \left(*\right)$.\\
		Xét hàm số $f\left(t\right)=\ln t+t$ với $t>0$.\\
		Do $f'\left(t\right)=\dfrac{1}{t}+1>0$ với mọi $t>0$ nên hàm số $f\left(t\right)$ đồng biến trên $\left(0;+\infty\right)$.\\
		Khi đó $\left(*\right)\Leftrightarrow f\left(x\right)=f\left(y\right)\Leftrightarrow x=y$.\\
		Ta có hệ phương trình: $ \left\{\begin{aligned}
			&x=y\\
			&x^2+y^2-6 x-2 y+6=0\\
		\end{aligned}\right.\Leftrightarrow \left\{\begin{aligned}
			&x=y\\
			&x^2-4x+3=0\\
		\end{aligned}\right.\Leftrightarrow \left[\begin{aligned}
			&x=y=1\\
			&x=y=3.\\
		\end{aligned}\right.$\\
		Vậy hệ phương trình có hai nghiệm: $\left(1; 1\right)$ và $\left(3;3\right)$.
	}
\end{vd}
%%=====Ví dụ 7
\begin{vd}%[1K6GK-0]
	Tìm $m$ để hệ phương trình sau có nghiệm duy nhất
	$$ \left\{\begin{aligned}
		&2^{|x|}+|x|=y+x^2+m\\
		&x^2+y^2=1.\\
	\end{aligned}\right.$$
	\loigiai{
		Nhận xét rằng nếu $x_0$ là nghiệm của hệ phương trình thì $-x_0$ cũng là nghiệm của hệ phương trình. Để hệ phương trình có nghiệm duy nhất thì $x_0=-x_0\Leftrightarrow x_0=0$.\\
		Với $x=0$ thế vào hệ phương trình ta được: $ \left\{\begin{aligned}
			&1=y+m\\
			&y^2=1\\
		\end{aligned}\right.\Leftrightarrow \left\{\begin{aligned}
			&m=0\\
			&y=1\\
		\end{aligned}\right.$ hoặc $\left\{\begin{aligned}
			&m=2\\
			&y=-1.\\
		\end{aligned}\right.$
		\begin{itemize}
			\item Với $m=0$ hệ phương trình thành $ \left\{\begin{aligned}
				&2^{|x|}+|x|=y+x^2\quad (1)\\
				&x^2+y^2=1\quad (2).\\
			\end{aligned}\right.$\\
			Từ (2) ta có $ \left\{\begin{aligned}
				&0 \leq|x| \leq 1\\
				&-1\leq y\leq 1\\
			\end{aligned}\right.\Leftrightarrow \left\{\begin{aligned}
				&|x| \geq x^2\\
				&2^{|x|} \geq 1 \geq y\\
			\end{aligned}\right.\Rightarrow 2^{|x|}+|x|\geq y+x^2$.\\
			Do đó $(1)\Leftarrow \left\{\begin{aligned}
				&|x|=x^2\\
				&2^{|x|}=y=1\\
			\end{aligned}\right.\Leftrightarrow \left\{\begin{aligned}
				&x=0\\
				&y=1\\
			\end{aligned}\right.$ thỏa mãn (2). Vậy $m=0$ ta nhận.
			\item Với $m=2$ hệ phương trình thành $ \left\{\begin{aligned}
				&2^{|x|}+|x|=y+x^2+2\\
				&x^2+y^2=1\\
			\end{aligned}\right.$.\\ Hệ này có ít nhất hai nghiệm $\left(1;0\right)$ và $\left(-1;0\right)$ nên $m=2$ ta loại.
		\end{itemize}
		Vậy $m=0$ hệ phương trình đã cho có nghiệm duy nhất.
	}
\end{vd}










\subsubsection{Bài tập rèn luyện}
\centerline{\fcolorbox{red}{yellow!50}{\bf {BÀI TẬP TỰ LUẬN }}}
\begin{bt}%[1K6BK-0]
	Giải hệ phương trình $\heva{&x^2+2y=4x-1\\&2\log_3(x-1)-\log_{\sqrt{3}}(y+1)=0} (x,y\in \mathbb{R})$.	
	\loigiai{
		Điều kiện $x>1$, $y>-1$. Hệ đã cho tương đương với 
		{\allowdisplaybreaks
			\begin{align*}
				&\heva{&x^2+2y=4x-1\\&\log_3(x-1)=\log_3(y+1)}\\
				\Leftrightarrow&\heva{&x^2-2x-3=0\\&y=x-2}\\
				\Leftrightarrow&\hoac{&x=-1,\ y=-3\\&x=3,\ y=1.}
		\end{align*}}	
		Đối chiếu với điều kiện ta được nghiệm $(x;y)$ của hệ đã cho là $(3;1)$.
	}
\end{bt}
\begin{bt}%[1K6BK-0]
	Giải hệ phương trình $\heva{& x^2-4x+y+2=0\\& 2\log_2 (x-2)-\log_{\sqrt{2}} y=0}$ $(x,y\in\mathbb{R})$.
	\loigiai{
		Điều kiện $x>2$, $y>1$ \qquad $(1)$.\\
		Từ hệ đã cho, ta có
		\allowdisplaybreaks
		\begin{eqnarray*}
			& & \heva{& x^2-4x+y+2=0\\& x-2=y} \Leftrightarrow \heva{& x^2-3x=0\\& y=x-2}\\
			&\Leftrightarrow & \heva{& x=0\\& y=-2} \text{ hoặc } \heva{& x=3\\& y=1.}
		\end{eqnarray*}
		Đối chiếu với điều kiện $(1)$, ta có nghiệm của hệ là $(x;y)=(3;1)$.
	}
\end{bt}
\begin{bt}%[1K6BK-0]
	Giải hệ phương trình $\heva{&
		\log_2\left(x^2+y^2\right)=1+\log_2(x y)\\&
		3^{x^2-x+y^2}=81
	}\quad(x, y\in\mathbb{R}).$
	\loigiai
	{\begin{itemize}
			\item Với điều kiện $xy>0$ (1), hệ đã cho tương đương $$\heva{&
				x^2+y^2=2x y\\&
				x^2-x y+y^2=4}\Leftrightarrow\heva{&
				x=y\\&
				y^2=4
			}\Leftrightarrow\heva{&
				x=y\\&
				y=\pm 2.
			}$$
			\item Kết hợp (1), hệ có nghiệm $(x ; y)=(2 ; 2)\ \text{và}\ (x ; y)=(-2 ;-2)$.
		\end{itemize}
	}
\end{bt}
%%=====Bài 1
\begin{bt}%[1K6KK-0]
	Giải hệ phương trình: $ \left\{\begin{aligned}
		&3^x-3^y=y-x\\
		&x^2+x y+y^2=12.\\
	\end{aligned}\right.$
	\loigiai{
		Ta có: $ \left\{\begin{aligned}
			&3^x-3^y=y-x\\
			&x^2+x y+y^2=12.\\
		\end{aligned}\right.\Leftrightarrow \left\{\begin{aligned}
			&3^x+x=3^y+y\quad (1)\\
			&x^2+x y+y^2=12\quad (2).\\
		\end{aligned}\right.$\\
		Xét hàm số $f\left(t\right)=3^t+t$ đồng biến trên $\mathbb{R}$.\\
		Khi đó $(2)\Leftrightarrow f\left(x\right)=f\left(y\right)\Leftrightarrow x=y$.\\
		Khi đó hệ phương trình thành $ \left\{\begin{aligned}
			&x=y\\
			&x^2+x y+y^2=12\\
		\end{aligned}\right.\Leftrightarrow \left\{\begin{aligned}
			&x=y\\
			&3x^2=12\\
		\end{aligned}\right.\Leftrightarrow \left\{\begin{aligned}
			&x=y\\
			&x=\pm 2.\\
		\end{aligned}\right.$\\
		Vậy nghiệm của hệ phương trình là $\left(-1; -1\right)$, $\left(1; 1\right)$.
	}
\end{bt}

%%=====Bài 1
\begin{bt}%[1K6KK-0]
	Giải hệ phương trình: $ \left\{\begin{aligned}
		&2^x-2^y=(y-x)\cdot(x y+2)\quad (1)\\
		&x^2+y^2=2\quad (2).
	\end{aligned}\right.$
	\loigiai{
		Thay (2) vào (1) ta được: $2^x-2^y=(y-x)\cdot\left(x^2+y^2+x y\right)\Leftrightarrow 2^x-2^y=y^3-x^3\Leftrightarrow 2^x+x^3=2^y+y^3\quad (*)$.\\
		Xét hàm số $f(t)=2^t+t^3$ có $f'(t)=2^t\ln 2+3t^2>0$ với mọi $t\in\mathbb{R}$. Do đó $f(t)$ đồng biến trên $\mathbb{R}$.\\
		Ta có: $(*)\Leftrightarrow f\left(x\right)=f\left(y\right)\Leftrightarrow x=y$.\\
		Khi đó hệ phương trình thành: $ \left\{\begin{aligned}
			&x=y\\
			&x^2+y^2=2\\
		\end{aligned}\right.\Leftrightarrow \left[\begin{aligned}
			&x=y=-1\\
			&x=y=1.\\
		\end{aligned}\right.$
	}
\end{bt}


\begin{bt}%[1K6GK-0]
	Chứng minh rằng với mọi $a>0$, hệ phương trình sau có nghiệm duy nhất
	$$\heva{&\mathrm{e}^x-\mathrm{e}^y=\ln(1+x)-\ln(1+y)\\&y-x=a.}$$
	\loigiai
	{Điều kiện $x>-1$ và $y>-1$. Khi đó hệ phương trình tương đương
		$$\heva{&\mathrm{e}^{x+a}-\mathrm{e}^x-\ln(1+x+a)+\ln(1+x)&\quad (1)\\&y=x+a.&\quad (2)}$$
		Hệ phương trình có nghiệm duy nhất khi và chỉ khi phương trình $(1)$ có nghiệm duy nhất trong khoảng $(-1;+\infty)$.\\
		Xét hàm số $f(x)=\mathrm{e}^{x+a}-\mathrm{e}^x-\ln(1+x+a)+\ln(1+x)$ trên khoảng $(-1;+\infty)$.
		\begin{itemize}
			\item Hàm số $f(x)$ liên tục trên $(-1;+\infty)$, $\lim\limits_{x \to (-1)^+} f(x)=-\infty$ và $\lim\limits_{x \to +\infty} f(x)=+\infty$.
			\item Với mọi $a>0$, ta có
			\allowdisplaybreaks
			\begin{eqnarray*}
				f'(x) &=& \mathrm{e}^{x+a}-\mathrm{e}^x+\dfrac{1}{1+x}-\dfrac{1}{1+x+a}\\
				&=& \mathrm{e}^{x}\left(\mathrm{e}^a-1\right) +\dfrac{a}{(1+x)(1+x+a)}>0,\, \forall x\in (-1;+\infty).
			\end{eqnarray*}
		\end{itemize}
		Do đó hàm số $f(t)$ đồng biến trên khoảng $(-1;+\infty)$.\\
		Suy ra phương trình $f(x)=0$ luôn có nghiệm duy nhất trong khoảng $(-1;+\infty)$ với mọi $a>0$.\\
		Vậy với mọi $a>0$, hệ phương trình đã cho luôn có nghiệm duy nhất.
	}
\end{bt}
\centerline{\fcolorbox{red}{yellow!50}{\bf {CÂU HỎI TRẮC NGHIỆM (Tầm 10 - 20 câu theo theo tỉ lệ 4:3:2:1)}}}
\Opensolutionfile{ans}[ans/ans-1K6-21-Dang11]
`
\begin{ex}%[1K6KK-0]
	Giải hệ phương trình $\heva{&x+2y=-1\\ & 4^{x+y^2}=16.}$
	\choice
	{$(x ; y)=(-1 ; 1)$, $(x ; y)=(3 ;-7)$}
	{\True $(x ; y)=(1 ;-1)$, $(x ; y)=(-7 ; 3)$}
	{$(x ; y)=(1 ; 1)$, $(x ; y)=(3 ; 7)$}
	{$(x ; y)=(-1 ; 1)$, $(x ; y)=(3 ; 7)$}
	\loigiai{Ta có 
		\begin{eqnarray*}
			& & \heva{&x+2y=-1\\ & 4^{x+y^2}=16}\\
			&\Leftrightarrow & \heva{&x+2y=-1\\ & 4^{x+y^2}=4^2}\\
			&\Leftrightarrow & \heva{&{x+2y=-1}\\ &{x+y^2=2}}\\
			&\Leftrightarrow & \heva{&{x=-2y-1}\\ &{y^2-2y-1=2}}\\
			&\Leftrightarrow & \heva{&{x=-2y-1}\\ &{y^2-2y-3=0}}\\
			&\Leftrightarrow & \heva{&x=-2y-1\\ &{\hoac{&y=-1\\ &y=3}}}\\
			&\Leftrightarrow & \hoac{&y=-1 ; x=1\\ &y=3 ; x=-7.}
		\end{eqnarray*}
	}
\end{ex}

\begin{ex}%[1K6KK-0]
	Giải hệ phương trình $\heva{&\log x-\log y=2\\ &x-10y=900}$.
	\choice
	{$\heva{&x=100\\ &y=10}$}
	{$\heva{&x=1800\\ & y=900}$}
	{\True $\heva{&x=1000\\ & y=10}$}
	{$\heva{&x=10\\ & y=1000}$}
	\loigiai{Điều kiện $x, y>0$. Hệ phương trình tương đương với
		$$\heva{&{\operatorname{l o g}\dfrac xy=2}\\ &{x-1 0y=9 0 0}}\Leftrightarrow\heva{&{\dfrac xy=1 0 0}\\ &{x-1 0y=9 0 0}}\Leftrightarrow\heva{&{x-1 0 0y=0}\\ &{x-1 0y=9 0 0}}\Leftrightarrow\heva{&x=1000\\ &y=10.}$$
	}
\end{ex}



%3
\begin{ex}%[1K6KK-0]
	Gọi $\left(x_0 ; y_0\right)$ là một nghiệm của hệ phương trình $\heva{&x+y=25\\ &\log_2x-\log_2y=2}$. Mệnh đề nào sau đây đúng?
	\choice
	{\True $x_0=4y_0$}
	{$x_0=4+y_0$}
	{$y_0=4x_0$}
	{$y_0=4+x_0$}
	\loigiai{ Điều kiện $\heva{&x>0\\ & y>0}$. Hệ phương trình tương đương với
		$$\heva{&{x+y=2 5}\\ &{\log_2\dfrac xy=2}}\Leftrightarrow\heva{&{x+y=2 5}\\ &{\dfrac xy=4}}\Leftrightarrow\heva{&{x+y=2 5}\\ &{x-4y=0}}\Leftrightarrow\heva{&x=20=x_0\\ &y=5=y_0}\Rightarrow x_0=4y_0.$$
	}
\end{ex}

%4
\begin{ex}%[1K6KK-0]
	Cặp số $(x; y)$ nào sau đây thỏa mãn hệ phương trình $\heva{&\log_4x+\log_4 2y=1+\log_4 9\\ & x+2y=20}$?
	\choice
	{$(x ; y)=(9 ; 2)$}
	{\True $(x ; y)=(18 ; 1)$}
	{$(x ; y)=(1 ; 18)$}
	{$(x ; y)=(16 ; 2)$}
	\loigiai{ Điều kiện $\heva{&x>0\\ & y>0}$. \\
		Ta có 
		\begin{eqnarray*}
			& & \heva{&\log_4x+\log_4 2y=1+\log_4 9\\ & x+2y=20}\\
			&\Leftrightarrow & \heva{&\log_4(2x y)=\log_4 36\\ & x+2y=20}\\
			&\Leftrightarrow & \heva{&{2x y=3 6}\\ &{x+2y=2 0}}\\
			&\Leftrightarrow & \heva{&{x y=1 8}\\ &{x=2 0-2y}}\\
			&\Leftrightarrow & \heva{&{2y^2-2 0y+1 8=0}\\ &{x=2 0-2y}}\\
			&\Leftrightarrow & \heva{& \hoac{& y=1 \\ & y=2} \\ & x=20-2y}\\
			&\Leftrightarrow & \hoac{&y=1 ; x=18\\ &y=9 ; x=2.}
		\end{eqnarray*}		
	}
\end{ex}

%5
\begin{ex}%[1K6KK-0]
	Hệ phương trình $\heva{&2^x\cdot 9^y=162\\ & 3^x\cdot 4^y=48}$ có tất cả bao nhiêu nghiệm $(x ; y)$?
	\choice
	{$0$}
	{\True $1$}
	{$2$}
	{$3$}
	\loigiai{ Nhân vế theo vế trong hệ phương trình, ta được
		$$6^x\cdot 36^y=162.48\Leftrightarrow 6^{x+2y}=6^5\Leftrightarrow x+2y=5.$$
		Thay $x=5-2y$ và phương trình thứ hai của hệ, ta có
		$$ 3^{5-2y}\cdot 4^y=48		 
		\Leftrightarrow\dfrac{3^5}{9^y}\cdot 4^y=2^4\cdot 3\Leftrightarrow\left(\dfrac 23\right)^{2y}=\left(\dfrac 23\right)^4\Leftrightarrow 2y=4\Leftrightarrow y=2\Rightarrow x=1 .$$
		Vậy hệ phương trình có duy nhất nghiệm $(x ; y)=(1 ; 2)$.
	}
\end{ex}


%6 
\begin{ex}%[1K6KK-0]
	Tìm tất cả các cặp số $(x ; y)$ thỏa mãn hệ phương trình $\heva{&6^x-2\cdot 3^y=2\\ & 6^x\cdot 3^y=12}$
	\choice
	{$(x ; y)=\left(1 ;\log_3 4\right)$}
	{$(x ; y)=\left(\log_6 2 ; 1\right)$}
	{\True $(x ; y)=\left(1 ;\log_3 2\right)$}
	{$(x ; y)=\left(1 ;\log_3 2\right)$, $(x ; y)=\left(\log_6 2 ; 1\right)$}
	\loigiai{ Đặt $\heva{&6^x=a>0\\ & 3^y=b>0}$. Hệ phương trình trở thành	
		$$\heva{&a-2b=2\\ &a b=12}\Leftrightarrow\heva{&{a=2b+2}\\ &{(2b+2)b=1 2}}\Leftrightarrow\heva{&a=2b+2\\ &b^2+b-6=0}\Leftrightarrow\heva{&a=2b+2\\ &{\hoac{&b=-3(\text{loại})\\ &b=2(\text{thỏa mãn})}}}\Leftrightarrow\heva{&a=6\\ &b=2.}$$	
		Suy ra $\heva{&6^x=6\\ & 3^y=2}\Leftrightarrow\heva{&x=1\\ & y=\log_3 2}$.
	}
\end{ex}

%7
\begin{ex}%[1K6KK-0]
	Gọi $\left(x_0 ; y_0\right)$ là một nghiệm của hệ phương trình $\heva{&\log_x y=2\\ &\log_{x+1}(y+23)=3}$. Mệnh đề nào sau đây đúng?
	\choice
	{$x_0=y_0$}
	{\True $x_0>y_0$}
	{$x_0<y_0$}
	{$x_0=y_0+2$}
	\loigiai{Điều kiện $\heva{&0<x\neq 1\\ & y>0}$. Hệ phương trình tương đương với
		\begin{eqnarray*}
			& & \heva{&y=x^2\\ & y+23=(x+1)^3}\\
			&\Leftrightarrow & \heva{&{y=x^2}\\ &{x^2+2 3=(x+1)^3}}\\
			&\Leftrightarrow & \heva{&{y=x^2}\\ &{x^3+2x^2+3x-2 2=0}}\\
			&\Leftrightarrow & \heva{&{y=x^2}\\ &{(x-2)(x^2+4x+1 1)=0}}\\
			&\Leftrightarrow & \heva{&x=2=x_0\\ &y=4=y_0.}	
		\end{eqnarray*}	
	}
\end{ex}

%8  
\begin{ex}%[1K6KK-0]
	Tìm tập nghiệm $S$ của hệ phương trình $\heva{&3^x=27\cdot3^y\\ &\log(x+2y)=\log 5+\log 3}$.
	\choice
	{\True $S=\{(7 ; 4)\}$}
	{$S=\{(4 ; 7)\}$}
	{$S=\{(6 ; 3)\}$}
	{$S=\{(9 ; 6)\}$}
	\loigiai{Điều kiện $x+2y>0$.\\
		Hệ phương trình tương đương với  $\heva{&3^x=3^3\cdot 3^y\\ &\log(x+2y)=\log 15}		
		\Leftrightarrow\heva{&{x=y+3}\\ &{x+2y=1 5}}\Leftrightarrow\heva{&x=7\\ &y=4.}$
	}
\end{ex}

%9
\begin{ex}%[1K6KK-0]
	Tìm tất cả các cặp số $(x; y)$ thỏa mãn $\dfrac{4^x}{2^y}=2$ và $\log (2x+2y)=1$.
	\choice
	{$(x ; y)=(4 ; 1)$}
	{\True $(x ; y)=(2 ; 3)$}
	{$(x ; y)=(3 ; 2)$}
	{$(x ; y)=(5 ; 9)$}
	\loigiai{Điều kiện $x+y>0$.\\
		Ta có
		$\heva{& \dfrac{4^x}{2^y}=2 \\ & \log (2x+2y)=1}\heva{& \dfrac{4^x}{2^y}=2 \\ & 2x+2y=10}\Leftrightarrow\heva{& 2x-y=1 \\ & x+y=5}\heva{& x=2 \\ & y=3.}$		
	}
\end{ex}



\begin{ex}%[1K6KK-0]
	Cho hệ phương trình $\heva{&\left(\dfrac 23\right)^{2x-y}+6\left(\dfrac 23\right)^{\tfrac{2x-y}2}-7=0\\ &3^{\log_9(x-y)}=1}$. Chọn khẳng định \textbf{đúng}?
	\choice
	{Điều kiện xác định của hệ phương trình là $x>y>0$}
	{Hệ phương trình đã cho có hai nghiệm $(x ; y)$}
	{\True Hệ phương trình đã cho có một nghiệm duy nhất $(x ; y)=(-1 ;-2)$}
	{Hệ phương trình đã cho vô nghiệm}
	\loigiai{Điều kiện $x-y>0\Leftrightarrow x>y$.\\
		Xét phương trình thứ nhất của hệ $\left(\dfrac 23\right)^{2x-y}+6\left(\dfrac 23\right)^{\tfrac{2x-y}2}-7=0$. Đặt $t=\left(\dfrac 23\right)^{\tfrac{2x-y}2}>0$. \\
		Khi đó ta có 
		$t^2+6t-7=0\Leftrightarrow\hoac{&t=1~(\text{thỏa mãn})\\ &t=-7~(\text{loại}).}$\\
		Do đó $\left(\dfrac 23\right)^{\tfrac{2x-y}2}=1\Leftrightarrow\dfrac{2x-y}2=0\Rightarrow 2x-y=0.$\\
		Xét phương trình thứ hai của hệ $3^{\log_9(x-y)}=1\Leftrightarrow 3^{\log_9(x-y)}=3^0\Leftrightarrow\log_9(x-y)=0\Leftrightarrow x-y=1$.\\
		Từ đó ta có $\heva{&{2x-y=0}\\ &{x-y=1}}\Leftrightarrow\heva{&x=-1\\ &y=-2}$ thỏa mãn điều kiện.\\
		Vậy hệ phương trình đã cho có một nghiệm duy nhất $(x ; y)=(-1 ;-2)$.
	}
\end{ex}





\Closesolutionfile{ans}
\begin{indapan}{10}
{ans/ans-1K6-21-Dang11}
\end{indapan}
\Opensolutionfile{ans}[ans/ans-1K6-21-Dang12]


\begin{dang}{Bất phương trình mũ và logarit bằng phương pháp đặt ẩn phụ}
	Khi giải bất phương trình mũ và logarit bằng phương pháp đặt ẩn phụ ta thường sử dụng một số phương pháp sau:
	\begin{itemize}
		\item Dùng ẩn phụ chuyển bất phương trình thành một bất phương trình với một ẩn phụ.\\
		(Ở đây ta coi các bất phương trình có dạng $f(x) \ge 0$, các trường hợp khác tiến hành tương tự).\\
		* Các phép đặt ẩn phụ với bất phương trình mũ thường gặp:
		\begin{enumerate}[+)]
			\item Bất phương trình có dạng: ${A_n}{a^{nf(x)}} + {A_{n - 1}}{a^{(n - 1)f(x)}} + ... + {A_1}{a^{f(x)}} + {A_0} \ge 0.$\\
			Đặt $t=a^{f(x)} (t>0)$, ta thu được bất phương trình: ${A_n}{t^n} + {A_{n - 1}}{t^{n - 1}} + ... + {A_1}t + {A_0} \ge 0$.
			\item Bất phương trình có dạng: $A{a^{2f(x)}} + B{{(ab)}^{f(x)}} + C{b^{2f(x)}} \ge 0.$\\
			Chia cả hai vế cho $b^{2f(x)}$ rồi đặt $t=\left(\dfrac{a}{b}\right)^{f(x)}$ ta được: $A.t^2+B.t+C\ge 0.$
			\item Bất phương trình có dạng: $A{a^{f(x)}} + B{b}^{g(x)} + C \ge 0$ trong đó, $a^{f(x)}b^{g(x)}=k$.\\
			Đặt $a^{f(x)}=t (t>0) \Rightarrow b^{g(x)}=k$.\\
			Ta thu được bất phương trình mới: $At+\dfrac{Bk}{t}+C \ge 0.$
		\end{enumerate}
		* Các phép đặt ẩn phụ thường gặp với bất phương trình logarit:
		\begin{enumerate}[+)]
			\item Đặt $t=\log_a{x} (x>0)$.
			\item Khi bất phương trình xuất hiện $a^{\log_b{x}}$, ta đặt $t=\log_b{x}$.
		\end{enumerate}
		\item Dùng ẩn phụ chuyển bất phương trình thành một bất phương trình với một ẩn phụ và hệ số vẫn chứa $x$.\\
		Trong trường hợp này, ta thường thu được một bất phương trình bậc hai theo ẩn phụ (hoặc vẫn theo $x$) có $\Delta$ là một số chính phương. 
		\item Dùng ẩn phụ chuyển bất phương trình thành một hệ bất phương trình với hai ẩn phụ.\\
		Bằng việc sử dụng hai ẩn phụ, ta đưa bất phương trình đã cho về một hệ gồm có:
		\begin{enumerate}[+)]
			\item Bất phương trình có được từ bất phương trình đầu bài.
			\item Phương trình (hoặc bất phương trình) có được từ việc đánh giá mối quan hệ của hai ẩn phụ.
		\end{enumerate}
	\end{itemize}
\end{dang}


\begin{vd}%[1K6YK-A]
	Giải các bất phương trình sau
	\begin{listEX}[2]
		\item $16^x-4^x-6\ge 0$.
		\item $32.4^x-18.2^x+1<0$.
		\item $2^x + 2^{3-x} \le 9$.
		\item $25^{\sqrt{x}} + 5 < 5^{\sqrt{x}+1} + 5^{\sqrt{x}}$
	\end{listEX}
	\loigiai{
		\begin{listEX}
			\item Đặt  $4^x=t$ với $t>0$. Khi đó, bất phương trình đề cho được viết lại
			$$t^2-t-6\ge 0
			\Leftrightarrow \hoac{&t\le -2\\&t\ge 3}
			\Leftrightarrow t\ge 3.$$
			Khi đó $x\ge \log_4 3$.
			Vậy tập nghiệm bất phương trình là $\left[ \log_4 3;+\infty\right) $.
			\item Đặt $ t=2^x, (t>0) $. Bất phương trình đã cho trở thành
			$32t^2-18t+1<0\Leftrightarrow \dfrac{1}{16}<t<\dfrac{1}{2}\Leftrightarrow \dfrac{1}{16}<2^x<\dfrac{1}{2}\Leftrightarrow -4<x<-1.$
			Tập nghiệm của bất phương trình đã cho là $ S=(-4 ;-1) $.
			\item $$\begin{array}{cl}
				(1) & \Leftrightarrow 2^x - \dfrac{8}{2^x}-9 \le 0\\
				& \Leftrightarrow \heva{& t=2^x > 0 \\ & t^2-9t+8 \le 0}\\
				& \Leftrightarrow \heva{& t=2^x > 0 \\ & 1 \le t \le 8} \Leftrightarrow 1 \le 2^x \le 8 \Leftrightarrow 0 \le x \le 3
			\end{array}$$
			Vậy tập nghiệm của bất phương trình là $x \in [0;3]$.
			\item Điều kiện: $x \ge 0$. Tập xác định: $\mathscr{D}=\left[ 0; +\infty\right)$.
			$$\begin{aligned}
				(*) & \Leftrightarrow \left(5^{\sqrt{x}} \right)^2 - 6\cdot 5^{\sqrt{x}} + 5 < 0 \\
				& \Leftrightarrow \heva{& t = 5^{\sqrt{x}} > 0 \\ & t^2-6t+5<0} \Leftrightarrow \heva{& t = 5^{\sqrt{x}} > 0 \\ & 1 < t< 5}\\
				& \Leftrightarrow 1 < t < 5 \Leftrightarrow 1 < 5^{\sqrt{x}} < 5 \Leftrightarrow 0 < \sqrt{x} < 1 \Leftrightarrow 0 < x < 1.
			\end{aligned}$$
			Kết hợp với điều kiện, tập nghiệm của bất phương trình là $x \in (0;1)$.
		\end{listEX}
	}
\end{vd}


\begin{vd}%[1K6BK-A]
	Giải các bất phương trình sau
	\begin{listEX}[2]
		\item $8^x+18^x=2\cdot27^x$.
		\item $6.4^{x}- 13.6^{x}+ 6.9^{x}= 0$.
		\item $3.8^{x}+ 4.12^{x}- 18^{x}- 2.27^{x}= 0$.
		\item $2\cdot 4^{\frac{1}{x}}+6^{\frac{1}{x}}=9^{\frac{1}{x}}$.
	\end{listEX}
	\loigiai{
		\begin{listEX}
			\item Tập xác định: $\mathscr{D}=\mathbb{R}$.
			$$\begin{array}{cl}
				(*)& \Leftrightarrow 1 + \left(\dfrac{3}{2}\right)^{2 x}= 2 \cdot \left(\dfrac{3}{2}\right)^{3 x}\\
				& \Leftrightarrow \heva{& t = \left(\dfrac{3}{2}\right)^{x}> 0 \\ & 2t^3-t^2-1=0} \Leftrightarrow \heva{& t = \left(\dfrac{3}{2}\right)^{x}> 0 \\ & 2t^3-t^2-1=0} \Leftrightarrow t= \left(\dfrac{3}{2}\right)^{x}= 1 \Leftrightarrow x= 0.
			\end{array}$$
			Vậy phương trình có một nghiệm là $x=0$.
			\item Tập xác định: $\mathscr{D}=\mathbb{R}$.
			$$\begin{array}{cl}
				(*)&\Leftrightarrow 6 - 13 \cdot \left(\dfrac{3}{2}\right)^{x}+ 6 \cdot \left(\dfrac{3}{2}\right)^{2 x}= 0\\
				& \heva{& t = \left(\dfrac{3}{2}\right)^{x}> 0 \\ & 6t^2-13t+6=0} \Leftrightarrow \hoac{& \left(\dfrac{3}{2}\right)^{x}= \dfrac{2}{3} \\ & \left(\dfrac{3}{2}\right)^{x}= \dfrac{3}{2}} \Leftrightarrow x = \pm 1.
			\end{array}$$
			Vậy phương trình có hai nghiệm: $x=\pm 1$.
			\item Tập xác định: $\mathscr{D}=\mathbb{R}$.
			$$\begin{array}{cl}
				(*)&\Leftrightarrow 3 + 4 \cdot \left(\dfrac{3}{2}\right)^{x}- \left(\dfrac{3}{2}\right)^{2 x}- 2 \cdot \left(\dfrac{3}{2}\right)^{3 x}= 0\\
				& \Leftrightarrow \heva{& t = \left(\dfrac{3}{2}\right)^{x}> 0 \\ & 2 t^{3}+ t^{3}- 4 t - 3 = 0} \Leftrightarrow \hoac{& t=\dfrac{3}{2} \\ & t = \left(\dfrac{3}{2}\right)^{x}=-1 && \text{(loại)}} \Leftrightarrow x= 1.
			\end{array}$$
			Vậy phương trình có nghiệm duy nhất là $x=1$.
			\item 	Điều kiện: $x \ne 0$.
			$$\begin{array}{cl}
				(*) & \Leftrightarrow 2\cdot\left(\dfrac{4}{9} \right)^{\frac{1}{x}}+\left(\dfrac{6}{9} \right)^{\frac{1}{x}} - 1 = 0\\
				& \Leftrightarrow 2 \cdot \left[\left(\dfrac{2}{3}\right)^{\frac{1}{x}}\right]^{2}+ \left(\dfrac{2}{3}\right)^{x}- 1 = 0\\
				& \Leftrightarrow \heva{& t = \left(\dfrac{2}{3}\right)^{x}> 0 \\ & 2t^2+t-1=0} \Leftrightarrow \hoac{& t = \left(\dfrac{2}{3}\right)^{x}= - 1&& \text{(loại)} \\ & t = \left(\frac{2}{3}\right)^{x}= \frac{1}{2}} \Leftrightarrow \left(\dfrac{2}{3}\right)^{x}= \dfrac{1}{2}\Leftrightarrow x = \log_{\frac{2}{3}}\dfrac{1}{2}.
			\end{array}$$
			So với điều kiện, nghiệm của phương trình là $x = \log_{\frac{2}{3}}\dfrac{1}{2}$.
		\end{listEX}
	}
\end{vd}

\begin{vd}[1K6BK-A]
	Giải các bất phương trình sau
	\begin{listEX}[2]
		\item $\log_{2}(x+1)+\log_{x+1}2\geq \dfrac{5}{2}$.
		\item $2\log_{5}x-\log_{x}125<1$.
		\item $\log_{3}x+\log_{3x}27\leq 3$.
		\item $\log_{4}(3^x-1)\cdot\log_{\tfrac{1}{4}}\dfrac{3^x-1}{16}\leq \dfrac{3}{4}$.
	\end{listEX}
	\loigiai{
		\begin{listEX}
			\item Điều kiện: $\heva{&x+1>0\\&x+1\ne 1}\Leftrightarrow\heva{&x>-1\\&x\ne 0}$.\\
			Bất phương trình đã cho tương đương với $\log_{2}(x+1)+\dfrac{1}{\log_{2}(x+1)}-\dfrac{5}{2}\geq 0$.\\
			Đặt $t=\log_{2}(x+1), t\ne 0$. Bất phương trình trở thành
			$$t+\dfrac{1}{t}-\dfrac{5}{2}\geq 0\Leftrightarrow 2t^2-5t+2\geq 0\Leftrightarrow \hoac{&t\leq \dfrac{1}{2}\\& t\geq 2.}$$
			Suy ra
			$$\hoac{&\log_{2}(x+1)\leq \dfrac{1}{2}\\& \log_{2}(x+1)\geq 2}\Leftrightarrow \hoac{&x+1\leq 2^{\tfrac{1}{2}}\\& x+1\geq 2^2}\Leftrightarrow \hoac{&x\leq \sqrt{2}-1\\& x\geq 3.}$$
			Kết hợp với điều kiện, tập nghiệm là $\left(-1;\sqrt{2}-1\right]\cup\left[3;+\infty\right)\setminus\{0\}$.
			\item Điều kiện: $0<x\ne 1$. Bất phương trình tương đương với
			$$2\log_{5}x-\dfrac{1}{\log_{125}x}-1<0\Leftrightarrow 2\log_{5}x-\dfrac{3}{\log_{5}x}-1<0.$$
			Đặt $t=\log_{5}x, t\ne 0$. Bất phương trình trở thành
			$$\dfrac{2t^2-t-3}{t}<0\Leftrightarrow \hoac{&t<-1\\& 0<t<\dfrac{3}{2}.}$$
			Suy ra
			$$\hoac{&\log_{5}x<-1\\& 0<\log_{5}x<\dfrac{3}{2}}\Leftrightarrow \hoac{&x<\dfrac{1}{5}\\& 1<x<5\sqrt{5}.}$$	
			Kết hợp với điều kiện, tập nghiệm của bất phương trình là: $x\in\left(0;\dfrac{1}{5}\right)\cup\left(1;5\sqrt{5}\right)$.
			\item Điều kiện: $0<3x\ne 1\Leftrightarrow 0<x\ne \dfrac{1}{3}$.
			Bất phương trình tương đương với
			$$\log_{3}x+\dfrac{\log_{3}27}{\log_{3}(3x)}\leq 3 \Leftrightarrow \log_{3}x+\dfrac{3}{1+\log_{3}x}\leq 3.$$
			Đặt $t=\log_{3}x, t\ne -1$. Bất phương trình trở thành
			$$t+\dfrac{3}{1+t}\leq 3\Leftrightarrow \dfrac{t^2-2t}{1+t}\leq 0\Leftrightarrow \heva{&t\ne -1\\&(t^2-2t)(1+t)\leq 0}\Leftrightarrow \hoac{&t< -1\\&0\leq t\leq 2.}$$
			Suy ra $$\hoac{&\log_{3}x<-1\\&0\leq \log_{3}x\leq 2}\Leftrightarrow \hoac{&\log_{3}x<\log_{3}\dfrac{1}{3}\\&\log_{3}1\leq \log_{3}x\leq \log_{3}9}\Leftrightarrow \hoac{&x<\dfrac{1}{3}\\&1\leq x\leq 9.}$$
			So sánh với điều kiện, tập nghiệm của bất phương trình là $\left(0;\dfrac{1}{3}\right)\cup[1;9]$.
			\item Điều kiện: $3^x-1> 0\Leftrightarrow 3^x> 1\Leftrightarrow x>0$. Bất phương trình đã cho tương đương với
			$$\log_{4}(3^x-1)\cdot\left[-\log_{4}(3^x-1)+\log_{4}16\right]-\dfrac{3}{4}\leq 0\Leftrightarrow -\log_{4}^{2}(3^x-1)+2\log_{4}(3^x-1)-\dfrac{3}{4}\leq 0.$$
			Đặt $t=\log_{4}(3^x-1)$. Bất phương trình trở thành
			$$4t^2-8t+3\geq 0\Leftrightarrow \hoac{&t\leq \dfrac{1}{2}\\& t\geq \dfrac{3}{2}.}$$
			Suy ra
			$$\hoac{&\log_{4}(3^x-1)\leq \dfrac{1}{2}\\& \log_{4}(3^x-1)\geq \dfrac{3}{2}}
			\Leftrightarrow \hoac{&3^x-1\leq 2\\& 3^x-1\geq 8}\Leftrightarrow \hoac{&3^x\leq 3\\& 3^x\geq 9}\Leftrightarrow \hoac{&x\leq 1\\& x\geq 2.}$$
			Kết hợp với điều kiện, tập nghiệm của bất phương trình là $(0;1]\cup [2;+\infty)$.
		\end{listEX}
	}
\end{vd}

\begin{vd}[1K6BK-A]
	Giải các bất phương trình sau
	\begin{listEX}[2]
		\item $\log_{\tfrac{1}{2}}^2x - 6\log_2x + 8\leq 0$.
		\item $\log_{\tfrac{1}{2}}^2x +\log_{\tfrac{1}{4}}x^2<0$.
		\item $\log_{x} 100-\dfrac{1}{2} \log_{100} x > 0$.
		\item $2\log_5\sqrt{x}-2=\log_x\dfrac{1}{5}$.
	\end{listEX}
	\loigiai{
		\begin{listEX}
			\item Điều kiện $x>0$. Bất phương trình tương đương với $\log_2^2x-6\log_2 x+8\le 0\hfill(1).$\\
			Đặt $t=\log_2 x$, ta được bất phương trình $t^2-6t+8\le 0\Leftrightarrow 2\le t\le 4$.\\
			Với $2\le t\le 4\Leftrightarrow 2\le \log_2 x\le 4\Leftrightarrow 4\le x\le 16$.\\
			Vậy, bất phương trình đã cho có nghiệm $x\in [4;16]$.
			\item Điều kiện $x>0$. Bất phương trình tương đương với
			$$\log_{\tfrac{1}{2}}^2x +\log_{\tfrac{1}{2}}x<0\Leftrightarrow -1<\log_{\tfrac{1}{2}}x<0\Leftrightarrow 2>x>1.$$
			Vậy, bất phương trình có nghiệm $x\in (1;2)$.
			\item Điều kiện $x>0$, $x\neq 1$. Đặt $\log_{100}x=t$, ta được bất phương trình
			$$\dfrac{1}{t}-\dfrac{t}{2}>0\Leftrightarrow \dfrac{2-t^2}{2t}>0\Leftrightarrow \hoac{& t<-\sqrt{2} \\ & 0<t<\sqrt{2}.}$$
			Suy ra
			$$\hoac{& \log_{100}x<-\sqrt{2} \\ & 0<\log_{100}x<\sqrt{2}}\Leftrightarrow \hoac{& 0<x<100^{-\sqrt{2}} \\ & 1<x<100^{\sqrt{2}}.}$$
			Vậy $x \in \left(0;100^{-\sqrt{2}}\right)\cup\left(1; 100^{\sqrt{2}} \right)$.
			\item Điều kiện: $0<x\neq 1$.\\
			Ta có $2\log_5\sqrt{x}-2=\log_x\dfrac{1}{5}\Leftrightarrow \log_5x-2=\log_x\dfrac{1}{5}\Leftrightarrow \log_5x-2-\dfrac{1}{\log_5x}=0$.\\
			Đặt $t=\log_5x$ ta được $t-2-\dfrac{1}{t}=0\Leftrightarrow t^2-2t+1=0\Leftrightarrow t=1\Rightarrow x=5$.
		\end{listEX}
	}
\end{vd}


\begin{vd}[1K6KK-A]
	Giải các bất phương trình sau
	\begin{listEX}[2]
		\item $(\sqrt{2}- 1)^{x}+(\sqrt{2}+ 1)^{x}- 2 \sqrt{2}= 0$.
		\item $(2 + \sqrt{3})^{x}+(2 - \sqrt{3})^{x}= 4$.
		\item $(\sqrt{10}+1)^{\log_3 x}-(\sqrt{10}-1)^{\log_3 x}\ge \dfrac{2x}{3}$. 
		\item $2^{\log_{\sqrt{2}}^2 x}+x^{2 \log_2 x}-20 \leq 0$.
	\end{listEX}
	\loigiai{
		\begin{listEX}
			\item Tập xác định: $\mathscr{D}=\mathbb{R}$.\\
			Ta có $(\sqrt{2}- 1)(\sqrt{2}+ 1)= 1 \Leftrightarrow(\sqrt{2}- 1)^{x}(\sqrt{2}+ 1)^{x}= 1^{x}= 1$.\\
			Đặt $t =(\sqrt{2}+ 1)^{x}> 0 \Rightarrow(\sqrt{2}- 1)^{x}= \dfrac{1}{t}$.
			$$\begin{array}{cl}
				(*)&\Leftrightarrow \dfrac{1}{t}+ t- 2 \sqrt{2}= 0 \Leftrightarrow t^{2}- 2 \sqrt{2}+ 1 = 0\\
				& \Leftrightarrow \hoac{& t=\sqrt{2}+1 \\ & t=\sqrt{2}-1} \Leftrightarrow \hoac{& (\sqrt{2}+ 1)^{x}= \sqrt{2}+ 1 \\ & (\sqrt{2}+ 1)^{x}= \sqrt{2}- 1} \Leftrightarrow \hoac{& x=1 \\ & x=-1.}
			\end{array}$$
			Vậy phương trình có hai nghiệm $x = \pm 1$.
			\item Tập xác định: $\mathscr{D}=\mathbb{R}$.\\
			Ta có $(2 + \sqrt{3})\cdot(2 - \sqrt{3})= 1 \Leftrightarrow(2 + \sqrt{3})^{x} \cdot (2 - \sqrt{3})^{x}= 1^{x}= 1$.\\
			Đặt $t =(2 + \sqrt{3})^{x}> 0 \Rightarrow(2 - \sqrt{3})^{x} = \dfrac{1}{(2 + \sqrt{3})^{x}}= \dfrac{1}{t}> 0$.
			$$(1)\Leftrightarrow t + \dfrac{1}{t}= 4 \Leftrightarrow t^2-4t+1=0 \Leftrightarrow \hoac{& t = 2 + \sqrt{3}> 0 \\ & t = 2 - \sqrt{3}> 0.}$$
			\begin{itemize}
				\item Với $t = 2 + \sqrt{3}$ thì $(2 + \sqrt{3})^{x}= 2 + \sqrt{3}\Leftrightarrow x = 1$.
				\item Với $t = 2- \sqrt{3}$ thì $(2 - \sqrt{3})^{- x}= 2 - \sqrt{3}\Leftrightarrow x = - 1$.
			\end{itemize}
			Vậy phương trình đã cho có hai nghiệm $x=\pm 1$.
			\item Ta có $x=3^{\log_3x}$ nên bất phương trình tương đương với
			$$\left(\dfrac{\sqrt{10}+1}{3}\right)^{\log_3x}-\left(\dfrac{\sqrt{10}-1}{3}\right)^{\log_3x}\ge \dfrac{2}{3}.$$
			Mà $\left(\dfrac{\sqrt{10}+1}{3}\right)^{\log_3x}\cdot \left(\dfrac{\sqrt{10}-1}{3}\right)^{\log_3x}=1$, nên đặt $\left(\dfrac{\sqrt{10}+1}{3}\right)^{\log_3x}=t>0$ ta được
			$$t-\dfrac{1}{t}\ge \dfrac{2}{3}\Leftrightarrow 3t^2-2t-3\ge 0\Rightarrow t\ge \dfrac{\sqrt{10}+1}{3}\Rightarrow \left(\dfrac{\sqrt{10}+1}{3}\right)^{\log_3x}\ge \dfrac{\sqrt{10}+1}{3}\Leftrightarrow \log_3x\ge 1\Leftrightarrow x\ge 3.$$
			Vậy $x \in[3;+\infty)$.
			\item Ta có $2^{\log_{\sqrt{2}}^2 x}=\left(2^{\log_{\sqrt{2}} x}\right)^{\log_{\sqrt{2}} x}=\left(x^2\right)^{\log_{\sqrt{2}} x}=x^{4\log_2 x}$. Đặt $t=x^{2\log_2 x}>0$. Ta được bất phương trình
			\begin{eqnarray*}
				&& t^2+t-20\le 0\Leftrightarrow -5\le t\le 4\\
				& \Rightarrow &  x^{2\log_2 x}\le 4 \Leftrightarrow 2\log_2x\le \log_x 4
				\Leftrightarrow 2\log_2x\le \dfrac{2}{\log_2 x}\Leftrightarrow \log_2^2 x\le 1
				\Leftrightarrow -1\le \log_2 x\le 1\Leftrightarrow \dfrac{1}{2}\le x\le 2.
			\end{eqnarray*}
			Vậy $x\in\left[\dfrac{1}{2};2\right]$.
		\end{listEX}	
	}
\end{vd}

\subsubsection{Bài tập tự luận}

\begin{bt}[1K6YK-A]
	Giải các bất phương trình sau
	\begin{listEX}[2]
		\item $(0{,}4)^x-(2{,}5)^{x+1}>1{,}5$.
		\item $3^{\sqrt{x}}+3^{\sqrt{x}-1}-3^{\sqrt{x}-2}<11$.
		\item $2^x+2^{-x}-3<0$.
		\item $4^{-x+0{,}5}-7\cdot 2^{-x}-4<0$.
	\end{listEX}
	\loigiai{
		\begin{listEX}
			\item Bất phương trình tương đương với $ (0{,}4)^x-2{,}5\cdot \left(\dfrac{1}{0{,}4}\right)^x - 1{,}5 > 0$.\\
			Đặt $u=(0{,}4)^x$, với $u>0$, ta có bất phương trình
			$$u-\dfrac{2{,}5}{u}-1{,5}>0\Leftrightarrow u^2-1{,}5u-2{,}5>0\Leftrightarrow\hoac{&u<-1\\ &u>2{,}5.}$$
			Đối chiếu điều kiện, suy ra $u>2{,}5$, từ đó ta có $(0{,}4)^x>2{,}5\Leftrightarrow x<-1$.\\
			Vậy tập nghiệm của bất phương trình là $S=(-\infty;-1)$.
			\item Điều kiện: $x\ge 0$.\\
			Đặt $u=3^{\sqrt{x}}$, với $u\ge 1$, ta có bất phương trình 
			$$u+\dfrac{u}{3}-\dfrac{u}{9}<11\Leftrightarrow \dfrac{11u}{9}<11\Leftrightarrow u<9.$$
			Với $u<9$, suy ra $3^{\sqrt{x}}<9\Leftrightarrow \sqrt{x}<2\Leftrightarrow x<4$.\\
			Đối chiếu điều kiện, tập nghiệm của bất phương trình là $S=[0;4)$.
			\item Đặt $u=2^x$, với $u>0$, ta có bất phương trình 
			$$u+\dfrac{1}{u}-3<0\Leftrightarrow u^2-3u+1<0\Leftrightarrow \dfrac{3-\sqrt{5}}{2}<u<\dfrac{3+\sqrt{5}}{2}.$$
			Từ đó suy ra $\log_2\dfrac{3-\sqrt{5}}{2}<x<\log_2\dfrac{3+\sqrt{5}}{2}$.\\
			Vậy tập nghiệm của bất phương trình là $S=\left(\log_2\dfrac{3-\sqrt{5}}{2};\log_2\dfrac{3+\sqrt{5}}{2}\right)$
			\item Bất phương trình tương đương với $2\cdot 4^{-x}-7\cdot 2^{-x}-4<0$.\\
			Đặt $u=2^{-x}$, với $u>0$, ta có bất phương trình 
			$$2u^2-7u-4<0\Leftrightarrow -\dfrac{1}{2}<u<4.$$
			Đối chiếu điều kiện, ta có $0<u<4$, từ đó suy ra $2^{-x}<4\Leftrightarrow -x<2\Leftrightarrow x>-2$.\\
			Vậy tập nghiệm của bất phương trình $S=(-2;+\infty)$.
		\end{listEX}
	}
\end{bt}
\begin{bt}[1K6BK-A]
	Giải các bất phương trình sau
	\begin{listEX}[2]
		\item $9^{x^2+x-1}+1\ge 10\cdot 3^{x^2+x-2}$.
		\item $4^x-2^{2(x-1)}+8^{\frac{2}{3}(x-2)}>52$.
		\item $25\cdot 2^x-10^x+5^x>25$.
		\item $49^{\frac{1}{x}}-35^{\frac{1}{x}}\le 25^{\frac{1}{x}}$.
	\end{listEX}
	\loigiai{
		\begin{listEX}
			\item Bất phương trình tương đương với 
			$$9\cdot 3^{2(x^2+x-2)}-10\cdot 3^{x^2+x-2}+1\ge 0.$$
			Đặt $t=3^{x^2+x-2}$, với $t>0$, ta có bất phương trình
			\begin{align*}
				9t^2-10t+1\ge 0\Leftrightarrow \hoac{&t\le \dfrac{1}{9}\\ &t\ge 1}
			\end{align*}
			\begin{itemize}
				\item[•] Với $t\le \dfrac{1}{9}$, suy ra $3^{x^2+x-2}\le 3^{-2}\Leftrightarrow x^2+x\le 0\Leftrightarrow -1\le x\le 0$.
				\item[•] Với $t\ge 1$, suy ra $3^{x^2+x-2}\ge 1\Leftrightarrow x^2+x-2\ge 0\Leftrightarrow \hoac{&x\le -2\\ &x\ge 1.}$
			\end{itemize}
			Vậy tập nghiệm của bất phương trình là $S=(-\infty;-2]\cup [-1;0]\cup [1;+\infty)$.
			\item Bất phương trình tương đương với 
			$$4^x-4^{x-1}+4^{x-2}>52\Leftrightarrow 4^x\left(1-\dfrac{1}{4}+\dfrac{1}{16}\right)>52\Leftrightarrow 4^x>64\Leftrightarrow x>3.$$
			Tập nghiệm của bất phương trình là $S=(3;+\infty)$.
			\item Bất phương trình tương đương với 
			\begin{align*}
				&\ \ 2^x(25-5^x)+(5^x-25)>0\Leftrightarrow (2^x-1)(25-5^x)>0\\
				\Leftrightarrow &\ \heva{&2^x>1\\ &25>5^x}\vee \heva{&2^x<1\\ &25<5^x}\Leftrightarrow \heva{&x>0\\ &x<2}\vee \heva{&x<0\\ &x>2}\Leftrightarrow 0<x<2.
			\end{align*}
			Tập nghiệm của bất phương trình là $S=(0;2)$.
			\item Điều kiện: $x\not=0$.\\
			Bất phương trình tương đương với
			$$\left(\dfrac{7}{5}\right)^{\frac{2}{x}}-\left(\dfrac{7}{5}\right)^{\frac{1}{x}}-1\le 0.$$
			Đặt $t=\left(\dfrac{7}{5}\right)^{\frac{1}{x}}$ với $0<t\not=1$, ta có bất phương trình
			$$t^2-t-1\le 0\Leftrightarrow \dfrac{1-\sqrt{5}}{2}\le t\le \dfrac{1+\sqrt{5}}{2}.$$
			Đối chiếu điều kiện ta có $0<t\le \dfrac{1+\sqrt{5}}{2}$ và $t\not=1$.\\
			Từ đó suy ra $\left(\dfrac{7}{5}\right)^{\frac{1}{x}}\le \dfrac{1+\sqrt{5}}{2}\Leftrightarrow \dfrac{1}{x}\le \log_{\frac{7}{5}}\dfrac{1+\sqrt{5}}{2}\Leftrightarrow x\in (-\infty;0)\cup \left[\log_{\frac{1+\sqrt{5}}{2}}\dfrac{7}{5};+\infty\right)$.\\
			Vậy tập nghiệm của bất phương trình là $S=(-\infty;0)\cup \left[\log_{\frac{1+\sqrt{5}}{2}}\dfrac{7}{5};+\infty\right)$.
		\end{listEX}
	}
\end{bt}



\begin{bt}[1K6BK-A]
	Giải các bất phương trình sau
	\begin{listEX}[2]
		\item $\left(\sqrt{3}+\sqrt{2}\right)^x+\left(\sqrt{3}-\sqrt{2}\right)^x\le 2$.
		\item $\left(3-\sqrt{5}\right)^{2x-x^2}+\left(3+\sqrt{5}\right)^{2x-x^2}\le 2^{1-x^2+2x}$.
	\end{listEX}
	\loigiai{
		\begin{listEX}
			\item Đặt $t=(\sqrt{3}+\sqrt{2})^x$ với $t>0$, suy ra $(\sqrt{3}-\sqrt{2})^x=\dfrac{1}{t}$.\\
			Từ đó ta có bất phương trình
			\begin{center}
				$t+\dfrac{1}{t}\le 2\Leftrightarrow t^2+1\le 2t\Leftrightarrow (t-1)^2\le 0\Leftrightarrow t=1$.
			\end{center}
			Với $t=0$, suy ra $(\sqrt{3}+\sqrt{2})^x=1\Leftrightarrow x=0$.\\
			Vậy nghiệm của bất phương trình là $x=0$.
			\item Chia hai vế của bất phương trình cho  $2^{2x-x^2}>0$, ta được
			$$\left(\dfrac{3-\sqrt{5}}{2}\right)^{2x-x^2}+\left(\dfrac{3+\sqrt{5}}{2}\right)^{2x-x^2}\le 2.$$
			Đặt $t=\left(\dfrac{3+\sqrt{5}}{2}\right)^{2x-x^2}\Rightarrow t>0$ và $\left(\dfrac{3-\sqrt{5}}{2}\right)^{2x-x^2}=\dfrac{1}{t}$. Khi đó bất phương trình trở thành
			$$\dfrac{1}{t}+t\le 2\Leftrightarrow t^2-2t+1\le 0\Leftrightarrow (t-1)^2\le 0\Leftrightarrow t=1.$$
			Suy ra $\left(\dfrac{3+\sqrt{5}}{2}\right)^{2x-x^2}=1\Leftrightarrow 2x-x^2=0\Leftrightarrow \hoac{& x=0 \\ & x=2.}$\\
			Vậy bất phương trình tập nghiệm là $T=\{0;2\}$.
		\end{listEX}
	}
\end{bt}
\begin{bt}[1K6KK-A]
	Giải các bất phương trình sau
	\begin{listEX}[2]
		\item $\log_{4}(3^x-1)\cdot\log_{\tfrac{1}{4}}\dfrac{3^x-1}{16}\leq \dfrac{3}{4}$.
		\item $\log_{2}x+\log_{3}x<1+\log_{2}x\cdot\log_{3}x$.
		\item $\log_{\sqrt{2}}^2\left(8 x^2\right)+ 3\log_{16}(4x)- 2\log_{\sqrt{2}}\left(2x^3\right)<\dfrac{51}{4}$.
		\item $\log_2^4 x-\log_{\tfrac{1}{2}}^2 \left(\dfrac{x^3}{8} \right)+9 \log_2 \left(\dfrac{32}{x^2} \right)<4\log_{\tfrac{1}{2}}^2x$.
	\end{listEX}
	\loigiai{
		\begin{listEX}
			\item Điều kiện: $3^x-1> 0\Leftrightarrow 3^x> 1\Leftrightarrow x>0$. Bất phương trình đã cho tương đương với
			$$\log_{4}(3^x-1)\cdot\left[-\log_{4}(3^x-1)+\log_{4}16\right]-\dfrac{3}{4}\leq 0\Leftrightarrow -\log_{4}^{2}(3^x-1)+2\log_{4}(3^x-1)-\dfrac{3}{4}\leq 0.$$
			Đặt $t=\log_{4}(3^x-1)$. Bất phương trình trở thành
			$$4t^2-8t+3\geq 0\Leftrightarrow \hoac{&t\leq \dfrac{1}{2}\\& t\geq \dfrac{3}{2}.}$$
			Suy ra
			$$\hoac{&\log_{4}(3^x-1)\leq \dfrac{1}{2}\\& \log_{4}(3^x-1)\geq \dfrac{3}{2}}
			\Leftrightarrow \hoac{&3^x-1\leq 2\\& 3^x-1\geq 8}\Leftrightarrow \hoac{&3^x\leq 3\\& 3^x\geq 9}\Leftrightarrow \hoac{&x\leq 1\\& x\geq 2.}$$
			Kết hợp với điều kiện, tập nghiệm của bất phương trình là $(0;1]\cup [2;+\infty)$.
			\item Điều kiện: $x>0$.\\
			Bất phương trình đã cho tương đương với $\log_{2}x+\dfrac{\log_{2}x}{\log_{2}3}-\log_{2}x\cdot\dfrac{\log_{2}x}{\log_{2}3}<0$.\\
			Đặt $t=\log_{2}x$, bất phương trình trở thành
			$$t+\dfrac{t}{\log_{2}3}-1-\dfrac{t^2}{\log_{2}3}<0\Leftrightarrow t^2-(1+\log_{2}3)t+\log_{2}3>0\Leftrightarrow \hoac{&t<1\\& t>\log_{2}3.}$$
			Suy ra $$\hoac{&\log_{2}x<1\\& \log_{2}x>\log_{2}3}\Leftrightarrow \hoac{&0<x<2\\& x>3.}$$
			Kết hợp với điều kiện, tập nghiệm của bất phương trình là $S=(0;2)\cup(3;+\infty)$.
			\item Điều kiện $x>0$. Bất phương trình tương đương với
			$$\left(6+2\log_{\sqrt{2}}x\right)^2+\dfrac{3}{8}\log_{\sqrt{2}}x+\dfrac{3}{2}-\left(4+6\log_{\sqrt{2}}x\right)<\dfrac{51}{4}\hfill(1).$$
			Đặt $t=\log_{\sqrt{2}}x$, ta được bất phương trình
			$$(6+2t)^2+\dfrac{3}{8}t++\dfrac{3}{2}-4-6t<\dfrac{51}{4}\Leftrightarrow 4t^2+\dfrac{147}{8}t+\dfrac{83}{4}<0\Leftrightarrow -\dfrac{83}{32}<t<-2.$$
			Với $-\dfrac{83}{32}<t<-2\Leftrightarrow -\dfrac{83}{32}<\log_{\sqrt{2}}x<-2\Leftrightarrow 2^{-\tfrac{83}{64}}<x<\dfrac{1}{2}$.\\
			Vậy, nghiệm của bất phương trình đã cho là $x\in\left(2^{-\tfrac{83}{64}};\dfrac{1}{2}\right)$.
			\item Điều kiện $x>0$. Bất phương trình tương đương với
			$$\log_2^2x-9(\log_2x-1)^2+9(5-2\log_2x)<4\log_2^2x.$$
			Đặt $\log_2x=t$, ta được phương trình
			$$t^4-13t^2+36<0\Leftrightarrow 4<t^2<9\Leftrightarrow\heva{& \hoac{& t>2 \\ & t<-2} \\ & -3<t<3}\Leftrightarrow \hoac{& -3<t<-2 \\ & 2<t<3.}$$
			Suy ra 
			$$\hoac{& -3<\log_2x<-2 \\ & 2<\log_2x<3}\Leftrightarrow\hoac{& \dfrac{1}{8}<x<\dfrac{1}{4} \\ & 4<x<8.}$$
			Vậy $x\in\left(\dfrac{1}{8};\dfrac{1}{4}\right)\cup(4;8)$.
		\end{listEX}
	}
\end{bt}




\begin{bt}[1K6GK-A]
	Giải các bất phương trình sau
	\begin{listEX}[2]
		\item $2 \left(1+\log_2 x \right) \log_4 x+\log_5 x < 0$.
		\item $(\sqrt{3}+1)^{\log_2 x}+x \cdot(\sqrt{3}-1)^{\log_2 x}=1+x^2$.
	\end{listEX}
	\loigiai{
		\begin{listEX}
			\item Đặt $\log_2 x=t\Rightarrow x=2^t$ ta được
			\begin{eqnarray*}
				&& (1+t)t+t\log_52<0\\
				& \Leftrightarrow & t(t+1+\log_52)<0\\
				& \Leftrightarrow & -1-\log_52<t<0\\
				& \Leftrightarrow & -1-\log_52<\log_2 x<0\\
				& \Leftrightarrow & \dfrac{1}{2\cdot 2^{\log_52}}<x<1.
			\end{eqnarray*}
			Vậy $x\in \left(\dfrac{1}{2\cdot 2^{\log_52}};1\right)$.
			\item Điều kiện $x>0$.\\
			Đặt $u=(\sqrt{3}+1)^{\log_2 x} > 0$; $v=(\sqrt{3}-1)^{\log_2 x} > 0$. Suy ra $uv=2^{\log_2x}=x$, ta được hệ
			$$\heva{& u+xv=1+x^2 \\ & uv=x}\Leftrightarrow \heva{& u+xv=1+x^2 \\ & u\cdot xv=x^2}\Rightarrow u+xv-u\cdot xv=1\Leftrightarrow (u-1)(1-xv)=0\Leftrightarrow \heva{& u=1 \\ & xv=1.}$$
			\begin{itemize}
				\item Với $u=1\Rightarrow (\sqrt{3}-1)^{\log_2x}=1\Leftrightarrow\log_2x=0\Leftrightarrow x=1$ (thỏa mãn).
				\item Với $xv=1\Rightarrow x(\sqrt{3}-1)^{\log_2x}=1\Leftrightarrow x^{\log_2(\sqrt{3}-1)+1}=1\Leftrightarrow \log_2(\sqrt{3}-1)+1=0$ (vô lý).
			\end{itemize}
			Vậy nghiệm phương trình là $x=1$.
		\end{listEX}
	}
\end{bt}




\subsubsection{Bài tập trắc nghiệm}



\begin{ex}%[1K6YK-A]
	Xét bất phương trình $2^{2x}-3\cdot2^{x+2}+32<0$. Nếu đặt $t=2^x$ thì bất phương trình trở thành bất phương trình nào sau đây?
	\choice
	{$t^2-3t+32<0$}
	{\True $t^2-12t+32<0$}
	{$t^2-6t+32<0$}
	{$t^2-16t+32<0$}
	\loigiai
	{
		Ta có $2^{2x}-3\cdot2^{x+2}+32<0 \Leftrightarrow 2^{2x}-12\cdot2^{x}+32<0$. \\
		Đặt $t=2^x$, bất phương trình trở thành $t^2-12t+32<0$.
	}
\end{ex}



\begin{ex}%[1K6BK-A]
	Bất phương trình $2^{x+1}+4\cdot 2^{-x}-9<0$ có bao nhiêu nghiệm nguyên?
	\choice
	{$3$}
	{$4$}
	{\True $2$}
	{Vô số}
	\loigiai{
		Ta có $2^{x+1}+4\cdot 2^{-x}-9<0\Leftrightarrow 2\cdot (2^x)^2-9\cdot 2^x+4<0\Leftrightarrow \dfrac{1}{2}<2^x<4\Leftrightarrow -1<x<2$.\\
		Vì $x$ là số nguyên nên $x\in \left\{0;1\right\}$. Vậy bất phương trình có hai nghiệm nguyên.
	}
\end{ex}






\begin{ex}%[1K6BK-A]
	Tập nghiệm của bất phương trình $32.4^x-18.2^x+1<0$ là \choice
	{$(-4 ; 0)$}
	{\True $ (-4 ;-1) $}
	{$ (-3 ; 1) $}
	{$ $(-5 ;-2)$ $}
	\loigiai{
		Đặt $ t=2^x, (t>0) $. Bất phương trình đã cho trở thành
		\[32t^2-18t+1<0\Leftrightarrow \dfrac{1}{16}<t<\dfrac{1}{2}\Leftrightarrow \dfrac{1}{16}<2^x<\dfrac{1}{2}\Leftrightarrow -4<x<-1.\]
		Tập nghiệm của bất phương trình đã cho là $ S=(-4 ;-1) $.
	}
\end{ex}


\begin{ex}%[1K6BK-A]
	Tập nghiệm của bất phương trình $3\cdot 9^x-10\cdot 3^x+3\le 0$ có dạng $S=[a;b]$ trong đó $a<b$. Giá trị của biểu thức $5b-2a$ bằng
	\choice
	{$\dfrac{43}{3}$}
	{$\dfrac{8}{3}$}
	{\True $7$}
	{$3$}
	\loigiai{
		Đặt $t=3^x>0$, bất phương trình trở thành $3t^2-10t+3\le 0\Leftrightarrow t\in \left[\dfrac{1}{3};3\right]$.\\
		Suy ra $\dfrac{1}{3}\le 3^x\le 3\Leftrightarrow -1\le x\le 1$.\\
		Tập nghiệm của bất phương trình là $[-1;1]$, do đó $a=-1$, $b=1$.\\
		Vậy $5b-2a=5+2=7$.
	}
\end{ex}


\begin{ex}%[1K6BK-A]
	Tập nghiệm của bất phương trình $6\cdot 9^x-13\cdot 6^x+6\cdot 4^x\le 0$ có dạng $S=[a;b]$. Giá trị của biểu thức $a^2+b^2$ bằng
	\choice
	{\True $2$}
	{$4$}
	{$5$}
	{$3$}
	\loigiai{
		Ta có 
		\begin{eqnarray*}
			&&6\cdot 9^x-13\cdot 6^x+6\cdot 4^x\le 0\\
			&\Leftrightarrow &6\cdot\left(\dfrac{3}{2}\right)^{2x}-13\cdot\left(\dfrac{3}{2}\right)^x+6\le 0\\
			&\Leftrightarrow & \dfrac{2}{3}\le \left(\dfrac{3}{2}\right)^x\le \dfrac{3}{2}\\
			&\Leftrightarrow & -1\le x\le 1.
		\end{eqnarray*}
		Suy ra $a=-1,\,b=1$, hay ta có $a^2+b^2=2$.
	}
\end{ex}



\begin{ex}%[1K6BK-A]
	Biết tập nghiệm của bất phương trình $2^x<3-\dfrac{2}{2^x}$ là khoảng $(a;b)$. Tổng $T=a+b$ bằng
	\choice
	{\True $1$}
	{$2$}
	{$3$}
	{$0$}
	\loigiai{Biến đổi bất phương trình ta được
		$$2^x<3-\dfrac{2}{2^x}\Leftrightarrow4^x-3\cdot2^x+2<0\Leftrightarrow1<2^x<2\Leftrightarrow0<x<1.$$
		Suy ra $a=0$, $b=1$ và $T=a+b=1$.}
\end{ex}
\begin{ex}%[1K6KK-A]
	Tập nghiệm của bất phương trình $\left(2+\sqrt{3}\right)^{x^2 - 2x + 1} + \left(2-\sqrt{3}\right)^{x^2 - 2x - 1}\leq\dfrac{4}{2 - \sqrt{3}}$ là đoạn $[a;b]$. Giá trị biểu thức $a+3b$ bằng
	\choice
	{$8+2\sqrt{3}$}
	{$2+4\sqrt{3}$}
	{$2+4\sqrt{2}$}
	{\True $4+2\sqrt{2}$}
	\loigiai
	{
		Ta có
		\begin{eqnarray*}
			&&\left(2+\sqrt{3}\right)^{x^2 - 2x + 1} + \left(2-\sqrt{3}\right)^{x^2 - 2x - 1}\leq\dfrac{4}{2 - \sqrt{3}}\\
			&\Leftrightarrow& \left(2+\sqrt{3}\right)^2\left(2 + \sqrt{3}\right)^{x^2 - 2x - 1} + \dfrac{1}{(2 + \sqrt{3})^{x^2 - 2x - 1}}\leq 4(2+\sqrt{3}).
		\end{eqnarray*}
		Đặt $t=\left(2+\sqrt{3}\right)^{x^2-2x-1}$ ($t>0$). Phương trình đã cho trở thành
		\begin{align*}
			\left(2+\sqrt{3}\right)^2 t+\dfrac{1}{t}\le 4\left(2+\sqrt{3}\right)\Leftrightarrow \left(7+4\sqrt{3}\right)t^2-4\left(2+\sqrt{3}\right)t+1\le 0\Leftrightarrow \dfrac{1}{7+4\sqrt{3}}\le t\le 1. 
		\end{align*}
		Suy ra $\dfrac{1}{7+4\sqrt{3}}\le \left(2+\sqrt{3}\right)^{x^2-2x-1}\le 1\Leftrightarrow -2\le x^2-2x-1\le 0\Leftrightarrow x\in \left[1-\sqrt{2};1+\sqrt{2}\right]$.\\
		Suy ra $a=1-\sqrt{2}$, $b=1+\sqrt{2}\Rightarrow a+3b=4+2\sqrt{2}$.
	}
\end{ex}


\begin{ex}%[1K6KK-A]
	Số nghiệm nguyên dương nhỏ hơn $ 10 $ của bất phương trình $ 3^{x+1} -2^{2x+1}-12^{\tfrac{x}{2}}<0 $ là 
	\choice
	{8}
	{10}
	{7}
	{\True 9}
	\loigiai{
		Ta biến đổi bất phương trình
		\begin{eqnarray*}
			& &3^{x+1} -2^{2x+1}-12^{\tfrac{x}{2}}<0\\
			&\Leftrightarrow& 3\cdot3^x-2\cdot 4^x-\sqrt{12}^x <0\\ &\Leftrightarrow& 3\cdot \left( \dfrac{3}{4}\right)^x-2-\left(\dfrac{\sqrt{12}}{4} \right)^x<0\\  
			&\Leftrightarrow& 3\cdot\left( \dfrac{\sqrt{3}}{2}\right)^{2x}- \left( \dfrac{\sqrt{3}}{2}\right)^{x} -2<0.\quad\quad (1)
		\end{eqnarray*}
		
		Đặt $ t=\left( \dfrac{\sqrt{3}}{2}\right)^{x}$, điều kiện $t>0$.\\
		Bất phương trình (1) trở thành $ 3t^2-t-2<0 \Leftrightarrow -\dfrac{2}{3}<t<1 $.\\
		Kết hợp điều kiện ta được $ 0<t<1 $. Suy ra  $ 0<\left( \dfrac{\sqrt{3}}{2}\right)^{x}<1 \Leftrightarrow x>0 $.\\
		Do đó nghiệm nguyên dương nhỏ hơn $ 10 $ của bất phương trình thuộc tập $  \left\lbrace 1;2;3; \cdots;8;9 \right\rbrace $.\\
		Vậy số nghiệm nguyên dương nhỏ hơn $ 10 $ của bất phương trình là $ 9 $.
	}
\end{ex}


\begin{ex}%[1K6KK-A]
	Biết rằng tập nghiệm của bất phương trình $\left(3+\sqrt{5}\right)^x+\left(3-\sqrt{5}\right)^x < 3\cdot 2^x$ là khoảng $(a;b)$. Hãy tính $b-a$
	\choice
	{\True $2$}
	{$3$}
	{$1$}
	{$4$}
	\loigiai{
		Ta có $\left(3+\sqrt{5}\right)^x+\left(3-\sqrt{5}\right)^x < 3\cdot 2^x \Leftrightarrow \left(\dfrac{3+\sqrt{5}}{2}\right)^x+\left(\dfrac{3-\sqrt{5}}{2}\right)^x < 3$.\\
		Đặt $t=\left(\dfrac{3+\sqrt{5}}{2}\right)^x > 0 \Rightarrow \left(\dfrac{3-\sqrt{5}}{2}\right)^x = \dfrac{1}{t}$.\\
		Bất phương trình đã cho tương đương với $$t + \dfrac{1}{t} < 3 \Leftrightarrow t^2 - 3t + 1 < 0 \Leftrightarrow \dfrac{3-\sqrt{5}}{2} < t < \dfrac{3+\sqrt{5}}{2} \Leftrightarrow -1<x< 1.$$
		Vậy $b-a =2$.
	}
\end{ex}


\begin{ex}%[1K6KK-A]
	Tìm tập nghiệm của bất phương trình $(\sqrt{2}+1)^x-(\sqrt{2}-1)^{x-2}\le 2(\sqrt{2}+1)$.
	\choice
	{$(-\infty;\sqrt{2}]$}
	{$[-2;+\infty)$}
	{\True$(-\infty;2]$}
	{$[-1;1]$}
	\loigiai{
		Từ bất phương trình trên đề bài ta nhân hai vế cho $\sqrt{2}-1 = \dfrac{1}{\sqrt{2} + 1} > 0$ thu được
		\[
		(\sqrt{2}+1)^{x-1}-(\sqrt{2}-1)^{x-1}\le 2 \quad (*)
		\]
		Đặt $t=(\sqrt{2}+1)^{x-1}$ với $t>0$. Bất phương trình đã cho trở thành
		\[
		t-\dfrac{1}{t}\le 2\Leftrightarrow t^2-2t-1\le 0 \Leftrightarrow 1-\sqrt{2}\le t\le 1+\sqrt{2}.
		\]
		Do $t > 0$ nên $t\le 1+\sqrt{2}$. Khi đó $(\sqrt{2}+1)^{x-1}\le 1+\sqrt{2}\Leftrightarrow x-1\le 1\Leftrightarrow x\le2$.
	}
\end{ex}



\begin{ex}%[1K6BK-A]
	Tập nghiệm $ S $ của bất phương trình $ \log_2^2 x - \log_2 x-2>0 $ là
	\choice
	{$ S=(-1;2)$}
	{$ S=(\infty;-1)\cup (2;+\infty)$}
	{\True $ S=\left(0;\dfrac{1}{2}\right) \cup (4;+\infty)$}
	{$ S=\left(\dfrac{1}{2};4\right)$}
	\loigiai{
		Điều kiện xác định của bất phương trình là $ x>0 $.\\
		Bât phương trình đã cho tương đương với	
		$$\hoac{& \log_2 x >2 \\ & \log_2 x <-1}\Leftrightarrow \hoac{& x>4 \\ & x<\dfrac{1}{2}.}$$
		Đối chiếu điều kiện, tập nghiệm của bất phương trình là $  S=\left(0;\dfrac{1}{2}\right) \cup (4;+\infty) $.
	}
\end{ex}


\begin{ex}%[1K6BK-A]
	Tập nghiệm của bất phương trình $\log_2^2x-5\log_2x-6\leq 0$ là
	\choice
	{$S=[64;+\infty)$}
	{\True $S=\left[\dfrac{1}{2};64\right]$}
	{$S=\left(0;\dfrac{1}{2}\right]$}
	{$S=\left(0;\dfrac{1}{2}\right]\cup[64;+\infty)$}
	\loigiai{
		Ta có $\log_2^2x-5\log_2x-6\leq 0\Leftrightarrow-1\leq\log_2x\leq 6\Leftrightarrow\dfrac{1}{2}\leq x\leq 64$.}
\end{ex}

\begin{ex}%[1K6BK-A]
	Cho bất phương trình $\log_2^2 (2x) - 4\log_2 x - 4 \leq 0$. Khi đặt $t = \log_2 x$ thì bất phương trình đã cho trở thành bất phương trình nào sau đây?
	\choice
	{$ t^2 - 4t - 3 \leq 0$}
	{\True $ t^2 - 2t - 3 \leq 0 $}
	{$ t^2 \leq 0 $}
	{ $t^2 - 4t - 4 \leq 0$ }
	\loigiai{
		Ta có $\log_2^2 (2x) - 4\log_2 x - 4 \leq 0 \Leftrightarrow (1+\log_2 x)^2 - 4\log_2 x - 4 \leq 0 \Leftrightarrow \log_2^2 x -2\log_2 x - 3 \leq 0$.\\
		Khi đặt $t = \log_2 x$ thì bất phương trình trở thành $t^2 - 2t - 3 \leq 0$.	
	}
\end{ex}


\begin{ex}%[1K6BK-A]
	Tập nghiệm của bất phương trình $\log_3^2x-2\log_{\sqrt{3}}x+3>0$ là
	\choice
	{$(3;27)$}
	{$(-\infty;3)\cup(27;+\infty)$}
	{\True $(0;3)\cup(27;+\infty)$}
	{$[3;27]$}
	\loigiai
	{
		Bất phương trình đã cho tương đương với
		\[\log_3^2x-4\log_{3}x+3>0\Leftrightarrow\hoac{& \log_3x>3 \\ & \log_3x<1}\Leftrightarrow\hoac{& x>27 \\ & 0<x<3.}\]
		Vậy tập nghiệm của bất phương trình đã cho là $(0;3)\cup(27;+\infty)$.
	}
\end{ex}

\begin{ex}%[1K6BK-A]
	Tìm số nghiệm nguyên của bất phương trình $\log_2^2x-8\log_2\sqrt{x}+3<0$.
	\choice
	{$1$}
	{\True $5$}
	{$4$}
	{$7$}
	\loigiai{
		$$ \log_2^2x-8\log_2\sqrt{x}+3<0\quad (1)$$
		Điều kiện $x>0$.\\
		$(1)\Leftrightarrow \log_2^2x-4\log_2x+3<0\Leftrightarrow 1<\log_2x<3\Leftrightarrow2<x<8.$\\
		Mà $x$ là số nguyên nên $x\in \left\lbrace 3;4;5;6;7\right\rbrace $.\\
		Vậy có $5$ giá trị nguyên là nghiệm của bất phương trình đã cho.
	}
\end{ex}

\begin{ex}%[1K6BK-A]
	Biết tập nghiệm của bất phương trình $\log_{\frac{1}{3}}^2 x+\log_3 (3x)-3 \le 0$ có dạng $S=\left[a;b\right]$. Tính $a+b$
	\choice
	{\True $\dfrac{28}{9}$}
	{$\dfrac{10}{3}$}
	{$3$}
	{$\dfrac{26}{9}$}
	\loigiai{
		Điều kiện $x>0$.\\
		Với điều kiện trên, ta có
		\begin{eqnarray*}
			\log_{\frac{1}{3}}^2 x+\log_3 (3x)-3 \le 0 &\Leftrightarrow& \left(\log_3 x\right)^2+1+\log_3 x-3 \le 0\\
			&\Leftrightarrow& \log_3^2 x+\log_3 x-2 \le 0\\
			&\Leftrightarrow & -2 \le \log_3 x \le 1\\
			&\Leftrightarrow & \dfrac{1}{9} \le x \le 3.
		\end{eqnarray*}
		Kết hợp với điều kiện, suy ra tập nghiệm của bất phương trình là $S=\left[\dfrac{1}{9};3\right]$, do đó $a+b=\dfrac{1}{9}+3=\dfrac{28}{9}$.
	}
\end{ex}











\begin{ex}%[1K6KK-A]
	Tập nghiệm của bất phương trình $4\log_9x+\dfrac{1}{\log_3x}>3$ là
	\choice
	{$S=\left(0;\dfrac{1}{2}\right)\cup \left(1;+\infty\right)$}
	{\True $S=\left(1;\sqrt{3}\right)\cup \left(3;+\infty\right)$}
	{$S=\left(1;+\infty\right)$}
	{$S=\left(3;+\infty\right)$}
	\loigiai{
		Điều kiện xác định $0<x\ne 1$.\\
		Với điều kiện đó ta có 
		\[4\log_9x+\dfrac{1}{\log_3x}>3\Leftrightarrow 2\log_3x+\dfrac{1}{\log_3x}>3.\]	
		Đặt $t=\log_3x$, với $t\ne 0$, ta có bất phương trình
		\[2t+\dfrac{1}{t}>3 \Leftrightarrow  \dfrac{2t^2-3t+1}{t}>0\Leftrightarrow \hoac{&0<t<\dfrac{1}{2}\\&t>1.}\]
		\begin{itemize}
			\item Với $0<t<\dfrac{1}{2}$, ta có $0<\log_3x<\dfrac{1}{2}\Leftrightarrow 1<x<\sqrt{3}$.
			\item Với $t>1$, ta có $\log_3x>1\Leftrightarrow x>3$.
		\end{itemize}
		Từ đó suy ra tập nghiệm của bất phương trình là $S=\left(1;\sqrt{3}\right)\cup \left(3;+\infty\right)$.
	}
\end{ex}

\begin{ex}%[1K6BK-A]
	Tập nghiệm của bất phương trình $\log\limits_2\left(3\cdot 2^x-2\right)<2x$ là 
	\choice
	{$(1;2)$}
	{\True $\left(\log\limits_2\dfrac{2}{3};0\right)\cup (1;+\infty)$}
	{$(-\infty;1)\cup (2;+\infty)$}
	{$(-\infty;0)\cup (1;+\infty)$}
	\loigiai
	{Điều kiện $3\cdot 2^x-2>0\Leftrightarrow 2^x>\dfrac{2}{3}\Leftrightarrow x>\log_2 \dfrac{2}{3}$.\\
		Khi đó, bất phương trình đã cho tương đương với
		$$
		3\cdot 2^x -2<2^{2x}\Leftrightarrow  \left(2^x\right)^2-3\cdot 2^x+2>0\Leftrightarrow\hoac{& 2^x<1 \\ & 2^x>2}\Leftrightarrow\hoac{& x<0 \\ & x>1.}$$
		Đối chiếu với điều kiện, tập nghiệm của bất phương trình đã cho là $\left(\log\limits_2\dfrac{2}{3};0\right)\cup (1;+\infty)$.
	}
\end{ex}


\begin{ex}%[1K6KK-A]
	Biết bất phương trình $\log_3\left(3^x-1\right)\log_{27}\left(3^{x+2}-9\right)\le1$ có tập nghiệm là đoạn $[a;b]$. Tổng $T=a+b$ bằng
	\choice
	{$T=3+\log_3112$}
	{$T=-2+\log_3112$}
	{$T=-2$}
	{\True $T=-3+\log_3112$}
	\loigiai{Biến đổi bất phương trình ta được
		\begin{eqnarray*}
			&&\log_3\left(3^x-1\right)\log_{27}\left(3^{x+2}-9\right)\le1\\
			&\Leftrightarrow&\dfrac{1}{3}\log_3\left(3^x-1\right)\log_3\left(9\left(3^x-1\right)\right)\le1\\
			&\Leftrightarrow&\dfrac{1}{3}\log_3\left(3^x-1\right)\left(2+\log_3\left(3^x-1\right)\right)\le1.
		\end{eqnarray*}
		Đặt $t=\log_3\left(3^x-1\right)$, bất phương trình trở thành
		\begin{eqnarray*}
			&&t(2+t)\le3\Leftrightarrow t^2+2t-3\le0\Leftrightarrow-3\le t\le1\\
			&\Leftrightarrow&\heva{&\log_3\left(3^x-1\right)\ge-3\\&\log_3\left(3^x-1\right)\le1}\Leftrightarrow\heva{&3^x-1\ge\dfrac{1}{27}\\&3^x-1\le3}\\
			&\Leftrightarrow&\heva{&3^x\ge\dfrac{28}{27}\\&3^x\le4}\Leftrightarrow\log_3\dfrac{28}{27}\le x\le\log_34.
		\end{eqnarray*}
		Vậy ta có $T=a+b=\log_3\dfrac{28}{27}+\log_34=\log_3\dfrac{112}{27}=-3+\log_3112$.}
\end{ex}




\begin{ex}%[1K6KK-A]
	Biết rằng bất phương trình $\log_2(5^x+2)+2\log_{5^x+2}2>3$ có tập nghiệm là $S=\left(\log_a b;+\infty\right)$, với $a$, $b$ là các số nguyên dương nhỏ hơn $6$ và $a\ne 1$. Tính $P=2a+3b$.
	\choice
	{\True $P=16$}
	{$P=7$}
	{$P=11$}
	{$P=18$}
	\loigiai{
		Ta có \begin{align*}
			&\log_2(5^x+2)+2\log_{5^x+2}2>3\\
			\Leftrightarrow &\log_2(5^x+2)+2\dfrac{1}{\log_2(5^x+2)}>3\\
			\Leftrightarrow &\log_2^2(5^x+2)-3\log_2(5^x+2)+2>0\\
			\Leftrightarrow &\hoac{&\log_2(5^x+2)<1\\&\log_2(5^x+2)>2}\\
			\Leftrightarrow &\hoac{&5^x+2<2\\&5^x+2>4}\\
			\Leftrightarrow &5^x>2\\
			\Leftrightarrow &x>\log_5 2
		\end{align*}
		Kết hợp với điều kiện $a$, $b$ là các số nguyên dương nhỏ hơn $6$ nên ta tìm được $a=5$, $b=2$.\\
		Và do đó, $P=16$.
		
	}
\end{ex}

\Closesolutionfile{ans}
\begin{indapan}{10}
	{ans/ans-1K6-21-Dang12}
\end{indapan}
\Opensolutionfile{ans}[ans/ans-1K6-21-Dang13]

\begin{dang}{Bất phương trình mũ giải bằng phương pháp logarit hoá}
	Bất phương trình mũ ${a^{f(x)}}>{b^{g(x)}}$ ta cần nhớ:
	\begin{itemize}
		\item Lấy logarit $2$ vế bất phương trình đã cho với cơ số $a > 1$ ta được: ${{\log }_a}{a^{f(x)}}>{{\log }_a}{b^{g(x)}}\Leftrightarrow f(x)>g(x){{\log }_a}b$.
		\item Lấy logarit $2$ vế bất phương trình đã cho với cơ số $0 < a < 1$ ta được: ${{\log }_a}{a^{f(x)}}>{{\log }_a}{b^{g(x)}}\Leftrightarrow f(x)<g(x){{\log }_a}b$.
	\end{itemize}
\end{dang}
\subsubsection{Ví dụ mẫu}
\begin{vd}%[1K6BK-8][TH]%[DCHT Toán 11 - KNTT -Nguyễn Thanh Sang] %[ID6 chương trình mới]
	Giải bất phương trình ${3^{x^2}}<2^x$.
	\dapso{ $0<x<{{\log }_3}2$.}
	\loigiai{
		Ta có: ${3^{x^2}}<2^x\Leftrightarrow {{\log }_3}{3^{x^2}}<{{\log }_3}{2^x}\Leftrightarrow {x^2}-x\log_3 2<0$ $\Leftrightarrow 0<x<{{\log }_3}2$.
	}
\end{vd}
\begin{vd}%[1K6KK-B][VD]%[DCHT Toán 11 - KNTT -Nguyễn Thanh Sang] %[ID6 chương trình mới]
	Giải bất phương trình ${3^{x^2-5x+6}}>{2^{x-2}}$.
	\dapso{$ x>3+{{\log }_3}2,x<2$.}
	\loigiai{
		Logarit cơ số $2$ cả 2 vế ta có:
		\begin{eqnarray*}
			{3^{x^2-5x+6}}>{2^{x-2}}&\Leftrightarrow& {{\log }_3}{3^{x^2-5x+6}}>{{\log }_3}{2^{x-2}}\\ 
			&\Leftrightarrow& \Leftrightarrow {x^2}-5x+6>(x-2){{\log }_3}2 \\ 
			&\Leftrightarrow& (x-2)(x-3-{{\log }_3}2)>0 \\ 
			&\Leftrightarrow& \hoac{&x>3+{{\log }_3}2\\ &x<2.}
		\end{eqnarray*}
	}
\end{vd}
\begin{vd}%[1K6KK-B][VD]%[DCHT Toán 11 - KNTT -Nguyễn Thanh Sang] %[ID6 chương trình mới]
	Giải bất phương trình ${{7.2}^{x^2}}>{{16.7}^{x-1}}$
	\dapso{$ x>2,x<{{\log }_2}7-2$.}
	\loigiai{
		Logarit cơ số $3$ cả 2 vế ta có:\\
		BPT $\Leftrightarrow{2^{{x^2}-4}}>{7^{x-2}}\Leftrightarrow{x^2}-4>(x-2){\log_2}7$\\
		$\Leftrightarrow (x-2)(x+2-{{\log }_2}7)>0 
		\Leftrightarrow	\hoac{&x>2\\ & x<{{\log }_2}7-2.}$}
\end{vd}
\begin{vd}%[1K6KK-B][VD]%[DCHT Toán 11 - KNTT -Nguyễn Thanh Sang] %[ID6 chương trình mới]
	Giải bất phương trình ${2^{x^2-1}}+{2^{x^2+2}}<{3^{x^2}}+{3^{x^2-1}}$
	\dapso{$x>\sqrt{3},x<-\sqrt{3}$.}
	\loigiai{
		BPT $\Leftrightarrow \dfrac{{2^{x^2}}}{2}+{{4.2}^{x^2}}<{3^{x^2}}+\dfrac{{3^{x^2}}}{3}\Leftrightarrow \dfrac{9}{2}{{.2}^{x^2}}<\dfrac{4}{3}{{.3}^{x^2}}$ \\
		$\Leftrightarrow {2^{{x^2}-3}}<{3^{{x^2}-3}}\Leftrightarrow{x^2}-3<\left({{x^2}-3}\right){\log_2}3$\\
		$\Leftrightarrow (x^2-3)(1-{{\log }_2}3)<0\Leftrightarrow {x^2}-3>0\Leftrightarrow 
		\hoac{&x>\sqrt{3}\\ &x<-\sqrt{3}.}$
	}
\end{vd}
\begin{vd}%[1K6KK-B][VD]%[DCHT Toán 11 - KNTT -Nguyễn Thanh Sang] %[ID6 chương trình mới]
	Giải bất phương trình ${2^{x-3}}>{3^{{x^2}-5x+6}}$
	\dapso{$3<x< \log_3 18$.}
	\loigiai{
		Logarit hóa hai vế của phương trình theo cơ số 2 ta được:
		\begin{eqnarray*}
			\left( 3 \right)\Leftrightarrow {{\log }_2}{2^{x-3}}>{{\log }_2}{3^{x^2-5x+6}}\\ 
			&\Leftrightarrow& \left( x-3 \right).\left( 1+2{{\log }_2}3-x{{\log }_2}3 \right)>0 \\ 
			&\Leftrightarrow& 	\hoac{
				& \heva{&x>3\\ & 1+2{{\log }_2}3-x{{\log }_2}3>0}\\
				& \heva{&x<3\\ &1+2{{\log }_2}3-x{{\log }_2}3<0}}\\
			&\Leftrightarrow& 	\hoac{
				& \heva{&x>3\\ & x<\dfrac{1+2{{\log }_2}3}{{{\log }_2}3}}\\
				& \heva{&x<3\\ & x>\dfrac{1+2{{\log }_2}3}{{{\log }_2}3}}}\\
			&\Leftrightarrow& 3<x< \log_3 18.
		\end{eqnarray*}
	}
\end{vd}
\begin{vd}%[1K6GK-B][VDC]%[DCHT Toán 11 - KNTT -Nguyễn Thanh Sang] %[ID6 chương trình mới]
	Có bao nhiêu giá trị nguyên dương của tham số $m$ để tập nghiệm của bất phương trình $\left( {3^{x+3}}-\sqrt{3} \right)\left( {3^x}-4m \right)<0$ chứa không quá $6$ số nguyên.
	\dapso{$20$.}
	\loigiai{
		Bất phương trình tương đương
		$\left( 27.3^x-\sqrt{3} \right)\left( 3^x-4m \right)<0$.\\
		$\Leftrightarrow \left( 3^x-\dfrac{\sqrt{3}}{27} \right)\left( {3^x}-4m \right)<0$.\\
		$\Leftrightarrow \dfrac{\sqrt{3}}{27}<3^x<4m$. \\
		Vì $m\in {\mathbb{N}}^{*}$ nên $4m>\dfrac{\sqrt{3}}{27}$.\\ Suy ra
		${{\log }_3}\dfrac{\sqrt{3}}{27}<x<{{\log }_3}4m
		\Leftrightarrow -\dfrac{5}{2}<x<{{\log }_3}4m$.\\
		Vậy tập nghiệm của bất phương trình là $\left( -\dfrac{5}{2};\,{{\log }_3}4m \right)$.\\
		Vì tập nghiệm chứa không quá $6$ số nguyên nên ${{\log }_3}4m\le 4\Leftrightarrow m\le \dfrac{81}{4}$.\\
		Do $m\in {{\mathbb{N}}^{*}}$ nên $m\in \left\{ 1\,;\,2\,;\,...\,;\,20 \right\}$. Vậy có tất cả $20$ số nguyên dương $m$ thỏa mãn yêu cầu đề bài.
	}
\end{vd}

\subsubsection{Bài tập rèn luyện}
%\centerline{\fcolorbox{red}{yellow!50}{\bf {BÀI TẬP TỰ LUẬN }}}
\begin{bt}%[1K6BK-B][TH]%[DCHT Toán 11 - KNTT -Nguyễn Thanh Sang] %[ID6 chương trình mới]
	Giải bất phương trình $ 3^{x}<\mathrm{e}^{x} $.
	\dapso{$ S=(-\infty;0) $.}
	\loigiai{
		Ta có $ 3^{x}<\mathrm{e}^{x}\Leftrightarrow x\ln 3<x\Leftrightarrow x(\ln 3-1)<0\Leftrightarrow x<0 $ vì $ 3>\mathrm{e} $.
	}
\end{bt}
%\begin{bt}%[1K6BK-B][TH]%[DCHT Toán 11 - KNTT -Nguyễn Thanh Sang] %[ID6 chương trình mới]
%	Giải bất phương trình ${3^{x^2}}<2^x$.
%	\dapso{$0<x<{{\log }_3}2$.}
%	\loigiai{
	%		Ta có ${3^{x^2}}<2^x\Leftrightarrow {{\log }_3}{3^{x^2}}<{{\log }_3}{2^x}\Leftrightarrow {x^2}-x\log {{ & }_3}2<0$ $\Leftrightarrow 0<x<{{\log }_3}2$.
	%	}
%\end{bt}
\begin{bt}%[1K6KK-B][VD]%[DCHT Toán 11 - KNTT -Nguyễn Thanh Sang] %[ID6 chương trình mới]
	Tìm số nghiệm nguyên của bất phương trình $\left(17-12\sqrt{2}\right)^x\geq \left(3+\sqrt{8}\right)^{x^2}$
	\dapso{$3$.}
	\loigiai{
		\begin{eqnarray*}
			&&\left(17-12\sqrt{2}\right)^x\geq \left(3+\sqrt{8}\right)^{x^2}\\
			&\Leftrightarrow &\log_{3+\sqrt{8}} \left(17-12\sqrt{2}\right)^x\geq \log_{3+\sqrt{8}} \left(3+\sqrt{8}\right)^{x^2}\\
			&\Leftrightarrow &-2x\geq x^2\\
			&\Leftrightarrow &x^2+2x\leq 0\\
			&\Leftrightarrow &-2\leq x\leq 0.
		\end{eqnarray*}
		Vì $x\in \mathbb{Z}$ nên $x\in \{-2,-1,0\}$.\\
		Vậy số nghiệm nguyên của bất phương trình là $3$.
	}
\end{bt}
\begin{bt}%[1K6KK-B][VD]%[DCHT Toán 11 - KNTT -Nguyễn Thanh Sang] %[ID6 chương trình mới]
	Số $5 \times 9^{2018}$ viết trong hệ thập phân có bao nhiêu chữ số.
	\dapso{$1927$.}
	\loigiai{
		Số $5 \times 9^{2018}$ viết trong hệ thập phân có $n$ chữ số khi và chỉ khi\\
		$10^{n-1} \le x <10^n \Leftrightarrow n-1 \le 2018 \log 9+ \log 5<n \Rightarrow n=1927$.
	}
\end{bt}
\begin{bt}%[1K6GK-B][VDC]%[DCHT Toán 11 - KNTT -Nguyễn Thanh Sang] %[ID6 chương trình mới]
	Tìm tất cả các giá trị của tham số $m$ sao cho bất phương trình ${3^{x^2}}.{m^{x+1}}-\dfrac{1}{3^x}<0$ có duy nhất một nghiệm nguyên.
	\dapso{$m\in \left[ \dfrac{1}{3};1 \right)\cup \left( 9;27 \right]$.}
	\loigiai{
		Điều kiện: $\heva{&m>0\\ &m\ne 1.}$\\
		Bất phương trình $\Leftrightarrow \,{3^{x^2}}.{m^{x+1}}<\dfrac{1}{3^x}$.\\
		Lấy logarit cơ số 3 hai vế ta được\\
		${{\log }_3}\left( {3^{x^2}}.{m^{x+1}} \right)<{{\log }_3}\left( \dfrac{1}{3^x} \right)$\\
		$\Leftrightarrow \,x^2+\left( x+1 \right){{\log }_3}m<-x$\\
		$\Leftrightarrow \,\left( x+1 \right).\left( x+{{\log }_3}m \right)<0$\\
		$\Leftrightarrow
		\hoac{&-1<x<-{{\log }_3}m\,\,(1)\\ &-1>x>-{{\log }_3}m\,\,(2).}$\\
		(1) có nghiệm nguyên duy nhất $\Leftrightarrow x=0$ thì $0<-{{\log }_3}m\le 1$ $\Leftrightarrow 1>m\ge \dfrac{1}{3}$.\\
		(2) có nghiệm nguyên duy nhất $\Leftrightarrow x=-2$ thì $-2>-{{\log }_3}m\ge -3$ $\Leftrightarrow 9<m\le 27$.
	}
\end{bt}
\centerline{\fcolorbox{red}{yellow!50}{\bf {CÂU HỎI TRẮC NGHIỆM (Tầm 10 - 20 câu theo theo tỉ lệ 4:3:2:1)}}}

%%==========Câu 1
\begin{ex}%[1K6BK-B]%[DCHT Toán 11 - KNTT -Nguyễn Thanh Sang] %[ID6 chương trình mới]
	Tập nghiệm của bất phương trình $2^x>{3^{x+1}}$ là
	\choice
	{ $\varnothing $}
	{\True $\left( -\infty ;{{\log }_{\tfrac{2}{3}}}3 \right)$}
	{ $\left( -\infty ;{{\log }_2}3 \right]$}
	{ $\left( {{\log }_{\tfrac{2}{3}}}3;+\infty \right)$}
	\loigiai
	{
		\begin{itemize}
			\item Cách 1: $2^x>{3^{x+1}}\Leftrightarrow x>{{\log }_2}\left( {3^{x+1}} \right)\Leftrightarrow x>\left( x+1 \right){{\log }_2}3\Leftrightarrow x\left( 1-{{\log }_2}3 \right)>{{\log }_2}3$\\
			$\Leftrightarrow x{{\log }_2}\dfrac{2}{3}>{{\log }_2}3\Leftrightarrow x<\dfrac{{{\log }_2}3}{{{\log }_2}\dfrac{2}{3}}\Leftrightarrow x<{{\log }_{\tfrac{2}{3}}}3$.
			\item Cách 2: $2^x>{3^{x+1}}\Leftrightarrow {{\left( \dfrac{2}{3} \right)}^x}>3\Leftrightarrow x<{{\log }_{\tfrac{2}{3}}}3$.
		\end{itemize}
	}
\end{ex}

%%==========Câu 2
\begin{ex}%[1K6BK-B]
	Tìm tập nghiệm của bất phương trình ${2^{x+1}}>{3^{x+2}}$.
	\choice
	{ $\left( -\infty ;{{\log }_{\tfrac{3}{2}}}\dfrac{9}{2} \right)$}
	{\True $\left( -\infty ;{{\log }_{\tfrac{2}{3}}}\dfrac{9}{2} \right)$}
	{ $\left( -\infty ;{{\log }_{\tfrac{2}{3}}}\dfrac{9}{2} \right]$}
	{ $\left( {{\log }_{tfrac{2}{3}}}\dfrac{9}{2};+\infty \right)$}
	\loigiai{
		Ta có ${2^{x+1}}>{3^{x+2}}\Leftrightarrow x+1>\left( x+2 \right){{\log }_2}3\Leftrightarrow \left( {{\log }_2}3-1 \right)x<1-2{{\log }_2}3$.\\
		$\Leftrightarrow x<{{\log }_{\tfrac{3}{2}}}\dfrac{2}{9}\Leftrightarrow x<{{\log }_{\tfrac{2}{3}}}\dfrac{9}{2}$.
	}
\end{ex}

%%==========Câu 3
\begin{ex}%[1K6BK-B]
	Tập nghiệm của bất phương trình ${2^{x^2}}<{3^{x^2}}$ là
	\choice
	{\True $R\setminus \left\{ 0 \right\}$}
	{ $\left( 0;+\infty \right)$}
	{ $\left( -\infty ;0 \right)$}
	{ $\mathbb{R}$}
	\loigiai{
		Đk: $x\in \mathbb{R}$.\\
		Lấy logrit cơ số 2 hai vế của bpt ta có: $x^2<x^2{{\log }_2}3\Leftrightarrow {x^2}\left( {{\log }_2}3-1 \right)>0\Leftrightarrow x\ne 0$.
	}
\end{ex}

%%==========Câu 4
\begin{ex}%[1K6KK-B]
	Số nghiệm nguyên của bất phương trình $3^x \cdot 5^{x^2}<1$ là
	\choice
	{\True $0$}
	{ $1$}
	{ $2$}
	{ $3$}
	\loigiai{
		Lấy logarit cơ số $3$ cả hai vế phương trình ta có:\\
		${{\log }_3}\left( {3^x} \cdot 5^{x^2} \right)<{{\log }_3}1\Leftrightarrow x+x^2{{\log }_3}5<0
		\Leftrightarrow -\dfrac{1}{{{\log }_3}5}<x<0$.\\
		Kết hợp $x\in \mathbb{Z}$ nên bất phương trình không có nghiệm nguyên.
	}
\end{ex}

%%==========Câu 5
\begin{ex}%[1K6KK-B]
	Tập nghiệm của bất phương trình ${{10}^{x^2}}<{{\text{e}}^x}$ là
	\choice
	{ $\left( 0;\sqrt[10]{\text{e}} \right)$}
	{ $\left( 0;e \right)$}
	{\True $\left( 0;\log e \right)$}
	{ $\left( 0;\ln 10 \right)$}
	\loigiai{
		Ta có ${{10}^{x^2}}<e^x\Leftrightarrow \log \left( {{10}^{x^2}} \right)<\log \left( {e^x} \right)\Leftrightarrow {x^2}-x.\log e<0\Leftrightarrow x\in \left( 0;\log e \right)$.
	}
\end{ex}

%%==========Câu 6
\begin{ex}%[1K6KK-B]
	Bất phương trình $ \left(0{,}2\right)^{x^2}\cdot 2^x\ge \dfrac{2}{5}$ tương đương với bất phương trình nào sau đây?
	\choice
	{$x\ge 1$}
	{$-x^2+x-{{\log }_2}\left( \dfrac{2}{5} \right)\ge 0$}
	{\True $x^2-x{{\log }_5}2+{{\log }_5}2-1\le 0$}
	{ $x^2-x{{\log }_5}2+{{\log }_5}2-1\ge 0$}
	\loigiai{
		${(0{,}2)}^{x^2}{{.2}^x}\ge \dfrac{2}{5}$\\
		$\Leftrightarrow {{\log }_5}\left[ {{\left( 0{,}2 \right)}^{x^2}}{{.2}^x} \right]\ge {{\log }_5}\left( \dfrac{2}{5} \right)$\\
		$\Leftrightarrow {{\log }_5}\left( 0,{2^{x^2}} \right)+{{\log }_5}{2^x}\ge {{\log }_5}2-{{\log }_5}5$\\
		$\Leftrightarrow {x^2}{{\log }_5}0,2+x{{\log }_5}2\ge {{\log }_5}2-1$\\
		$\Leftrightarrow -x^2+x{{\log }_5}2-{{\log }_5}2+1\ge 0$\\
		$\Leftrightarrow {x^2}-x{{\log }_5}2+{{\log }_5}2-1\le 0$.
	}
\end{ex}

%%==========Câu 7
\begin{ex}%[1K6KK-B]
	Tập nghiệm của bất phương trình $4^x+{4^{x+2}}+{4^{x+4}}\ge {5^x}+{5^{x+2}}+{5^{x+4}}$ là
	\choice
	{\True $T=\left( -\infty ;\,{{\log }_{\tfrac{4}{5}}}\dfrac{31}{13} \right]$}
	{ $T=\left[ {{\log }_{\tfrac{4}{5}}}\dfrac{31}{13};\,+\infty \right)$}
	{ $T=\left( -\infty ;\,{{\log }_{\tfrac{4}{5}}}\dfrac{31}{13} \right)$}
	{ $T=\left( {{\log }_{\tfrac{4}{5}}}\dfrac{31}{13};\,+\infty \right)$}
	\loigiai{
		Bất phương trình đã cho tương đương\\
		$4^x+16\cdot {4^x}+256\cdot {4^x}\ge {5^x}+25\cdot {5^x}+625\cdot {5^x} \\
		\Leftrightarrow 273\cdot {4^x}\ge 651\cdot {5^x}$ \\
		$\Leftrightarrow {4^x}\ge \dfrac{31}{13}\cdot {5^x} \\
		\Leftrightarrow {{\left( \dfrac{4}{5} \right)}^x}\ge \dfrac{31}{13} \\
		\Leftrightarrow x\le {{\log }_{\tfrac{4}{5}}}\left( \dfrac{31}{13} \right)$. \\
		Vậy tập nghiệm của bất phương trình đã cho là $T=\left( -\infty ;\,{{\log }_{\tfrac{4}{5}}}\dfrac{31}{13} \right]$.
	}
\end{ex}

%%==========Câu 8
\begin{ex}%[1K6KK-B]
	Tập nghiệm của bất phương trình ${5^{{{\log}_{\tfrac 13}}\left({\dfrac{{x-2}}2}\right)}}<1$ là
	\choice
	{$\left(1;+\infty\right)$}
	{\True $\left(-\infty;0\right)$}
	{$\left(0;2\right)$}
	{$\left({0;+\infty}\right)$}
	\loigiai{
		Điều kiện: $\dfrac{{x-2}}x>0\Leftrightarrow \hoac{&x<0\\ &x>2.}$ \\
		$ {5^{{{\log}_{\tfrac 13}}\left({\dfrac{{x-2}}x}\right)}}<1\Leftrightarrow{\log_{\tfrac 13}}\left({\dfrac{{x-2}}x}\right)<0\Leftrightarrow\dfrac{{x-2}}x>1\Leftrightarrow x<0$.
	}
\end{ex}

%%==========Câu 9
\begin{ex}%[1K6KK-B]
	Tìm số nghiệm nguyên của bất phương trình ${2^{x-2020}}{{.3}^{2022x}}>{3^{x^2+4040}}$.
	\choice
	{ $2020$}
	{\True $2017$}
	{ $2018$}
	{ $2019$}
	\loigiai{
		${2^{x-2020}}{{.3}^{2022x}}>{3^{x^2+4040}} \\
		\Leftrightarrow {2^{x-2020}}>{3^{x^2-2022x+4040}} \\
		\Leftrightarrow {{\log }_3}{2^{x-2020}}>{{\log }_3}{3^{x^2-2022x+4040}} \\
		\Leftrightarrow \left( x-2020 \right){{\log }_3}2>x^2-2022x+4040$ \\
		$\Leftrightarrow \left( x-2020 \right){{\log }_3}2>\left( x-2020 \right)\left( x-2 \right) \\
		\Leftrightarrow \left( x-2020 \right)\left( x-2-{{\log }_3}2 \right)<0 \\
		\Leftrightarrow 2+{{\log }_3}2<x<2020$.\\
		Vậy bất phương trình có tập nghiệm $S=\left( 2+\log_32;2020 \right)$, suy ra tập nghiệm nguyên bất phương trình là $\left\{ 3;4;...;2019 \right\}$. \\
		Do đó bất phương trình có $2017$ nghiệm nguyên.
	}
\end{ex}

%%==========Câu 10
\begin{ex}%[1K6GK-B]
	Có bao nhiêu số nguyên $y\in \left[ -2022;2022 \right]$ để bất phương trình ${{\left( 3x \right)}^{y+\dfrac{{{\log }_3}x}{10}}}\ge {3^{\dfrac{11}{10}{{\log }_3}x}}$ có nghiệm đúng với mọi số thực $x\in \left( 1;9 \right)$?
	\choice
	{ $4044$}
	{ $4026$}
	{\True $2022$}
	{ $2023$}
	\loigiai{
		${{\left( 3x \right)}^{y+\dfrac{{{\log }_3}x}{10}}}\ge {3^{\dfrac{11}{10}{{\log }_3}x}} \\
		\Leftrightarrow {{\log }_3}{{\left( 3x \right)}^{y+\dfrac{{{\log }_3}x}{10}}}\ge {{\log }_3}{3^{\dfrac{11}{10}{{\log }_3}x}} \\
		\Leftrightarrow \left( y+\dfrac{{{\log }_3}x}{10} \right)\left( 1+{{\log }_3}x \right)\ge \dfrac{11}{10}{{\log }_3}x$ \\
		$\Leftrightarrow \left( 10y+{{\log }_3}x \right)\left( 1+{{\log }_3}x \right)\ge 11{{\log }_3}x \\
		\Leftrightarrow 10y\ge \dfrac{11{{\log }_3}x}{1+{{\log }_3}x}-{{\log }_3}x \\
		\Leftrightarrow 10y\ge \dfrac{10{{\log }_3}x-{{\left( {{\log }_3}x \right)}^2}}{1+{{\log }_3}x}$ \\
		Đặt $t={{\log }_3}x\,\,$, vì $x\in \left( 1;9 \right)$ nên $t\in \left( 0;2 \right)$.\\
		Khi đó $10y\ge \dfrac{10t-t^2}{1+t}$.\\
		Xét $f\left( t \right)=\dfrac{10t-t^2}{1+t}\Rightarrow f'\left( t \right)=\dfrac{-t^2-2t+10}{{{\left( 1+t \right)}^2}}>0\,,\,\,\,\,\forall t\in \left( 0;2 \right)$\\
		Vậy để bất phương trình ${{\left( 3x \right)}^{y+\dfrac{{{\log }_3}x}{10}}}\ge {3^{\dfrac{11}{10}{{\log }_3}x}}$ có nghiệm đúng với mọi số thực $x\in \left( 1;9 \right)$ khi và khi bất phương trình $10y\ge \dfrac{10t-t^2}{1+t}$ nghiệm đúng $\forall t\in \left( 0;2 \right)$.\\
		Khi đó ta được: $10y\ge f\left(2\right)=\dfrac{{16}}3\Rightarrow y\ge\dfrac 8{{15}}$.\\
		Vì $y\in \mathbb{Z};\,y\in \left[ -2022;2022 \right]\Rightarrow y\in \left\{ 1;2;\ldots;2022 \right\}$. Vậy có $2022$  giá trị của $y$ thỏa mãn đề bài.
	}
\end{ex}

%%==========Câu 11
\begin{ex}%[1K6GK-B]
	Có bao nhiêu số nguyên $m$ để tập nghiệm của bất phương trình ${3^{x^2}}.{m^{x+1}}-{3^{-x}}<0$ chứa đúng một số nguyên.
	\choice
	{\True $18$}
	{ $19$}
	{ $17$}
	{ $20$}
	\loigiai{
		Điều kiện $m>0$.\\
		Ta có ${3^{x^2}}.{m^{x+1}}-{3^{-x}}<0\Leftrightarrow {3^{x^2}}.{m^{x+1}}<{3^{-x}}\Leftrightarrow {3^{x^2+x}}.{m^{x+1}}<1$\\
		$\Leftrightarrow \left( {x^2}+x \right)+\left( x+1 \right){{\log }_3}m<0\Leftrightarrow \left( x+1 \right)\left( x+{{\log }_3}m \right)<0$.
		\begin{itemize}
			\item 
			Nếu $m>3\Leftrightarrow {{\log }_3}m>1\Rightarrow $tập nghiệm của bất phương trình là $-{{\log }_3} m<x<-1$.\\
			Khi đó, bất phương trình có đúng một nghiệm nguyên khi \\
			$-3\le -{{\log }_3}m<-2\Leftrightarrow 2<{{\log }_3}m\le 3\Leftrightarrow 9<m\le 27$. \\
			Vì $m\in \mathbb{Z}$ nên $m\in \left\{ 10,11,...,27 \right\}$ nên có $18$ số nguyên.
			\item 
			Nếu $0<m<3\Leftrightarrow {{\log }_3}m<1\Rightarrow $tập nghiệm của bất phương trình là $-1<x<-{{\log }_3}m$.\\
			Khi đó, bất phương trình có đúng một nghiệm nguyên khi \\
			$0<-{{\log }_3}m\le 1\Leftrightarrow -1\le {{\log }_3}m<0\Leftrightarrow \dfrac{1}{3}\le m<1$.\\
			Vì $m\in \mathbb{Z}$ nên $m\in \varnothing $.
		\end{itemize}
		Vậy có $18$ số nguyên $m$ thoả mãn yêu cầu bài toán.
	}
\end{ex}

\begin{dang}{Bất phương trình logarit giải bằng phương pháp mũ hoá}
	Giải bất phương trình logarit cơ bản  ${{\log }_a}x>b$ với $a>0,a\ne 1$ ta cần nhớ:
	\begin{itemize}
		\item Nếu $a>1$ thì ${{\log }_a}x>b\Leftrightarrow x>a^b$.
		\item Nếu $0<a<1$ thì ${{\log }_a}x>b\Leftrightarrow 0<x<a^b$.
	\end{itemize}
\end{dang}
\subsubsection{Ví dụ mẫu}
\begin{vd}[TH]%[DCHT Toán 11 - KNTT -Nguyễn Thanh Sang] %[ID6 chương trình mới]
	Giải bất phương trình ${{\log }_3}\left( 36-x^2 \right)\le 3$.
	\dapso{$S=\left( -6;-3 \right]\cup \left[ 3;6 \right)$.}
	\loigiai{
		Điều kiện: $36-x^2>0\Leftrightarrow -6<x<6$.\\
		Ta có: ${{\log }_3}\left( 36-x^2 \right)\le 3\Leftrightarrow 36-x^2\le 27\Leftrightarrow 9-x^2\le 0\Leftrightarrow 
		\heva{&x\le -3\\ &x\ge 3.}$\\
		Kết hợp với điều kiện ta có tập nghiệm của bất phương trình là $S=\left( -6;-3 \right]\cup \left[ 3;6 \right)$.
	}
\end{vd}
\begin{vd}%[1K6BK-B][TH]%[DCHT Toán 11 - KNTT -Nguyễn Thanh Sang] %[ID6 chương trình mới]
	Giải bất phương trình ${{\log }_3}\left( 13-x^2 \right)\ge 2$.
	\dapso{$S=\left[ -2;2 \right]$.}
	\loigiai{
		Điều kiện: $13-x^2>0\Leftrightarrow -\sqrt{13}<x<\sqrt{13}$.\\
		Khi đó ${{\log }_3}\left( 13-x^2 \right)\ge 2\Leftrightarrow 13-x^2\ge 9\Leftrightarrow {x^2}\le 4$ $\Leftrightarrow -2\le x\le 2$\\
		Vậy tập nghiệm của bất phương trình ${{\log }_3}\left( 13-x^2 \right)\ge 2$ là $\left[ -2;2 \right]$.
	}
\end{vd}
\begin{vd}%[1K6BK-B][VD]%[DCHT Toán 11 - KNTT -Nguyễn Thanh Sang] %[ID6 chương trình mới]
	Giải bất phương trình ${{\log }_3}\left( {{\log }_{\tfrac{1}{2}}}x \right)<1$.
	\dapso{$S=\left( \dfrac{1}{8};1 \right)$.}
	\loigiai{
		Điều kiện:
		$	\heva{&{{\log }_{\tfrac{1}{2}}}x>0\\ &x>0}$
		$\Leftrightarrow \heva{&x<1\\ &x>0.}$\\
		Ta có: ${{\log }_3}\left( {{\log }_{\tfrac{1}{2}}}x \right)<1$ 
		$\Leftrightarrow {{\log }_{\tfrac{1}{2}}}x<3^1$ 
		$\Leftrightarrow x>{{\left( \dfrac{1}{2} \right)}^3}$ 
		$\Leftrightarrow x>\dfrac{1}{8}$.\\
		Vậy tập nghiệm của bất phương trình là $S=\left( \dfrac{1}{8};1 \right)$.
	}
\end{vd}
\begin{vd}%[1K6KK-B][VD]%[DCHT Toán 11 - KNTT -Nguyễn Thanh Sang] %[ID6 chương trình mới]
	Tìm số nghiệm nguyên của bất phương trình ${{\log }_5}(x^2-11x+43)<2$.
	\dapso{$6$.}
	\loigiai{
		Điều kiện: $x^2-11x+43>0\,\,\forall x$.\\
		Ta có: ${{\log }_5}(x^2-11x+43)<2\Leftrightarrow {x^2}-11x+43<25\Leftrightarrow {x^2}-11x+18<0\Leftrightarrow 2<x<9$.\\
		Vậy nghiệm của bất phương trình là $S=\left( 2;9 \right)$.\\
		Kết hợp $x\in \mathbb{Z}\Rightarrow x=\left\{ 3;4;5;6;7;8 \right\}$ nên bất phương trình đã cho có $6$ nghiệm nguyên.
	}
\end{vd}
\begin{vd}%[1K6GK-B][VDC]%[DCHT Toán 11 - KNTT -Nguyễn Thanh Sang] %[ID6 chương trình mới]
	Có bao nhiêu số nguyên $x$ sao cho ứng với mỗi $x$ có không quá $242$ số nguyên $y$. thỏa mãn ${{\log }_4}\left( {x^2}+y \right)\ge {{\log }_3}\left( x+y \right)$.
	\dapso{$56$.}
	\loigiai{
		Điều kiện: $\heva{&{x^2}+y>0\\ &x+y>0.}$\\
		Đặt ${{\log }_3}\left( x+y \right)=t$, ta có 
		$\heva{&{x^2}+y\ge {4^t}\\ &x+y=3^t}$
		$\Leftrightarrow $
		$\heva{&x^2-x\ge 4^t-3^t\\ &y=3^t-x.}$\\
		Nhận xét rằng hàm số $f\left( t \right)=4^t-3^t$ đồng biến trên khoảng $\left( 0;+\infty \right)$ và $f\left( t \right)>0$ với mọi $t>0$.\\
		Gọi $n\in \mathbb{Z}$ thỏa $4^n-3^n=x^2-x$, khi đó ta có $t\le n$.\\
		Từ đó, ta có: $-x<y=3^t-x\le {3^n}-x$.\\
		Mặt khác, vì có không quá $242$ số nguyên $y$ thỏa mãn đề bài nên $3^n\le 242\Leftrightarrow n\le {{\log }_3}242$.\\
		Từ đó, suy ra $x^2-x\le {4^{{{\log }_3}242}}-242\Leftrightarrow -27,4\le x\le 28,4$.\\
		Mà $x\in \mathbb{Z}$ nên $x\in \left\{ -27,\,\,-26,\,\,...,\,\,27,\,\,28 \right\}$.\\
		Vậy có $56$ giá trị nguyên của $x$ thỏa yêu cầu đề bài.
	}
\end{vd}

\subsubsection{Bài tập rèn luyện}
%\centerline{\fcolorbox{red}{yellow!50}{\bf {BÀI TẬP TỰ LUẬN }}}
\begin{bt}%[1K6BK-B][TH]%[DCHT Toán 11 - KNTT -Nguyễn Thanh Sang] %[ID6 chương trình mới]
	Giải bất phương trình $\log_{3}(31-x^2)\geq 3$.
	\dapso{$[-2;2]$.}
	\loigiai{
		Điều kiện: $31-x^2>0\Leftrightarrow -\sqrt{31}<x<\sqrt{31}$.\\
		Ta có:
		$\log_{3}(31-x^2)\geq 3\Leftrightarrow 31-x^2\geq 3^3\Leftrightarrow 4\geq x^2\Leftrightarrow -2\leq x\leq 2$.
		Vậy $S=[-2;2]$.
	}
\end{bt}
\begin{bt}%[1K6BK-B][TH]%[DCHT Toán 11 - KNTT -Nguyễn Thanh Sang] %[ID6 chương trình mới]
	Giải bất phương trình $\log_2\left(5-2^{x+1}\right)\ge 1-x$ ta được tập nghiệm là đoạn $[a,b]$. Tính $b-a$.
	\dapso{$2$.}
	\loigiai{
		Điều kiện: $5-2^{x+1}>0\Leftrightarrow x+1>\log_2 5\Leftrightarrow x>\log_2 5-1$.\\
		Bất phương trình tương đương với
		\begin{eqnarray*}
			5-2^{x+1}\ge 2^{1-x} &\Leftrightarrow& 5\cdot 2^x-2\cdot 4^x\ge 2\\
			&\Leftrightarrow& \dfrac{1}{2}\le 2^x\le 2\\
			&\Leftrightarrow& -1\le x\le 1.
		\end{eqnarray*}
		Kết hợp với điều kiện ta có $x\in[-1;1]$, khi đó $a=-1,b=1$ nên $b-a=2$.
	}
\end{bt}
\begin{bt}%[1K6BK-B][TH]%[DCHT Toán 11 - KNTT -Nguyễn Thanh Sang] %[ID6 chương trình mới]
	Giải bất phương trình ${{\log }_2}(1-{{\log }_2}x)<1$.
	\dapso{$\dfrac{1}{2}<x<2$.}
	\loigiai{
		Điều kiện 
		$\heva{&x>0\\ &1-{{\log }_2}x>0}$
		$\Leftrightarrow \heva{&x>0\\ &1-{{\log }_2}x>0}
		\Leftrightarrow \heva{&x>0\\ &x<2}
		\Leftrightarrow 0<x<2$.\\
		Ta có: ${{\log }_2}(1-{{\log }_2}x)<1
		\Leftrightarrow 1-{{\log }_2}x<2
		\Leftrightarrow 1-{{\log }_2}x<2
		\Leftrightarrow {{\log }_2}x>-1
		\Leftrightarrow x>\dfrac{1}{2}$.\\
		Kết hợp với điều kiện ta được nghiệm của bất phương trình $\dfrac{1}{2}<x<2$.
	}
\end{bt}

\begin{bt}%[1K6KK-B][VD]%[DCHT Toán 11 - KNTT -Nguyễn Thanh Sang] %[ID6 chương trình mới]
	Giải bất phương trình ${\log_x}\left[{{{\log}_9}\left({{3^x}-9}\right)}\right]<1$.
	\dapso{$\left({{{\log}_2}10;+\infty}\right)$.}
	\loigiai{
		Điều kiện: $\heva{&0<x\ne 1\\ &{3^x}-9>0\\ &{\log_9}\left({{3^x}-9}\right)>0}\Leftrightarrow{3^x}>10\Leftrightarrow x>{\log_3}10$.	\\
		Khi đó bất phương trình đã cho tương đương với:\\
		$ {\log_4}\left({{2^x}-4}\right)\le x\Leftrightarrow{2^x}-4\le{4^x}\Leftrightarrow{4^x}-{2^x}+4\ge 0\, ~ \forall x\in \mathbb{R}$.\\
		So với điều kiện ta thu được tập nghiệm $\left({{{\log}_2}10;+\infty}\right)$.
	}
\end{bt}
\centerline{\fcolorbox{red}{yellow!50}{\bf {CÂU HỎI TRẮC NGHIỆM (Tầm 10 - 20 câu theo theo tỉ lệ 4:3:2:1)}}}


%%==========Câu 12
\begin{ex}%[1K6BK-B]
	Nghiệm của bất phương trình ${{\log }_2}\left( {3^x}-2 \right)<0$ là
	\choice
	{ $x>1$}
	{ $x<1$}
	{ $0<x<1$}
	{\True ${{\log }_3}2<x<1$}
	\loigiai{
		Điều kiện $3^x-2>0\Leftrightarrow x>{{\log }_3}2$.\\
		Ta có ${{\log }_2}\left( {3^x}-2 \right)<0\Leftrightarrow {3^x}-2<1
		\Leftrightarrow {3^x}<3\Leftrightarrow x<1$.
	}
\end{ex}

%%==========Câu 14
\begin{ex}%[1K6KK-B]% 
	Giải bất phương trình ${{\log }_2}(1-2{{\log }_9}x)<1$.
	\loigiai{
		Điều kiện 
		$\heva{&x>0\\ &1-2{{\log }_9}x>0}$
		$\Leftrightarrow \heva{&x>0\\ &1-{{\log }_3}x>0}
		\Leftrightarrow \heva{&x>0\\ &x<3}
		\Leftrightarrow 0<x<3$.\\
		Ta có: ${{\log }_2}(1-2{{\log }_9}x)<1
		\Leftrightarrow 1-2{{\log }_9}x<2
		\Leftrightarrow 1-{{\log }_3}x<2
		\Leftrightarrow {{\log }_3}x>-1
		\Leftrightarrow x>\dfrac{1}{3}$.\\
		Kết hợp với điều kiện ta được nghiệm của bất phương trình $\dfrac{1}{3}<x<3$.}
\end{ex}


%%==========Câu 15
\begin{ex}%[1K6KK-B]% 
	Tập nghiệm của bất phương trình ${{\log }_3}\left( 18-x^2 \right)\ge 2$ là
	\choice
	{ $\left( -\infty \,;\,3 \right]$}
	{ $\left( 0\,;\,3 \right]$}
	{\True $\left[ -3\,;\,3 \right]$}
	{ $\left( -\infty \,;\,-3 \right]\cup \left[ 3\,;\,+\infty \right)$}
	\loigiai{
		Điều kiện: $18-x^2>0\Leftrightarrow x\in \left( -3\sqrt{2}\,;\,3\sqrt{2} \right)$.\\
		Khi đó ta có: ${{\log }_3}\left( 18-x^2 \right)\ge 2\Leftrightarrow 18-x^2\ge 9\Leftrightarrow -3\le x\le 3$.\\
		Kết hợp với điều kiện ta được tập ngiệm của bất phương trình đã cho là $\left[ -3\,;\,3 \right]$.}
\end{ex}


%%==========Câu 16
\begin{ex}%[1K6KK-B]%
	Số nghiệm nguyên của bất phương trình ${{\log }_{\tfrac{1}{2}}}(x^2+2x-8)\ge -4$ là
	\choice
	{\True $4$}
	{ $5$}
	{ $10$}
	{ $11$}
	\loigiai{
		Điều kiện: $x^2+2x-8>0
		\Leftrightarrow \heva{&x>2\\ &x<-4.}$\\
		BPT $x^2+2x-8\le {{\left( \dfrac{1}{2} \right)}^{-4}}=16
		\Leftrightarrow {x^2}+2x-24\le 0
		\Leftrightarrow -6\le x\le 4$.\\
		So với điều kiện ta có: $ \hoac{&-6\le x<-4 \\ &2<x\le 4.}$\\
		Kết hợp $x\in \mathbb{Z}$ suy ra bất phương trình đã cho có $4$ nghiệm nguyên.}
\end{ex}

%%==========Câu 17
\begin{ex}%[1K6KK-B]% 
	Tổng các nghiệm nguyên của bất phương trình ${{\log }_2}(x^2+3x)\le 2$ là
	\choice
	{ $T=-7$}
	{ $T=-6$}
	{\True $T=-3$}
	{ $T=-4$}
	\loigiai{
		Điều kiện: $x^2+3x>0\Leftrightarrow \heva{&x>0\\ &x<-3.}$\\
		Ta có: ${{\log }_2}(x^2+3x)\le 2
		\Leftrightarrow {x^2}+3x\le 4
		\Leftrightarrow -4\le x\le 1$.\\
		Kết hợp với điều kiện ta có tập nghiệm của bất phương trình là: $S=\left[ -4;-3 \right)\cup \left( 0;1 \right]$.\\
		Kết hợp $x\in \mathbb{Z}\Rightarrow x=\left\{ -4;1 \right\}\Rightarrow T=-3$.}
\end{ex}

%%==========Câu 18
\begin{ex}%[1K6KK-B]% 
	Tổng các nghiệm nguyên của bất phương trình ${{\log }_{\tfrac{1}{2}}}(x^2-4x+6)>-2$
	\choice
	{ $T=7$}
	{\True $T=6$}
	{ $T=5$}
	{ $T=3$}
	\loigiai{
		Điều kiện $x^2-4x+6>0\,\,\forall x$.\\
		Ta có:$\log{_{\tfrac 12}}({x^2}-4x+6)>-2
		\Leftrightarrow{x^2}-4x+6<{\left({\dfrac 12}\right)^{-2}}=4
		\Leftrightarrow{x^2}-4x+2<0
		\Leftrightarrow 2-\sqrt 2<x<2+\sqrt 2 $.\\
		Kết hợp $x\in \mathbb{Z}\Rightarrow x=\left\{ 1;2;3 \right\}\Rightarrow T=6$.}
\end{ex}

%%==========Câu 19
\begin{ex}%[1K6KK-B]% 
	Có tất cả bao nhiêu số nguyên $x$ thỏa mãn bất phương trình ${{\log }_{\tfrac{1}{2}}}\left[ {{\log }_2}\left( 2-x^2 \right) \right]>0$?
	\choice
	{ Vô số}
	{ $1$}
	{\True $0$}
	{ $2$}
	\loigiai{
		Điều kiện $ \heva{&{{\log }_2}\left( 2-x^2 \right)>0\\ & 2-x^2>0}
		\Leftrightarrow \heva{&2-x^2>1\\ & 2-x^2>0}
		\Leftrightarrow 2-x^2>1\Leftrightarrow -1<x<1$.\\
		Khi đó ${{\log }_{\tfrac{1}{2}}}\left[ {{\log }_2}\left( 2-x^2 \right) \right]>0$
		$ \Leftrightarrow{\log_2}\left({2-{x^2}}\right)<1
		\Leftrightarrow 2-x^2<2
		\Leftrightarrow {x^2}>0
		\Leftrightarrow x\ne 0$.\\
		Kết hợp với giả thiết $x$ là số nguyên ta thấy không có số nguyên $x$ nào thỏa mãn bất phương trình đã cho.}
\end{ex}

%%==========Câu 20
\begin{ex}%[1K6GK-B]% 
	Có bao nhiêu số nguyên $x$ sao cho ứng với mỗi $x$ có không quá $127$ số nguyên $y$ thỏa mãn ${{\log }_3}\left( {x^2}+y \right)\ge {{\log }_2}\left( x+y \right)$?
	\choice
	{ $89$}
	{ $46$}
	{ $45$}
	{\True $90$}
	\loigiai{
		Ta có ${{\log }_3}\left( {x^2}+y \right)\ge {{\log }_2}\left( x+y \right)\left( 1 \right)$\\
		Đặt $t=x+y\in \mathbb{N}*$ (do $x,y\in \mathbb{Z},x+y>0$)\\
		$(1)\Leftrightarrow {{\log }_3}\left( {x^2}-x+t \right)\ge {{\log }_2}t\Leftrightarrow g(t)={{\log }_2}t-{{\log }_3}\left( {x^2}-x+t \right)\le 0\left( 2 \right)$ \\
		Đạo hàm ${g}'(t)=\dfrac{1}{t\ln 2}-\dfrac{1}{\left( {x^2}-x+t \right)\ln 3}>0$ với mọi $y$. \\
		Do đó $g\left( t \right)$ đồng biến trên $\left[ 1;+\infty \right)$.\\
		Vì mỗi $x$ nguyên có không quá $127$ giá trị $t\in \mathbb{N}^{*}$ nên ta có\\
		$g(128)>0\Leftrightarrow {{\log }_2}128-{{\log }_3}\left( {x^2}-x+128 \right)>0$
		$\Leftrightarrow {x^2}-x+128<3^7\Leftrightarrow -44,8\le x\le 45,8$.\\
		Như vậy có $90$ giá trị thỏa yêu cầu bài toán.}
\end{ex}

%%==========Câu 21
\begin{ex}%[1K6GK-B]% 
	Có bao nhiêu số nguyên $m$ để bất phương trình ${{\log }_2}\left( {x^2}-2x-m \right)\le 4$ nghiệm đúng với mọi $x\in \left[ 0;4 \right]$.
	\choice
	{ $6$}
	{\True $7$}
	{ $5$}
	{ $9$}
	\loigiai{
		Ta có $\heva{&{x^2}-2x-m>0\\ &{x^2}-2x-m\le 16},\forall x\in\left[{0;4}\right]
		\Leftrightarrow\heva{&m<f(x)={x^2}-2\\&m\ge h(x)={x^2}-2x-16},\forall x\in\left[{0;4}\right]$\\
		$\Leftrightarrow\heva{&m<\mathop{\min}\limits_{\left[{0;4}\right]}f(x)=-1\\ &m\ge\mathop{\max}\limits_{\left[{0;4}\right]}h(x)=-8}\Leftrightarrow-8\le m<-1 $.\\
		Mà $m\in \mathbb{Z}$ nên ta chọn $m\in \left\{ -8;-7;-6;-5;-4;-3;-2 \right\}$.\\
		Vậy có $7$ giá trị nguyên của tham số $m$ thỏa mãn yêu cầu của đề bài.
	}
\end{ex}
\Closesolutionfile{ans}
\begin{indapan}{10}
	{ans/ans-1K6-21-Dang13}
\end{indapan}
\Opensolutionfile{ans}[ans/ans-1K6-21-Dang14]
\begin{dang}{Giải bất phương trình bằng phương pháp hàm số và phương pháp đánh giá}%LÊ VĂN VĨ
	
	\textbf{Để giải bất phương trình bằng phương pháp hàm số ta làm các bước sau}
	\begin{enumerate}[\textbf{Bước} 1.]
		\item Biến đổi bất phương trình về dạng $f(u)>f(v)$, với $u=u(x)$, $v=v(x)$.
		\item  Chứng minh $f(x)$ là hàm đồng biến (hoặc nghịch biến).
		\item  Khi đó  $f(u)>f(v)\Leftrightarrow u(x)>v(x)$ và giải bất phương trình này để tìm $x$.
	\end{enumerate}
	\textbf{Để giải bất phương trình bằng phương pháp hàm số ta làm các bước sau}
	\begin{enumerate}[\textbf{Bước} 1.]
		\item Biến đổi bất phương trình về dạng $f(x)\geq g(x)$.
		\item  Chứng minh $f(x)\leq A$ và $g(x)\geq A$.
		\item  $f(x)\geq g(x) \Leftrightarrow \heva{&f(x)= A\\&g(x)=A}$.
	\end{enumerate}
	\begin{note} 
		\begin{enumerate}
			\item Nếu $f(x)$ và $g(x)$ là các hàm đồng biến thì $f(x)+g(x)$ đồng biến.
			\item Nếu $f(x)$ và $g(x)$ là các hàm đồng biến thì $f(g(x))$ đồng biến.
			\item Nếu $f(x)$ và $g(x)$ là các hàm nghịch biến thì $f(g(x))$ đồng biến.
			\item Nếu $f(x)$ là hàm đồng biến và $g(x)$ là  hàm nghịch biến thì $f(g(x))$ nghịch biến.
		\end{enumerate}
	\end{note}
\end{dang}
\subsubsection{Ví dụ mẫu}
\begin{vd}%[1K6KK-C]
	Giải bất phương trình: $3^x+5^x\geq 17\cdot 2^{x-1}$.
	
	\loigiai{
		Ta có 
		\allowdisplaybreaks
		\begin{eqnarray*}
			&&3^x+5^x\geq 17\cdot 2^{x-1}\Leftrightarrow 3^x+5^x\geq \dfrac{17}{2}\cdot 2^x  \Leftrightarrow \left(\dfrac{3}{2}\right)^x+\left(\dfrac{5}{2}\right)^x \geq \dfrac{17}{2}.
		\end{eqnarray*}
		Xét $f(x)=\left(\dfrac{3}{2}\right)^x+\left(\dfrac{5}{2}\right)^x$, ta có $f(x)$ là hàm đồng biến trên $\mathbb{R}$. Do đó 
		\allowdisplaybreaks
		\begin{eqnarray*}
			&&\left(\dfrac{3}{2}\right)^x+\left(\dfrac{5}{2}\right)^x \geq \dfrac{17}{2} \Leftrightarrow f(x)\geq f(2)\Leftrightarrow x\geq 2.
		\end{eqnarray*}
		Vậy tập nghiệm của bất phương trình là $S=[2;+\infty)$.
	}
\end{vd}

\begin{vd}%[1K6KK-C]%
	Giải bất phương trình ${{\log }_2}x\ge {{\log }_5}\left( x^2+{x^{{{\log }_2}3}} \right).$
	\loigiai{Điều kiện: $\heva{&x>0 \\& x^{2}+x^{\log _{2} 3}>0} \Leftrightarrow x>0.$\\
		Đặt $t=\log _{2}x$ thì $x=2^{t}$. Bất phương trình đã cho trở thành
		\allowdisplaybreaks
		\begin{eqnarray*}
			t \ge\log _{5}\left(4^{t}+3^{t}\right) &\Leftrightarrow& 5^{t} \ge4^{t}+3^{t}\\
			&\Leftrightarrow&\left(\frac{4}{5}\right)^{t}+\left(\frac{3}{5}\right)^{t} \le 1=\left(\frac{4}{5}\right)^{2}+\left(\frac{3}{5}\right)^{2}\\ &\Leftrightarrow& f(t) \leqslant f(2)\;\text{với}\; f(t)=\left(\frac{4}{5}\right)^{t}+\left(\frac{3}{5}\right)^{t}\\
			&\Leftrightarrow& t \ge2\;\left(\text{Vì}\; f(t) \;\text{nghịch biến trên}\; \mathbb{R}).\right.
		\end{eqnarray*}
		Do đó, $\log _{2} x \ge\log _{5}\left(x^{2}+x^{\log _{2} 3}\right) \Leftrightarrow \log _{2} x \ge2 \Leftrightarrow x \ge4$.\\
		Vậy bất phương trình đã cho có tập nghiệm là $S=[4 ;+\infty)$.}
\end{vd}
\subsubsection{Bài tập rèn luyện}
\begin{bt}%[1K6GK-C]
	Giải bất phương trình $\log_2^2x-2046\log_2x+2023^2+x^2-2^{2024}x+2^{2046} \leq 0$.
	\loigiai{Ta có
		\allowdisplaybreaks
		\begin{eqnarray*}
			&&\log_2^2x-2046\log_2x+2023^2+x^2-2^{2024}x+2^{2046} \leq 0\\
			&\Leftrightarrow& \left(\log_2x-2023\right)^2+(x-2^{2023})^2 \leq 0\\
			&\Leftrightarrow& \heva{&\log_2x-2023=0\\&x-2^{2023}=0}\\
			&\Leftrightarrow& x=2^{2023}.
		\end{eqnarray*}
		Vậy bất phương trình có tập nghiệm $S=\left\lbrace 2^{2023} \right\rbrace $.
		
	}
\end{bt}
\begin{bt}%Câu 199%[Trần Quốc Tráng]%[1K6GK-C]
	Giải bất phương trình $
	1-4^x+2^{x+1}\geq 2^x+2^{-x}
	$.
	\loigiai{
		Ta có\\
		$1-4^x+2^{x+1}\geq 2^x+2^{-x} \Leftrightarrow 2-\left(4^x-2 \cdot 2^x+1\right)\geq 2^x+2^{-x}  \Leftrightarrow 2-\left(2^x-1\right)^2>2^x+2^{-x}$.\\
		$2-\left(2^x-1\right)^2 \leq 2$. Dấu \lq\lq $=$\rq\rq \, xảy ra khi $2^x-1=0 \Leftrightarrow x=0$.\\
		$2^x+2^{-x} \geq 2 \sqrt{2^x \cdot 2^{-x}}=2$. Dấu \lq\lq=\rq\rq \, xảy ra khi $2^x=2^{-x}\Leftrightarrow x=0$.\\
		Do đó 
		\allowdisplaybreaks
		\begin{eqnarray*}
			&&2-\left(2^x-1\right)^2>2^x+2^{-x} \Leftrightarrow \heva{&2-\left(2^x-1\right)^2=2\\&2^x+2^{-x}=2 }\Leftrightarrow x=0.
		\end{eqnarray*}
		Vậy bất phương trình có tập nghiệm $S=\left\lbrace 0 \right\rbrace $.
	}
	
\end{bt}
\begin{bt}%[1K6GK-C]
	Giải bất phương trình $\log_3(9-\sqrt{x-1})\geq\log_2\left(x^2-2
	x+5\right)$.
	\loigiai{
		Điều kiện $\heva{&x-1\geq 0\\&9-\sqrt{x-1}>0\\&x^2-2 x+5>0} \Leftrightarrow\heva{&x \geq 1\\&x<82\\&\forall x\in \mathbf{R}} \Leftrightarrow 1\leq x<82$.\\
		Ta có 
		\allowdisplaybreaks
		\begin{eqnarray*}
			&&\log_3(9-\sqrt{x-1}) \leq \log_3 9=2\\
			&& \log_2\left(x^2-2 x+5\right)=\log_2\left[(x-1)^2+4\right] \geq \log_2 4=2
		\end{eqnarray*}
		Do đó \allowdisplaybreaks
		\begin{eqnarray*}
			&&\log_3(9-\sqrt{x-1})\geq\log_2\left(x^2-2
			x+5\right)\Leftrightarrow \heva{&\log_3(9-\sqrt{x-1})=2\\&\log_2\left(x^2-2 x+5\right)=2}\Leftrightarrow x=1.
		\end{eqnarray*}
		Vậy bất phương trình có tập nghiệm $S=\left\lbrace 1 \right\rbrace $.
	}	
	
\end{bt}
\begin{bt}%[1K6GK-C]
	Giải bất phương trình
	$
	2^{x^2+1}\leq 2-\sqrt{x}
	$.
	\loigiai{
		Điều kiện: $x\geq 0$. Ta có 
		\allowdisplaybreaks
		\begin{eqnarray*}
			&&2^{x^2+1}\geq 2^1=2\\ &&2-\sqrt{x}\leq 2.  
		\end{eqnarray*}
		Do đó $
		2^{x^2+1}\leq 2-\sqrt{x}
		\Leftrightarrow
		\heva{&2^{x^2+1}=2\\&2-\sqrt{x}=2}\Leftrightarrow x=0
		$. \\
		Vậy: Bất phương trình có nghiệm duy nhất $x=0$.
		
	}
\end{bt}

\begin{bt}%[1K6GK-C]
	Tìm số nghiệm nguyên trên $[-2024;2024]$ của bất 
	phương trình $3\cdot 2^x+4\cdot 3^x+5\cdot 4^x>6\cdot 5^x$.
	
	\loigiai{
		
		Ta có
		\allowdisplaybreaks
		\begin{eqnarray*}
			3\cdot 2^x+4.3^x+5\cdot 4^x>6\cdot5^x \Leftrightarrow 3\left(\dfrac{2}{5}\right)^x+4\left(\dfrac{3}{5}\right)^x+5\left(\dfrac{4}{5}\right)^x-6>0
		\end{eqnarray*}
		Xét hàm số $f(x)=3\left(\dfrac{2}{5}\right)^x+4\left(\dfrac{3}{5}\right)^x+5\left(\dfrac{4}{5}\right)^x-6, \forall x \in \mathbb{R}$ là hàm số $f(x)$ nghịch biến trên $\mathbb{R}$ suy ra phương trình $f(x)=0$ có nhiều nhất một nghiệm \quad (1).\\
		Mặt khác $f(1) \cdot f(2)=\dfrac{8}{5} \cdot\left(-\dfrac{22}{25}\right)=-\dfrac{176}{125}<0$ nên  $f(x)=0$ có ít nhất một nghiệm $x_0\in (1; 2)$ \quad (2). \\
		Từ (1) và (2), ta suy ra phương trình $f(x)=0$ có duy nhất một nghiệm $x_0\in (1; 2)$. 
		Do dó 
		\allowdisplaybreaks
		\begin{eqnarray*}
			&&3\left(\dfrac{2}{5}\right)^x+4\left(\dfrac{3}{5}\right)^x+5\left(\dfrac{4}{5}\right)^x-6>0\\
			&\Leftrightarrow& f(x)>f(x_0)\Leftrightarrow x<x_0\in (1;2) \\
			&\Rightarrow& x\in \left\lbrace -2024;\, -2023;\, \cdots;\, 0\right\rbrace.
		\end{eqnarray*}
		Vậy bất phương trình đã cho có $0+2024+1=2025$ nghiệm nguyên trên $[-2024;2024]$.
	}
	
\end{bt}

\begin{bt}%[1K6GK-C]
	Gọi $S$ là tập hợp tất cả các giá trị của tham số $m$ để trong tất cả các cặp $(x;y)$ thỏa mãn \linebreak $\log_{x^2+y^2+2}(4x+4y-2) \geq 1$ đồng thời tồn tại duy nhất cặp $(x;y)$ sao cho $4x+3y-2m=0$. Tính tổng các giá trị của $S$.
	
	\loigiai{
		Ta có
		{\allowdisplaybreaks
			\begin{align*}
				&\,\log_{x^2+y^2+2}(4x+4y-2) \geq 1 \\
				\Leftrightarrow&\, \heva{&4x+4y-2>0\\&4x+4y-2 \geq x^2+y^2+2}\\ \Leftrightarrow&\, 4x+4y-2 \geq x^2+y^2+2\\ \Leftrightarrow&\, 4 \geq (x-2)^2+(y-2)^2.\tag{1}
			\end{align*}
		}Do tồn tại duy nhất cặp $(x;y)$ sao cho $4x+3y-2m=0.\hfill(2)$\\
		Nên $\heva{&(x-2)^2+(y-2)^2 \leq 4\\&4x+3y-2m=0}$ có nghiệm duy nhất.\\
		Hay đường thẳng $(2)$ tiếp xúc với hình tròn $(1)$, nghĩa là
		\[d_{\left(I; \Delta \right)}=\dfrac{|8+6-2m|}{5}=2 \Leftrightarrow |14-2m|=10 \Leftrightarrow \hoac{&m=12\\&m=2.}\]
		Vậy tổng các giá trị của $S$ là $14$.}
\end{bt}

\begin{bt}%[1K6KK-C]
	Tìm tập hợp tất cả các giá trị của tham số $m$ để bất phương trình $12^x+(2-m)6^x+3^x>0$ thỏa mãn với mọi $x$ dương.
	
	\loigiai{
		Ta có
		\begin{eqnarray*}
			&& 12^x+(2-m)6^x+3^x>0\\
			&\Leftrightarrow& \dfrac{12^x+(2-m)6^x+3^x}{6^x}>0\\
			&\Leftrightarrow& 2^x+2-m+\dfrac{1}{2^x}>0\\
			&\Leftrightarrow& 2^x+\dfrac{1}{2^x}>m-2.
		\end{eqnarray*}
		Theo bất đẳng thức AM-GM, thì $2^x+\dfrac{1}{2^x}\geq2$ với mọi $x$, dấu bằng xảy ra khi $x=0$.\\ Do đó $2^x+\dfrac{1}{2^x}>2$ với mọi $x>0$. Để bất đẳng thức thỏa mãn với mọi $x>0$ thì $2\ge m-2$, tức là $m\le4$.
	}
\end{bt}

\begin{bt}%[1K6KK-C]
	Tìm số nghiệm nguyên của bất phương trình $2024^{2x^2-4x+9}-2024^{x^2+5x+1}-(x-1)(8-x)<0$.
	
	\loigiai{
		Đặt $\heva{&a=2x^2-4x+9\\&b=x^2+5x+1}\Rightarrow a-b=x^2-9x+8$. Bất phương trình đã cho trở thành
		\allowdisplaybreaks
		\begin{eqnarray*}
			&&2024^a-2024^b+(a-b)<0\Leftrightarrow 2024^a+a<2024^b+b\\& \Leftrightarrow& f(a)<f(b)\Leftrightarrow a<b\Leftrightarrow a-b<0\\
			& \Leftrightarrow& x^2-9x+8<0\Leftrightarrow 1<x<8,
		\end{eqnarray*}
		trong đó  hàm số $f(t)=2024^t+t$ trên $\mathbb{R}$, là hàm   đồng biến trên $\mathbb{R}$. \\
		Vậy các nghiệm nguyên của bất phương trình đã cho là $\{2;3;4;5;6;7\}$.
	}
\end{bt}
\begin{bt}%[1K6KK-C]
	Tìm số nghiệm nguyên của bất phương trình $\log_3 \left( 2x^2-x+9\right) -\log_3 \left( x^2+3\right) <-x^2+x-6$.
	
	\loigiai{
		Đặt $\heva{&a=2x^2-4x+9\\&b=x^2+5x+1}\Rightarrow a-b=x^2-9x+8$. Bất phương trình đã cho trở thành
		\allowdisplaybreaks
		\begin{eqnarray*}
			&&\log_3 \left( 2x^2-x+9\right) -\log_3 \left( x^2+3\right) <-x^2+x-6\\
			&\Leftrightarrow &\log_3 \left( 2x^2-x+9\right)+2x^2-x+9 <\log_3 \left( x^2+3\right) +x^2+3\\& \Leftrightarrow& f(2x^2-x+9)<f(x^2+3)\Leftrightarrow 2x^2-x+9<x^2+3\Leftrightarrow a-b<0\\
			& \Leftrightarrow& x^2-x+6<0\Leftrightarrow -2<x<3,
		\end{eqnarray*}
		trong đó  hàm số $f(t)=\log_3 t+t$ trên $(0;+\infty)$, là hàm   đồng biến trên $(0;+\infty)$. \\
		Vậy các nghiệm nguyên của bất phương trình đã cho là $\{-2;\,0;,\,1;\, 2\}$.
	}
\end{bt}
\begin{bt}%[1K6KK-C]
	Tìm số nghiệm nguyên của bất phương trình $x+\log _2 \sqrt{x+1}+\log _3 \sqrt{x+9}>1$.
	
	\loigiai{
		Điều kiện $x>-1$. Ta có
		\allowdisplaybreaks
		\begin{eqnarray*}
			&&x+\log _2 \sqrt{x+1}+\log _3 \sqrt{x+9}>1	\\& \Leftrightarrow& x+\frac{1}{2} \log _2(x+1)+\frac{1}{2} \log _3(x+9)>1 \\
			&\Leftrightarrow&  g(x)>g(0) \Leftrightarrow x>0,
		\end{eqnarray*}
		trong đó  hàm số $g(x)=2 x+\log _2(x+1) +\log _3(x+9)$ trên $(-1;+\infty)$, là hàm   đồng biến trên $(-1;+\infty)$. \\
		Vậy tập nghiệm  của bất phương trình đã cho là $S=(0 ;+\infty)$.
	}
\end{bt}
\subsubsection{Bài tập trắc nghiệm}
\begin{ex}%[1K6KK-C]
	Tìm nghiệm của bất phương trình: $3^x+2^x>13^{\tfrac{x}{2}}$.
	\choice
	{$x>1$}
	{\True $x<2$}
	{$x>2$}
	{$x<1$}
	\loigiai{
		Ta có 
		\allowdisplaybreaks
		\begin{eqnarray*}
			&&3^x+2^x>13^{\tfrac{x}{2}}\Leftrightarrow 9^{\tfrac{x}{2}}+4^{\tfrac{x}{2}}>13^{\tfrac{x}{2}} \Leftrightarrow \left(\dfrac{9}{13}\right)^{\tfrac{x}{2}}+\left(\dfrac{4}{13}\right)^{\tfrac{x}{2}}>1.
		\end{eqnarray*}
		Xét $f(x)=\left(\dfrac{9}{13}\right)^{\tfrac{x}{2}}+\left(\dfrac{4}{13}\right)^{\tfrac{x}{2}}$, ta có $f(x)$ là hàm nghịch biến trên $\mathbb{R}$. Do đó 
		\allowdisplaybreaks
		\begin{eqnarray*}
			&&\left(\dfrac{9}{13}\right)^{\tfrac{x}{2}}+\left(\dfrac{4}{13}\right)^{\tfrac{x}{2}}>1 \Leftrightarrow f(x)>f(2)\Leftrightarrow x<2.
		\end{eqnarray*}
	}
\end{ex}
\begin{ex}%[1K6KK-C]
	Tìm tập nghiệm của bất phương trình ${{\log }_5}x\ge {{\log }_{13}}\left( x+{x^{{{\log }_5}12}} \right).$
	\choice
	{\True $[4 ;+\infty)$}
	{$(4 ;+\infty)$}
	{$(0 ;4]$}
	{$(0 ;+\infty)$}
	\loigiai{Điều kiện: $\heva{&x>0 \\& x+{x^{{{\log }_5}12}}>0} \Leftrightarrow x>0.$\\
		Đặt $t=\log _{52}x$ thì $x=5^{t}$. Bất phương trình đã cho trở thành
		\allowdisplaybreaks
		\begin{eqnarray*}
			t \ge\log _{13}\left(5^{t}+12^{t}\right) &\Leftrightarrow& 13^{t} \ge 5^{t}+12^{t}\\
			&\Leftrightarrow&\left(\frac{5}{13}\right)^{t}+\left(\frac{12}{13}\right)^{t} \le 1=\left(\dfrac{5}{13}\right)^{2}+\left(\dfrac{12}{13}\right)^{2}\\ &\Leftrightarrow& f(t) \leqslant f(2)\;\text{với}\; f(t)=\left(\dfrac{5}{13}\right)^{t}+\left(\frac{12}{13}\right)^{t}\\
			&\Leftrightarrow& t \ge2\;\left(\text{Vì}\; f(t) \;\text{nghịch biến trên}\; \mathbb{R}).\right.
		\end{eqnarray*}
		Do đó, ${{\log }_5}x\ge {{\log }_{13}}\left( x+{x^{{{\log }_5}12}} \right) \Leftrightarrow \log _{5} x \ge2 \Leftrightarrow x \ge 25$.\\
		Vậy bất phương trình đã cho có tập nghiệm là $S=[25 ;+\infty)$.}
\end{ex}

\begin{ex}%[1K6KK-C]
	Trong tất cả các cặp số thực $(x;y)$ thỏa mãn $\log_{x^2+y^2+3}(2x+2y+5) \geq 1$, có bao nhiêu giá trị thực của $m$ để tồn tại duy nhất cặp $(x;y)$ sao cho $x^2+y^2+4x+6y+13-m=0$?
	\choice
	{$1$}
	{\True $2$}
	{$3$}
	{$0$}
	\loigiai{
		Ta có
		\allowdisplaybreaks
		\begin{eqnarray*}
			&&\log_{x^2+y^2+3}(2x+2y+5) \geq 1\\
			&\Leftrightarrow& 2x+2y+5 \geq x^2+y^2+3\\
			&\Leftrightarrow& x^2+y^2-2x-2y-2 \leq 0\\&\Leftrightarrow& (x-1)^2+(y-1)^2 \leq 4.\quad(1)
		\end{eqnarray*}
		$(1)$ là hình tròn $(C)$ tâm $I_1(1;1)$, bán kính $R_1=2$.\\
		Mặt khác $x^2+y^2+4x+6y+13-m=0\Leftrightarrow (x+2)^2+(y+3)^2=m\quad (2)$.\\
		Với $m=0$, $(2) \Leftrightarrow \heva{&x=-2\\&y=-3}$. Ta thấy $(x;y)=(-2;-3)$ không thỏa mãn bất phương trình $(1)$.\\
		Với $m<0$, không tồn tại cặp $(x;y)$ thỏa mãn $(2)$.\\
		Với $m>0$ thì phương trình $(2)$ là phương trình đường tròn $(C')$ tâm $I_2(-2;-3)$, bán kính $R_2=\sqrt{m}$.
		\begin{center}
			\begin{tikzpicture}[line join = round, line cap = round,>=stealth,font=\footnotesize,scale=0.7]
				\def\xt{-10}\def\xp{6}\def\yd{-11}\def\yt{5}
				\draw[->] (\xt,0)--(\xp,0) node[below]{$x$};
				\draw[->] (0,\yd)--(0,\yt) node[left]{$y$};
				\coordinate (I_1) at (1,1);
				\draw[name path=circle,pattern=north east lines, pattern color=gray!50] (I_1) circle(2);
				\coordinate (I_2) at (-2,-3);
				\path[name path=aa] (I_1)--(I_2);
				\draw [name intersections={of=aa and circle}]
				(intersection-1) coordinate (P);
				\coordinate (Q) at ($(I_1)!-1!(P)$);
				\draw let \p1=($(I_2)-(P)$) in (I_2) circle ({veclen(\x1,\y1)});
				\draw let \p1=($(I_2)-(Q)$) in (I_2) circle ({veclen(\x1,\y1)});
				\foreach \diem/\goc in {I_1/0,I_2/-90} \fill[black](\diem) circle (1pt) ($(\diem)+(\goc:3mm)$) node{$\diem$};
				\draw[dashed] (1,0) node[below]{$1$}--(1,1)--(0,1) node[left]{$1$};
				\draw[dashed] (-2,0) node[above]{$-2$}--(-2,-3)--(0,-3) node[right]{$-3$};
				\foreach \diem in {P,Q} \fill (\diem) circle(1.5pt);
			\end{tikzpicture}
		\end{center}
		Tồn tại duy nhất cặp số $(x;y)$ thỏa mãn hệ $(1)$ và $(2)$ khi và chỉ khi $(C)$ và $(C')$ có một điểm chung duy nhất $\Leftrightarrow$ hình tròn $(C)$ và đường tròn $(C')$ tiếp xúc ngoài với nhau, hoặc hình tròn $(C)$ nằm trong $(C')$ và tiếp xúc trong với nhau $\Leftrightarrow \hoac{&I_1I_2=R_1+R_2\\&I_1I_2=R_2-R_1}\Leftrightarrow \hoac{&5=\sqrt{m}+2\\&5=\sqrt{m}-2}\Leftrightarrow \hoac{&m=9\\&m=49.}$\\
		Vậy có $2$ giá trị của $m$ thỏa mãn yêu cầu bài toán.}
\end{ex}

\begin{ex}%[1K6GK-C]
	Trong các nghiệm $(x;y)$ thỏa mãn bất phương trình $\log_{x^2+2y^2}(2x+y) \geq 1$. Giá trị lớn nhất của biểu thức $T=2x+y$ bằng
	\choice
	{$\dfrac{9}{4}$}
	{\True $\dfrac{9}{2}$}
	{$\dfrac{9}{8}$}
	{$9$}
	\loigiai{
		\textbf{Trường hợp 1:} $x^2+2y^2>1$. Đặt $\sqrt{2}y=z$. Suy ra $ x^2+z^2>1$ \quad (1).
		\allowdisplaybreaks
		\begin{eqnarray*}
			\log_{x^2+2y^2}(2x+y) \geq 1&\Leftrightarrow& 2x+y \geq x^2+2y^2\\
			&\Leftrightarrow& 2x+\dfrac{z}{\sqrt{2}} \geq x^2+z^2\\
			&\Leftrightarrow& (x-1)^2+\left(z-\dfrac{1}{2\sqrt{2}}\right)^2 \leq \dfrac{9}{8}.\quad (2)
		\end{eqnarray*}
		Tập hợp các điểm $M(x;z)$ là miền $(H)$ bao gồm miền ngoài của hình tròn $(C_1) \colon x^2+z^2=1$ và miền trong của hình tròn $(C_2) \colon (x-1)^2+\left(z-\dfrac{1}{2\sqrt{2}}\right)^2=\dfrac{9}{8}$.
		\begin{center}
			\begin{tikzpicture}[line join = round, line cap = round,>=stealth,font=\footnotesize,scale=1]
				\def\xt{-1.5}\def\xp{3.5}\def\yd{-2}\def\yt{2.5}
				\pgfmathsetmacro{\tam}{1/(2*sqrt(2))}
				\pgfmathsetmacro{\r}{3/(2*sqrt(2))}
				\fill[pattern=north east lines] (1,\tam) circle(\r);
				\draw[fill=white] (0,0) circle(1);
				\draw (1,\tam) circle(\r);
				\draw[->] (\xt,0)--(\xp,0) node[below]{$x$};
				\draw[->] (0,\yd)--(0,\yt) node[left]{$z$};
				\clip (\xt+0.1,\yd+0.1) rectangle (\xp-0.1,\yt-0.1);
				\draw[samples=200,domain=\xt:\xp,smooth] plot (\x, {-2*sqrt(2)*(\x)+9/sqrt(2)});
				\node at (2.5,-1)[right]{$d$};
				\node at (-1,-1)[]{$(C_1)$};
				\node at (1,1.5)[above]{$(C_2)$};
				\fill[white] (1.5,0.5) circle (7pt) node[black]{$(H)$};
				\fill (0,0) circle (1pt) node[below left]{$O$};
			\end{tikzpicture}
		\end{center}
		Hệ $\heva{&T=2x+\dfrac{z}{\sqrt{2}}\\&(x-1)^2+\left(z-\dfrac{1}{2\sqrt{2}}\right)^2 \leq \dfrac{9}{8}\\&x^2+z^2>1}$ có nghiệm khi đường thẳng $d \colon 2x+\dfrac{z}{\sqrt{2}}-T=0$ có điểm chung với miền $(H)$.\\
		Để $T$ đạt giá trị lớn nhất thì đường thẳng $d \colon 2x+\dfrac{z}{\sqrt{2}}-T=0$ tiếp xúc với đường tròn $(C_2)$\\
		$\Leftrightarrow \mathrm{d}(I;d)=\dfrac{3}{2\sqrt{2}}$ với $I\left(1;\dfrac{1}{2\sqrt{2}}\right)$ là tâm của đường tròn $(C_2)$.\\
		$\Leftrightarrow \dfrac{\left|2+\dfrac{1}{4}-T\right|}{\sqrt{4+\dfrac{1}{2}}}=\dfrac{3}{2\sqrt{2}}\Leftrightarrow \left|T-\dfrac{9}{4}\right|=\dfrac{9}{4}\Leftrightarrow \hoac{&T=0 (l)\\&T=\dfrac{9}{2}}$\\
		\textbf{Trường hợp 2:} $0<x^2+2y^2<1$.\\
		$\log_{x^2+2y^2}(2x+y) \geq 1\Leftrightarrow 2x+y \leq x^2+2y^2\Leftrightarrow T=2x+y<1$ (loại).\\
		Vậy giá trị lớn nhất của biểu thức $T=2x+y$ là $\max T=\dfrac{9}{2}$.}
\end{ex}

\begin{ex}%[1K6GK-C]
	Số nghiệm nguyên trên $[-2024;2024]$  của bất phương trình $3^{\log_7(x+4)}\geq x$ là
	\choice
	{\True $5$}
	{$1$}
	{$2$}
	{$3$}
	\loigiai{	
		Điều kiện  $x>-4$.
		\begin{enumerate}[Trường hợp 1.]
			\item $-4<x\leq 0$ thì bất phương trình luôn thỏa mãn.
			\item $x>0$, ta có 
			\allowdisplaybreaks
			\begin{eqnarray*}
				3^{\log_7(x+4)}\geq x\Leftrightarrow\log_7(x+4)\geq \log_3 x.
			\end{eqnarray*}
			Đặt $\log_3 x=t\Rightarrow x=3^t$.		Ta có 
			\allowdisplaybreaks
			\begin{eqnarray*}
				&&\log_7(x+4)\geq \log_3 x\Leftrightarrow \log_7\left( 3^t+4\right) \geq t\\
				& \Leftrightarrow& 3^t+4\geq 7^t\Leftrightarrow \left(\dfrac{3}{7}\right)^t+4\left(\dfrac{1}{7}\right)^t\geq 1\\
				&\Leftrightarrow& f(x)\geq f(1) \Leftrightarrow x\leq 1,
			\end{eqnarray*}
			trong đó $f(t)=\left(\dfrac{3}{7}\right)^t+4\left(\dfrac{1}{7}\right)^t, t \in \mathbb{R}$ là hàm 
			nghịch biến trên  $\mathbb{R}$.
		\end{enumerate}	
		Suy ra bất  phương trình có tập  nghiệm là $S=(-4;1]$. Do đó số  nghiệm nguyên trên $[-2024;2024]$  của bất phương trình $3^{\log_7(x+4)}\geq x$ là $5$.}
	
\end{ex}
\begin{ex}%[1K6GK-C]
	Tập nghiệm của bất phương trình $\log _{2}\left(x \sqrt{x^{2}+2}+4-x^{2}\right)+2 x+\sqrt{x^{2}+2} \leq 1$ là $(-\sqrt{a} ;-\sqrt{b}]$.
	Khi đó tích  $ab$  bằng
	\choice
	{ $\dfrac{12}{5}$}
	{ $\dfrac{5}{12}$}
	{$\dfrac{15}{16}$}
	{ \True   $\dfrac{16}{15}$}
	\loigiai{Ta có $$x(\sqrt{x^{2}+2}-x)+4=\dfrac{2 x}{\sqrt{x^{2}+2}+x}+4=\dfrac{2(3 x+2 \sqrt{x^{2}+2})}{\sqrt{x^{2}+2}+x}.$$
		Nhận xét  $\sqrt{x^{2}+2}+x>|x|+x \geq 0, \;  \forall x \in \mathbb{R}$.\\
		Khi đó $$x \sqrt{x^{2}+2}+4-x^{2}>0 \Leftrightarrow 2 \sqrt{x^{2}+2}>-3 x \quad \quad \quad (*)$$
		Với điều kiện (*) bất phương trình đã cho tương đương
		\allowdisplaybreaks
		\begin{eqnarray*}
			& &\log _{2} \frac{2(3 x+2 \sqrt{x^{2}+2})}{\sqrt{x^{2}+2}+x}+2 x+\sqrt{x^{2}+2} \leq 1\\
			&\Leftrightarrow&  1+\log _{2}(3 x+2 \sqrt{x^{2}+2})-\log _{2}(x+\sqrt{x^{2}+2})+2 x+\sqrt{x^{2}+2} \leq 1\\
			&\Leftrightarrow& 3 x+2 \sqrt{x^{2}+2}+\log _{2}(3 x+2 \sqrt{x^{2}+2}) \leq x+\sqrt{x^{2}+2}+\log _{2}(x+\sqrt{x^{2}+2})\\
			&\Leftrightarrow&f(3 x+2 \sqrt{x^{2}+2}) \leq f(x+\sqrt{x^{2}+2}) \quad \quad \quad (1)
		\end{eqnarray*}
		Xét hàm số $f(t)=t+\log _{2} t$ trên khoảng  $(0 ;+\infty)$. \\
		Ta có
		$f(t)$ luôn đồng biến trên khoảng  $(0 ;+\infty)$.\\
		Do đó
		\allowdisplaybreaks
		\begin{eqnarray*}
			(1) &\Leftrightarrow & 3 x+2 \sqrt{x^{2}+2} \leq x+\sqrt{x^{2}+2}\\
			&\Leftrightarrow& \sqrt{x^{2}+2} \leq-2 x\\
			&\Leftrightarrow&\heva{&x \leq 0 \\& x^{2}+2 \leq 4 x^{2}} \\
			&\Leftrightarrow&\heva{& \hoac{&x\leq -\sqrt{\dfrac{2}{3}} \\&x \geq \sqrt{\dfrac{2}{3}}} \\&x\leq 0}\\
			&\Leftrightarrow&x \leq -\sqrt{\dfrac{2}{3}}.
		\end{eqnarray*}
		Với điều  kiện trên thì  (*) $\Leftrightarrow 4\left(x^{2}+2\right)>9 x^{2} \Leftrightarrow-\sqrt{\dfrac{8}{5}}<x<\sqrt{\dfrac{8}{5}}$.\\
		Kết hợp (**) ta được tập nghiệm của bất phương trình (1)  là $\left(-\sqrt{\dfrac{8}{5}} ;-\sqrt{\dfrac{2}{3}}\right]$.\\
		Vậy $a b=\dfrac{16}{15}$.
	}
\end{ex}

\begin{ex}%[1K6KK-C]
	Bất phương trình $\log_6\left(\sqrt[3]{x}+\sqrt[6]{x}\right) \geq \log_{64}x$ có nghiệm là
	\choice
	{$x<0$}
	{\True $0<x \leq 64$}
	{$x \geq 64$}
	{$x\leq 64$}
	\loigiai{
		Điều kiện $x>0$.\\
		Đặt $t=\log_{64}x \Rightarrow x=64^t$.\\
		Khi đó bất phương trình đã cho trở thành\\
		$\log_6\left(4^t+2^t\right)\geq t \Leftrightarrow 4^t+2^t \geq 6^t \Leftrightarrow \left(\dfrac{2}{3}\right)^x+\left(\dfrac{1}{3}\right)^x \geq 1$.\\
		Xét hàm số $f(t)=\left(\dfrac{2}{3}\right)^x+\left(\dfrac{1}{3}\right)^x$ là hàm số nghịch biến và $f(1)=1$.\\
		Do đó bất phương trình có nghiệm $t<1 \Leftrightarrow \log_{64}x \leq 1 \Leftrightarrow 0<x\leq 64$.
	}
\end{ex}
\begin{ex}%[1K6GK-C]
	Số nghiệm nguyên trên $[-2024;2024]$  của bất phương trình $2 x^2-10 x+10>\log _2 \dfrac{2 x-1}{(x-2)^2}$ là
	\choice
	{\True $2022$}
	{$2024$}
	{$2023$}
	{$2025$}
	\loigiai{	
		Điều kiện $x>\dfrac{1}{2}, x \neq 2$.
		Ta có  
		\allowdisplaybreaks
		\begin{eqnarray*}
			&&2 x^2-10 x+10>\log _2 \dfrac{2 x-1}{(x-2)^2}\\
			&\Leftrightarrow& 2(x-2)^2+\log _2(x-2)^2>2 . \dfrac{2 x-1}{2}+\log _2 \dfrac{2 x-1}{2}\\
			&\Leftrightarrow& f\left( (x-2)^2\right) \geq f\left(\dfrac{2 x-1}{2} \right)  \Leftrightarrow (x-2)^2\geq \dfrac{2 x-1}{2}\\
			&\Leftrightarrow& x>\frac{5+\sqrt{7}}{2} ; \quad \frac{5-\sqrt{7}}{2}>x>\frac{1}{2},
		\end{eqnarray*}
		trong đó $f(x)=2 x+\log _2 x$ là hàm 
		đồng biến trên  $(0 ;+\infty)$.\\
		Suy ra bất  phương trình có   nghiệm nguyên là $x\in \{1;\, 4;\, 5;\,\cdots 2024\}$. Do đó số  nghiệm nguyên trên $[-2024;2024]$  của bất phương trình  là $2022$.}
	
\end{ex}
\begin{ex}%[1K6GK-C]
	Giải bất phương trình $\log _2 \dfrac{x+1}{4\left( \sqrt{x}+2\right)} <2(\sqrt{x}-x)$ ta được tập nghiệm $S=\left[a ; \dfrac{b+\sqrt{c}}{2}\right)$, với $\mathrm{a}, \mathrm{b}, \mathrm{c}$ là các số nguyên dương. Tính giá trị biểu thức $T=a+b+c$
	\choice
	{$T=6$}
	{$T=5$}
	{\True $T=8$}
	{$T=7$}
	
	\loigiai{	
		Điều kiện $x \geq 0$.
		Ta có  
		\allowdisplaybreaks
		\begin{eqnarray*}
			&&\log _2 \dfrac{x+1}{4\left( \sqrt{x}+2\right)} <2(\sqrt{x}-x)\\& \Leftrightarrow& \log _2(x+1)<2(\sqrt{x}-x)+2+\log _2(\sqrt{x}+2) \\
			& \Leftrightarrow& \log _2(x+1)+2 x<\log _2[(\sqrt{x}+1)+1]+2(\sqrt{x}+1)\\& \Leftrightarrow& f(x) 	 <f(\sqrt{x}+1)\Leftrightarrow x<\sqrt{x}+1\\ &\Leftrightarrow& x \hoac{&0\leq x<1\\
				&\heva{&x\geq 1\\&x^2-3x+1<0}}\Leftrightarrow
			0\leq  x<\frac{3+\sqrt{5}}{2},
		\end{eqnarray*}
		trong đó $f(x)=\log _2 (x+1)+2x$ là hàm 
		đồng biến trên  $(-1 ;+\infty)$.\\
		Suy ra $a=0$; $b=3$; $c=5$ $\Rightarrow T=8$.}
	
\end{ex}
\begin{ex}%[1K6GK-C]
	Gọi S là tập hợp số nguyên $\mathrm{x}$ thỏa mãn $\log _2 \dfrac{x^2+x+2}{2 x^2-3 x+5} \geq x^2-4 x+3$. Tổng các phần tử của tập hợp $S$ là
	\choice
	{\True $6$}
	{$9$}
	{$8$}
	{$7$}
	\loigiai{	
		Ta có  
		\allowdisplaybreaks
		\begin{eqnarray*}
			&&\log _2 \dfrac{x^2+x+2}{2 x^2-3 x+5} \geq x^2-4 x+3\\& \Leftrightarrow&\log _2\left(x^2+x+2\right)-\log _2\left(2 x^2-3 x+5\right) \geq\left(2 x^2-3 x+5\right)  -\left(x^2+x+2\right) \\
			& \Leftrightarrow& \log _2\left(x^2+x+2\right)+\left(x^2+x+2\right) \geq \log _2\left(2 x^2-3 x+5\right) +\left(2 x^2-3 x+5\right)\\& \Leftrightarrow& f(x^2+x+2) 	 \geq f(2 x^2-3 x+5)\Leftrightarrow x^2+x+2\geq  2x^2-3 x+5\\ &\Leftrightarrow&  x^2-4 x+3\leq 0\Leftrightarrow
			1\leq x\leq 3,
		\end{eqnarray*}
		trong đó $f(x)=\log _2 (x)+x$ là hàm 
		đồng biến trên  $(0 ;+\infty)$.\\
		Suy ra $S=\{1;\, 2;\, 3\}$. Vậy tổng các phần tử của tập hợp $S$ là $6$.}
	
\end{ex}

\Closesolutionfile{ans}
% \begin{indapan}{10}
% 	{ans/ans-1K6-21-Dang14}
% \end{indapan}


